\documentclass[11pt]{amsart}

\usepackage[T1]{fontenc}
\usepackage[utf8]{inputenc}
\usepackage{lmodern}
\usepackage[a4paper,margin=1.05in]{geometry}
\usepackage{microtype}
\usepackage{amsmath,amssymb,amsthm,mathtools,mathrsfs}
\usepackage{enumitem}
\usepackage{booktabs}
\usepackage[dvipsnames]{xcolor}
\usepackage[colorlinks=true,linkcolor=MidnightBlue,citecolor=ForestGreen,urlcolor=BrickRed,
  pdftitle={Split Support Geometry, Version 11},
  pdfauthor={The Clankers},
  pdfsubject={Universal support-amplitude semirings, pre-Boolean support, F-not-two coefficients, spectral ordinal sums, projective support geometry, adelic localization, Frobenius perfection, gerbe quantization, monomial Euler factors, and geometric Langlands}]{hyperref}
\usepackage{array}
\usepackage{longtable}
\usepackage{tikz-cd}

\numberwithin{equation}{section}

\newtheorem{theorem}{Theorem}[section]
\newtheorem{proposition}[theorem]{Proposition}
\newtheorem{corollary}[theorem]{Corollary}
\newtheorem{lemma}[theorem]{Lemma}
\newtheorem{conjecture}[theorem]{Conjecture}
\theoremstyle{definition}
\newtheorem{definition}[theorem]{Definition}
\newtheorem{example}[theorem]{Example}
\newtheorem{remark}[theorem]{Remark}
\newtheorem{warning}[theorem]{Warning}

\newcommand{\B}{\mathbf B}
\newcommand{\Z}{\mathbf Z}
\newcommand{\Q}{\mathbf Q}
\newcommand{\R}{\mathbf R}
\newcommand{\C}{\mathbf C}
\newcommand{\N}{\mathbf N}
\newcommand{\F}{\mathbf F}
\newcommand{\Fone}{\mathbf F_1}
\newcommand{\Fnot}{\mathbf F_{\neg 2}}
\newcommand{\cA}{\mathcal A}
\newcommand{\cH}{\mathcal H}
\newcommand{\cR}{\mathcal R}
\newcommand{\cL}{\mathcal L}
\newcommand{\Frob}{\operatorname{Frob}}
\newcommand{\dist}{\operatorname{dist}}
\newcommand{\Clop}{\operatorname{Clop}}
\newcommand{\No}{\mathbf{No}}
\newcommand{\Spec}{\operatorname{Spec}}
\newcommand{\Proj}{\operatorname{Proj}}
\newcommand{\supp}{\operatorname{supp}}
\newcommand{\Supp}{\operatorname{Supp}}
\newcommand{\im}{\operatorname{Im}}
\newcommand{\kerp}{\operatorname{Ker}}
\newcommand{\Hom}{\operatorname{Hom}}
\newcommand{\End}{\operatorname{End}}
\newcommand{\Aut}{\operatorname{Aut}}
\newcommand{\GL}{\operatorname{GL}}
\newcommand{\PGL}{\operatorname{PGL}}
\newcommand{\Frac}{\operatorname{Frac}}
\newcommand{\ord}{\operatorname{ord}}
\newcommand{\st}{\operatorname{st}}
\newcommand{\gr}{\operatorname{gr}}
\newcommand{\cO}{\mathcal O}
\newcommand{\cM}{\mathcal M}
\newcommand{\cU}{\mathcal U}
\newcommand{\zero}{\mathbf 0}
\newcommand{\Pum}{\mathbf P_{\mathrm{um}}}
\newcommand{\Sp}{\operatorname{Sp}}
\newcommand{\id}{\operatorname{id}}
\newcommand{\colim}{\operatorname*{colim}}
\newcommand{\liminv}{\operatorname*{lim}}
\newcommand{\eps}{\varepsilon}
\newcommand{\taup}{\tau}
\newcommand{\eR}{e}
\newcommand{\ideal}[1]{\mathfrak #1}
\newcommand{\Picf}{\operatorname{Pic}_{\mathrm f}}
\newcommand{\Picar}{\operatorname{Pic}_{\mathrm{ar}}}
\newcommand{\Jacar}{\operatorname{Jac}_{\mathrm{ar}}}
\newcommand{\Idem}{\operatorname{Idem}}
\newcommand{\Ann}{\operatorname{Ann}}
\newcommand{\Spl}{\operatorname{Spl}}
\newcommand{\Nbar}{\overline{\N}}
\newcommand{\Ext}{\operatorname{Ext}}
\newcommand{\Tor}{\operatorname{Tor}}
\newcommand{\Gal}{\operatorname{Gal}}
\newcommand{\Heis}{\operatorname{Heis}}
\newcommand{\Trg}{\operatorname{Trg}}
\newcommand{\CD}{\operatorname{CD}}
\newcommand{\Mat}{\operatorname{Mat}}
\newcommand{\Vect}{\operatorname{Vect}}
\newcommand{\LagCorr}{\operatorname{LagCorr}}
\newcommand{\UHF}{\mathcal Q_{\Omega}}
\newcommand{\Ad}{\operatorname{Ad}}
\newcommand{\EM}{\operatorname{EM}}
\newcommand{\Tr}{\operatorname{Tr}}
\newcommand{\tr}{\operatorname{tr}}
\newcommand{\detsp}{\operatorname{det}_{\mathrm{sp}}}
\newcommand{\charsp}{\operatorname{char}_{\mathrm{sp}}}
\newcommand{\LocSys}{\operatorname{LocSys}}
\newcommand{\Bun}{\operatorname{Bun}}
\newcommand{\Higgs}{\operatorname{Higgs}}
\newcommand{\QCoh}{\operatorname{QCoh}}
\newcommand{\IndCoh}{\operatorname{IndCoh}}
\newcommand{\Nilp}{\operatorname{Nilp}}
\newcommand{\Gm}{\mathbf G_m}
\newcommand{\Nm}{\operatorname N}
\newcommand{\Char}{\operatorname{Char}}
\newcommand{\SpecLat}{\operatorname{Spec}_{\mathrm{lat}}}
\newcommand{\perL}{\operatorname{per}_{L}}
\newcommand{\Hol}{\operatorname{Hol}}


\DeclareMathOperator{\Bool}{Bool}
\DeclareMathOperator{\Mon}{Mon}
\DeclareMathOperator{\Ring}{Ring}
\DeclareMathOperator{\SemiRing}{SemiRing}

\title[Split Support Geometry]
{Split Support Geometry\\[4pt]
\large Universal Zero Fibres, Arithmetic Curves, and Frobenius-Perfect Quantization}
\author{The Clankers}
\date{June 2026\\Pre-publication draft, Version 11}

\begin{document}

\begin{abstract}
A classical zero can arise from absence or from a present datum of zero
amplitude.  For a bounded distributive lattice $L$ and a nonzero
commutative ring $R$, we retain this distinction in the support-amplitude
semiring
\[
G_L(R)=\{(0,\lambda):\lambda\in L\}\cup\{(r,1_L):r\in R\}.
\]
Its amplitude projection is the universal map to a ring, its ring-null
fibre is $L$, and localization at $z_\lambda=(0,\lambda)$ is
$\downarrow\lambda$.  Ideals and prime ideals are classified explicitly,
and the Zariski spectrum is the topological ordinal sum
\[
\Spec G_L(R)\cong\SpecLat(L)\blacktriangleright\Spec R.
\]
For finite-dimensional data,
$\dim G_L(R)=\dim\SpecLat(L)+\dim R+1$.

This support fibre is pre-Boolean.  We write
\[
\Fnot(L):=(L,\vee,\wedge)
\]
for the support coefficient object before two-valued branch selection.
Booleanization is not an intrinsic collapse of the theory but the choice of
a prime branch $L\to\mathbf B$.  For every branch $\epsilon$ there is a
base-change map
\[
G_L(R)\longrightarrow G_{\mathbf B}(R)
\]
which preserves amplitudes and sends the lattice support through
$\epsilon$; on spectra it selects one support generization and then the
full arithmetic closed stratum.  Thus the construction separates three
levels: multiplicative $\mathbf F_1$ data, unresolved pre-Boolean support
$\Fnot(L)$, and Boolean two-valued incidence.

Over a field $K$, unimodular projective space decomposes by support
profiles:
\[
\Pum^n(G_L(K))\cong
\bigsqcup_{\varnothing\neq A\subseteq\{0,\dots,n\}}
\mathbf P^{|A|-1}(K)\times(L\setminus\{1_L\})^{n+1-|A|}.
\]
If $|L|=\ell$, then
\[
\#\Pum^n(G_L(\mathbf F_q))
=\frac{(q+\ell-1)^{n+1}-\ell^{n+1}}{q-1},
\]
with a closed point-counting zeta product.  The split determinant has
ordinary amplitude and lattice-permanent support, while every invertible
matrix remains monomial:
\[
\GL_n(G_L(K))\cong(K^\times)^n\rtimes S_n.
\]
Permutation-cycle lengths and cycle weights classify monomial matrices up
to normalizer conjugacy and compute their characteristic polynomials,
traces, and Euler factors.

For a proper Chen--Moriwaki adelic curve
$S=(K,(\Omega,\mathcal A,\nu),\phi)$, the measure algebra
$\mathscr B_S$ is a complete Boolean support lattice and the product
formula defines the defect transform
\[
\mathscr D_S(f)(a)=\int_\Omega f(\omega)\log|a|_\omega\,d\nu.
\]
A measurable support $A$ is balanced exactly when the restricted adelic
curve $S|_A$ is proper, and support localization gives
\[
G_{\mathscr B_S}(K)[z_A^{-1}]
\cong\downarrow[A]\cong\mathscr B_{S|_A}.
\]
Weighted heights and Arakelov degrees descend projectively precisely on
$\ker\mathscr D_S$.  The Connes--Consani arithmetic Jacobian is
$\operatorname{Pic}_{\mathrm f}\times\mathbf B$, and rigidified
denominator points form a Frobenius tower with globally perfect endpoint
$H_\Omega=\mathbf Q$, stalk $\mathbf F_1[T^{\mathbf Q_+}]$, and carry
extension $0\to\mathbf Z\to\mathbf Q\to\mathbf Q/\mathbf Z\to0$.

Finite denominator phases admit nondegenerate quadratic refinements,
FHLT gerbe quantization, finite Heisenberg algebras, and theta modules.  The
matrix tower has universal UHF limit $\mathcal Q_\Omega$ with
\[
K_0(\mathcal Q_\Omega)\cong(\mathbf Q,\mathbf Q_{\ge0},1),\qquad
\mathcal D_\Omega\cong C(\widehat{\mathbf Z})\cong C^*(\mathbf Q/\mathbf Z).
\]
The identity $\GL_n(G_L(K))\cong N_{\GL_n}(T)(K)$ supplies a monomial
spectral chart: local systems are rank-one local systems on finite covers,
and every loop determinant factors over the cycles of the covering
monodromy.  On the Connes--Consani scaling-site covers these are Frobenius
cycles.  Split characteristic coefficients define a support-refined
Hitchin map, and pullback along the normalizer chart gives a monomial
restriction of the coarse geometric Langlands functor.  Finally, the
reciprocal-root skeleton of the perfect stalk has Beurling realization and
satisfies $x_{\Lambda_Q}\le d_Q^2$; density at the target section is the
Nyman--B\'aez-Duarte criterion and is equivalent to the Riemann hypothesis.
\end{abstract}

\maketitle
\tableofcontents

\section{Introduction: the support-amplitude thesis}
\label{sec:introduction}

The paper studies a single structural operation and its consequences.  A
forgetful map often sends two geometrically different states to the same
classical zero:
\[
\text{absence}\longmapsto0,
\qquad
\text{present zero amplitude}\longmapsto0.
\]
The split-support construction retains the fibre of this map.  Its scalar
form is $G(R)=R\sqcup\{\tau\}$, with actual semiring zero $\tau$ and
supported ring zero $e=0_R$.  Its natural generalization replaces the
binary support bit by a bounded distributive lattice $L$ and produces
$G_L(R)$.  The amplitude layer remains algebraic; the support layer becomes
idempotent and functorial.

The central thesis is that this retained support fibre is not auxiliary
notation.  It has its own ideal theory, spectrum, localizations,
projective geometry, invariant theory, adelic balance law, and quantized
Frobenius tower.  Once these are computed, several constructions that
appear separately in arithmetic geometry, noncommutative geometry,
gerbe quantization, and geometric Langlands become instances of the same
support-amplitude mechanism.

A second thesis is that the support fibre should not be Booleanized too
early.  The Boolean semiring $\mathbf B$ is the first two-valued branch.
The object which precedes this branch is the idempotent support semiring
$\Fnot(L)=(L,\vee,\wedge)$.  The notation $\Fnot$ is only a mnemonic for
``not yet two-valued'': it records unresolved support, not a field with a
specified number of elements.  The formal statements use bounded
lattices, frames, spectra, and branch maps.

\subsection*{Four structural theses}

\begin{enumerate}[label=\textbf{T\arabic*.},leftmargin=3.8em]
\item \textbf{The operative support base is pre-Boolean.}
The coefficient object $\Fnot(L)$ is the unresolved support semiring of a
bounded distributive lattice or frame.  Boolean coefficients arise from a
branch $L\to\mathbf B$, equivalently from a prime support point.  Thus
$\mathbf B$ is not the starting support object; it is the first branch of
one.

\item \textbf{The zero fibre is a geometric coefficient object.}
For every bounded distributive lattice $L$, the semiring $G_L(R)$ has
universal ring reflection $p_L$, support projection $\chi_L$, and
ring-null fibre $L$.  Its prime spectrum is the spectral ordinal sum
\[
\SpecLat(L)\blacktriangleright\Spec R.
\]
Thus the support fibre contributes a genuine open spectral stratum and an
exact Krull-dimension jump.

\item \textbf{Support survives geometric and invariant quotients.}
Projective points decompose by support profiles, finite-field counts and
zeta functions are explicit, and matrix invariants carry two simultaneous
values: an ordinary amplitude invariant and a lattice-valued incidence
invariant.  The full linear symmetry remains the monomial normalizer of a
torus.  Weighted cycle data classify individual monomial operators and
compute their local Euler factors.

\item \textbf{Arithmetic properness is a support-localization condition.}
For an adelic curve, the measure algebra is a complete Boolean support
lattice and the product formula is the vanishing of a defect transform.
Localizing at a measurable support produces the support algebra of the
restricted place family; this restriction remains proper exactly on the
balanced locus.  The arithmetic Jacobian, the split Abel--Jacobi fibre, and
the rigidified denominator tower are concrete realizations of this
principle.

\item \textbf{Frobenius perfection admits finite and operator-algebraic
quantization.}
The denominator tower carries finite quadratic phases, gerbe actions,
Heisenberg extensions, theta modules, and Lagrangian transition
correspondences.  Its matrix realization converges to the universal UHF
algebra, while its diagonal is the cyclotomic algebra
$C^*(\Q/\Z)$.  The same normalizer structure defines a spectral chart for
geometric Langlands and factorizes monodromy determinants over covering
cycles.
\end{enumerate}

These theses give the work its internal direction.  The early algebraic
sections are used later as exact input: the spectrum theorem controls the
new generic strata, lattice localization becomes restriction of an adelic
place family, the determinant theorem becomes a support-refined Hitchin
map, and the monomial normal form becomes a periodic-orbit Euler-factor
formula.  The synthesis therefore produces additional theorems rather than
only parallel descriptions.

\subsection*{Principal theorem clusters}

\begin{center}
\renewcommand{\arraystretch}{1.25}
\begin{tabular}{>{\raggedright\arraybackslash}p{0.18\textwidth}>{\raggedright\arraybackslash}p{0.72\textwidth}}
\toprule
\textbf{Layer} & \textbf{Main output}\\
\midrule
Support algebra & Universal ring reflection, Boolean characters, complete
ideal classification, spectral ordinal sum, support localizations, and
nonassociative coefficient extensions.\\
Projective geometry & Lattice-profile decomposition of projective space,
finite-field point counts and zeta products, monomial linear group, split
determinants, characteristic coefficients, and cycle invariants.\\
Infinitesimal geometry & Boolean-cubical monads, hyperreal and surcomplex
residue fibres, projective zero and infinity, and the separation of support
activation from infinitesimal displacement.\\
Arithmetic geometry & Split Abel--Jacobi fibre, arithmetic Jacobian,
adelic defect transform, balanced localization, rigidified denominator
points, Frobenius perfection, and finite-adelic duality.\\
Quantization and spectral descent & Finite gerbes, Heisenberg and theta
modules, UHF limit, monomial local systems on covers, periodic-orbit Euler
factors, split Hitchin map, and geometric Langlands restriction.\\
Analytic endpoint & Beurling realization of reciprocal roots and the exact
Nyman--B\'aez-Duarte density criterion at the Frobenius-perfect stalk.\\
\bottomrule
\end{tabular}
\end{center}

\subsection*{Why the construction is useful}

The main gain is the availability of exact invariants at points where
ordinary algebra records only zero.  In matrix theory the support invariant
is a lattice permanent or cycle-cover condition.  In adelic geometry it is
a measurable place support together with its product-formula defect.  In
finite quantization it is the distinction between an absent sector and a
defined zero vector.  In spectral geometry it records which sheets and
cycles occur before their amplitudes are combined.

This makes the framework compatible with established theories without
reducing them to a common vocabulary.  The spherical $\Fone$-algebra, the
characteristic-one scaling sheaf, the Boolean support semiring, the
Arakelov metric, the gerbe anomaly line, and the complex amplitude retain
their own operations.  The paper constructs the morphisms among these
layers and computes the fibres that those morphisms forget.

\subsection*{Organization}

The text is divided into five parts.  Part~I develops support-amplitude
algebra and its spectrum.  Part~II treats ideals, projective geometry, and
explicit invariant coordinates.  Part~III develops infinitesimal and
projective analysis.  Part~IV places the construction on adelic and
absolute arithmetic curves and computes Frobenius perfection.  Part~V
constructs the gerbe, UHF, monomial spectral, and analytic realizations.
The final part records the logical source interfaces and condenses the
results into a theorem-level synthesis.

\part{Support-amplitude algebra}

\section{Split-zero algebra and universal maps}

Throughout, a semiring is commutative, unital, and has a specified zero; homomorphisms preserve both $0$ and $1$.  The one-element semiring, in which $0=1$, is denoted by $\zero$ and is allowed.

\begin{definition}[Split-zero globalization]
Let $R$ be a nonzero commutative ring.  Define
\[
S:=G(R):=R\sqcup\{\tau\},
\qquad \tau\notin R,
\]
with operations
\[
r\oplus s=r+s,
\qquad
r\oplus\tau=\tau\oplus r=r,
\qquad
\tau\oplus\tau=\tau,
\]
and
\[
rs=rs,
\qquad
r\tau=\tau r=\tau,
\qquad
\tau^2=\tau.
\]
We write
\[
e:=0_R\in R\subset S.
\]
Thus $0_S=\tau$ and $e\neq\tau$.
\end{definition}

\begin{proposition}[Pair normal form]\label{prop:pair-normal-form}
There is a canonical semiring isomorphism
\[
G(R)\xrightarrow{\sim}
D_R:=\{(0,0)\}\cup(R\times\{1\})
\subseteq R\times\B
\]
given by
\[
\tau\longmapsto(0,0),
\qquad
r\longmapsto(r,1).
\]
Under this isomorphism,
\[
e\longmapsto(0,1).
\]
\end{proposition}

\begin{proof}
The set $D_R$ is closed under the componentwise operations of $R\times\B$.  The displayed map is bijective, and the transported operations are exactly the defining operations of $G(R)$.
\end{proof}

\begin{definition}[Reflection and support]
The ring reflection and support character are
\[
p_R:G(R)\longrightarrow R,
\qquad
p_R(\tau)=0,
\quad
p_R(r)=r,
\]
and
\[
\chi_R:G(R)\longrightarrow\B,
\qquad
\chi_R(\tau)=0,
\quad
\chi_R(r)=1.
\]
\end{definition}

The following two universal properties were proved in \cite{Secondary}; they are recorded because they control every later construction.

\begin{theorem}[Universal ring reflection and Boolean character]\label{thm:universal-maps}
Let $R$ be a nonzero commutative ring.
\begin{enumerate}[label=\textup{(\alph*)}]
\item For every commutative ring $A$, composition with $p_R$ induces a bijection
\[
\Hom_{\Ring}(R,A)
\xrightarrow{\sim}
\Hom_{\SemiRing}(G(R),A).
\]
\item The support map $\chi_R$ is the unique semiring homomorphism $G(R)\to\B$.
\end{enumerate}
\end{theorem}

\begin{proof}
Part (a) is Theorem 2.2 of \cite{Secondary}.  For completeness, if $f:G(R)\to A$ has ring target, then
\[
f(e)=f(e\oplus e)=f(e)+f(e),
\]
so subtraction gives $f(e)=0_A$.  Hence $f|_R$ is a ring homomorphism and $f=(f|_R)\circ p_R$.

For part (b), let $u:G(R)\to\B$.  Since $1\oplus(-1)=e$ in $G(R)$,
\[
u(e)=u(1)+u(-1)=1+u(-1)=1.
\]
If $u(r)=0$ for some supported $r\in R$, then
\[
1=u(e)=u(re)=u(r)u(e)=0,
\]
a contradiction.  Hence every supported element maps to $1$, while $\tau$ maps to $0$.
\end{proof}


\section{Synchronized products, independent mixed support, and fixed loci}
\label{sec:synchronized-products}

The product of two split globalizations and the globalization of a product
ring are different objects.  The difference is exactly the availability of
mixed support.

\begin{theorem}[Boolean fibre-product theorem]
\label{thm:boolean-fibre-product}
For nonzero commutative rings $A$ and $B$, the map
\[
\Phi:G(A\times B)\longrightarrow G(A)\times_{\B}G(B)
\]
defined by
\[
\Phi(\tau)=(\tau,\tau),
\qquad
\Phi(a,b)=(a,b)
\]
is a canonical semiring isomorphism.  The fibre product is taken with
respect to the support characters $\chi_A$ and $\chi_B$.
\end{theorem}

\begin{proof}
A pair $(x,y)\in G(A)\times G(B)$ lies in the fibre product precisely when
$\chi_A(x)=\chi_B(y)$.  If this common value is $0$, then
$(x,y)=(\tau,\tau)$.  If it is $1$, then $x=a\in A$ and $y=b\in B$.
Consequently the fibre product is the disjoint union
\[
\{(\tau,\tau)\}\sqcup(A\times B),
\]
and the displayed map is bijective.  Addition and multiplication are
componentwise on both sides and agree with the defining operations of
$G(A\times B)$, so $\Phi$ is a semiring isomorphism.
\end{proof}

\begin{corollary}[Exact source of mixed support]
\label{cor:mixed-support-source}
Inside the independent product $G(A)\times G(B)$ one has the four support
faces
\[
\{(\tau,\tau)\},\qquad A\times\{\tau\},\qquad
\{\tau\}\times B,\qquad A\times B.
\]
The image of $G(A\times B)$ consists exactly of the first and fourth
faces.  Hence the two middle faces are the genuinely mixed-support states.
Equivalently, the support square
\[
G(A)\times G(B)\longrightarrow\B^2
\]
restricts on $G(A\times B)$ to the diagonal copy of $\B$.
\end{corollary}

\begin{proof}
This is the support-mask decomposition of the Cartesian product together
with Theorem~\ref{thm:boolean-fibre-product}.
\end{proof}

The support maps fit into the commutative pullback diagram
\[
\begin{tikzcd}[column sep=large,row sep=large]
G(A\times B) \arrow[r,hook] \arrow[d,"\chi_{A\times B}"']
& G(A)\times G(B) \arrow[d,"\chi_A\times\chi_B"]\\
\B \arrow[r,"\Delta"'] & \B^2,
\end{tikzcd}
\]
where $\Delta(\epsilon)=(\epsilon,\epsilon)$.  The two off-diagonal
Boolean vertices are exactly the mixed-support faces omitted by the upper
left object.

\begin{theorem}[Fixed locus of a supported unit]
\label{thm:fixed-locus-unit}
Let $R$ be a commutative ring and let $u\in R^\times$.  Then
\[
\operatorname{Fix}_{G(R)}(u)
:=\{x\in G(R):ux=x\}
=\{\tau\}\sqcup\Ann_R(u-1).
\]
In particular, if $K$ is a field, then
\[
\operatorname{Fix}_{G(K)}(u)=
\begin{cases}
G(K),&u=1,\\
\{\tau,e\}\cong\B,&u\ne1.
\end{cases}
\]
\end{theorem}

\begin{proof}
The element $\tau$ is fixed because it is multiplicatively absorbing.  A
supported element $r\in R$ is fixed precisely when $ur=r$, equivalently
$(u-1)r=0$.  This proves the first formula.  Over a field, the annihilator
of $u-1$ is all of $K$ when $u=1$ and is $\{0\}$ otherwise.
\end{proof}

\begin{corollary}[The Boolean involution locus]
\label{cor:boolean-involution-locus}
If $R$ is an integral domain of characteristic different from $2$, then
\[
\operatorname{Fix}_{G(R)}(-1)=\{\tau,e\}\cong\B.
\]
\end{corollary}

\begin{proof}
Apply Theorem~\ref{thm:fixed-locus-unit} and use
$\Ann_R(-2)=\{0\}$.
\end{proof}

\section{Idempotent localization and the division trichotomy}

We use localization only through its universal property.  If $A$ is a semiring and $M\subseteq A$ is multiplicatively closed, a localization of $A$ at $M$ is a semiring $M^{-1}A$ with a map $\lambda:A\to M^{-1}A$ such that every $\lambda(m)$ is invertible and every map from $A$ that inverts $M$ factors uniquely through $\lambda$.

\begin{theorem}[Localization at an idempotent]\label{thm:idempotent-localization}
Let $A$ be a commutative semiring and let $a\in A$ satisfy $a^2=a$.  Then
\[
A[a^{-1}]\cong aA,
\]
where $aA=\{ax:x\in A\}$ is regarded as a semiring with zero $a0_A$ and multiplicative identity $a$.
\end{theorem}

\begin{proof}
The subset $aA$ is closed under addition and multiplication because
\[
ax+ay=a(x+y),
\qquad
(ax)(ay)=a^2xy=axy.
\]
The element $a$ is its multiplicative identity.  Define
\[
\pi_a:A\longrightarrow aA,
\qquad
x\longmapsto ax.
\]
Then
\[
\pi_a(a)=a^2=a=1_{aA},
\]
so $a$ becomes invertible.

Let $f:A\to T$ be a semiring homomorphism such that $f(a)$ is invertible.  Since $f(a)$ is idempotent,
\[
f(a)^2=f(a).
\]
Multiplication by $f(a)^{-1}$ gives $f(a)=1_T$.  Define
\[
\overline f:aA\longrightarrow T,
\qquad
\overline f(ax)=f(x).
\]
If $ax=ay$, then
\[
f(a)f(x)=f(a)f(y),
\]
and $f(a)=1_T$, so $f(x)=f(y)$.  Thus $\overline f$ is well defined.  It is a semiring homomorphism and satisfies $f=\overline f\circ\pi_a$.  Uniqueness follows because every element of $aA$ has the form $\pi_a(x)$.
\end{proof}

\begin{corollary}[The two zero localizations]\label{cor:two-zero-localizations}
For every nonzero commutative ring $R$,
\[
G(R)[e^{-1}]\cong\{\tau,e\}\cong\B,
\]
and
\[
G(R)[\tau^{-1}]\cong\{\tau\}\cong\zero.
\]
Under the first isomorphism, the localization map is the support character:
\[
\frac{x}{e}=\chi_R(x).
\]
\end{corollary}

\begin{proof}
Both $e$ and $\tau$ are idempotent.  Moreover,
\[
eG(R)=\{\tau,e\},
\qquad
\tau G(R)=\{\tau\}.
\]
In $eG(R)$ the element $e$ is the identity and $\tau$ is the zero, so $eG(R)\cong\B$.  The localization map is $x\mapsto ex$, which sends $\tau$ to $\tau$ and every supported element to $e$; this is exactly $\chi_R$ after the identification $\tau\leftrightarrow0$, $e\leftrightarrow1$.
\end{proof}

\begin{theorem}[Compatibility with ordinary ring localization]\label{thm:ring-localization}
Let $M\subseteq R$ be a multiplicatively closed set with $0\notin M$.  Then the canonical map $R\to M^{-1}R$ induces a canonical semiring isomorphism
\[
M^{-1}G(R)\xrightarrow{\sim}G(M^{-1}R).
\]
\end{theorem}

\begin{proof}
Let
\[
\lambda:G(R)\longrightarrow G(M^{-1}R)
\]
be defined by $\lambda(\tau)=\tau$ and $\lambda(r)=r/1$.  Every $m\in M$ maps to the supported unit $m/1$.

Let $f:G(R)\to T$ be a semiring homomorphism such that every $f(m)$ is invertible.  Define
\[
\overline f(\tau):=0_T,
\qquad
\overline f(r/m):=f(r)f(m)^{-1}.
\]
If $r/m=s/n$ in $M^{-1}R$, then there exists $u\in M$ with
\[
u(nr-ms)=0
\]
in $R$, equivalently $unr=ums$.  Applying $f$ gives
\[
f(u)f(n)f(r)=f(u)f(m)f(s).
\]
Since $f(u)$, $f(m)$, and $f(n)$ are invertible, one obtains
\[
f(r)f(m)^{-1}=f(s)f(n)^{-1}.
\]
Thus $\overline f$ is well defined.  The usual fraction calculations show that it is a semiring homomorphism.  It is the unique map satisfying $\overline f\circ\lambda=f$, because the supported fractions and $\tau$ exhaust $G(M^{-1}R)$.
\end{proof}

\begin{corollary}[Classical fraction field versus Boolean fraction residue]\label{cor:fraction-comparison}
Let $R$ be an integral domain with fraction field $K$.
\begin{enumerate}[label=\textup{(\alph*)}]
\item Localizing at the classically nonzero elements gives
\[
G(R)[(R\setminus\{0\})^{-1}]
\cong
G(K).
\]
\item Localizing at all semiring-nonzero elements, namely at $R=G(R)\setminus\{\tau\}$, gives
\[
G(R)[R^{-1}]
\cong
\B.
\]
\end{enumerate}
\end{corollary}

\begin{proof}
Part (a) is Theorem \ref{thm:ring-localization}.  For part (b), the denominator set contains $e$, so the localization factors through $G(R)[e^{-1}]\cong\B$.  In $\B$, every supported element of $G(R)$ already maps to $1$, hence is invertible.  Therefore $\B$ has the universal property of the localization at all supported elements.
\end{proof}

\begin{remark}[Element localization versus localization at a prime]\label{rem:prime-v-element}
For $R=\Z$, let
\[
P_0=\{\tau,0\}.
\]
The stalk localization at $P_0$ inverts the complement $\Z\setminus\{0\}$ and is therefore $G(\Q)$.  Principal localization at the element $e=0$ instead gives $\B$.  Thus ``localizing at the generic zero prime'' and ``placing the zero element in the denominator set'' are distinct operations.
\end{remark}

\begin{theorem}[Strict identification with a nonzero amplitude]\label{thm:strict-identification-collapse}
Let $K$ be a field, let $S=G(K)$, and let $h\in K^\times$.  Let $\sim$ be the smallest semiring congruence on $S$ satisfying
\[
e\sim h.
\]
Then
\[
S/{\sim}\cong\B.
\]
\end{theorem}

\begin{proof}
In the quotient, the class of $h$ is invertible, so the class of $e$ is invertible.  Since $e^2=e$, its class is an invertible idempotent and therefore equals $1$.  For every supported $x\in K$,
\[
ex=e
\]
in $S$, hence in the quotient
\[
x=1x=ex=e=1.
\]
Thus all supported elements lie in one congruence class.  The support character $\chi_K:S\to\B$ satisfies $\chi_K(e)=\chi_K(h)=1$, so it coequalizes the imposed relation and distinguishes $\tau$ from the supported class.  Hence the quotient has exactly two elements and is Boolean.
\end{proof}

\begin{theorem}[Division trichotomy]\label{thm:division-trichotomy}
Let $R$ be an integral domain with fraction field $K$.  The three universal division operations are
\[
\boxed{
\begin{aligned}
\text{invert }R\setminus\{0\}
&\quad\Longrightarrow\quad G(K),\\
\text{invert }e=0_R
&\quad\Longrightarrow\quad \B,\\
\text{invert }\tau=0_{G(R)}
&\quad\Longrightarrow\quad \zero.
\end{aligned}}
\]
The first preserves amplitude and support, the second preserves only support, and the third preserves no distinction.
\end{theorem}

\begin{proof}
This is the conjunction of Corollaries \ref{cor:two-zero-localizations} and \ref{cor:fraction-comparison}.
\end{proof}

\section{Ring-null infinitesimals and Boolean cubes}

The preceding localization theorem suggests an intrinsic infinitesimal notion which is independent of order, topology, or valuation.

\begin{definition}[Absolute ring-null halo]\label{def:absolute-halo}
Let $A$ be a commutative semiring admitting a universal ring reflection
\[
\rho_A:A\longrightarrow A^{\mathrm{ring}}.
\]
Define
\[
\mathfrak I_{\mathrm{abs}}(A):=\rho_A^{-1}(0).
\]
An element of $\mathfrak I_{\mathrm{abs}}(A)\setminus\{0_A\}$ is called an \emph{absolute ring-null infinitesimal}.
\end{definition}

\begin{proposition}[Intrinsic characterization]\label{prop:absolute-characterization}
For $x\in A$,
\[
x\in\mathfrak I_{\mathrm{abs}}(A)
\]
if and only if every semiring homomorphism from $A$ to a commutative ring sends $x$ to zero.
\end{proposition}

\begin{proof}
Every map from $A$ to a ring factors uniquely through $\rho_A$.  Hence $\rho_A(x)=0$ implies that every ring-valued map kills $x$.  Conversely, apply the assumed condition to the ring-valued map $\rho_A$ itself.
\end{proof}

\begin{corollary}[The scalar split-zero infinitesimal]\label{cor:scalar-absolute-halo}
For $S=G(R)$,
\[
\mathfrak I_{\mathrm{abs}}(S)=\{\tau,e\}.
\]
Thus $e$ is the unique nonzero absolute ring-null infinitesimal of $S$.
\end{corollary}

\begin{proof}
The universal ring reflection is $p_R$ by Theorem \ref{thm:universal-maps}, and
\[
p_R^{-1}(0_R)=\{\tau,e\}.
\]
The semiring zero is $\tau$, so $e$ is the unique nonzero member of this fibre.
\end{proof}

Let $I$ be a finite set and write
\[
S^I=\prod_{i\in I}S.
\]
For $A\subseteq I$, define the supported-zero mask
\[
(\eps_A)_i=
\begin{cases}
e,&i\in A,\\
\tau,&i\notin A.
\end{cases}
\]

\begin{theorem}[Finite-product ring reflection]\label{thm:product-ring-reflection}
The coordinatewise map
\[
p_I:=p_R^I:S^I\longrightarrow R^I
\]
is the universal ring reflection of $S^I$.  Its kernel is
\[
p_I^{-1}(0)=\{\eps_A:A\subseteq I\}.
\]
Moreover,
\[
\eps_A+\eps_B=\eps_{A\cup B},
\qquad
\eps_A\eps_B=\eps_{A\cap B}.
\]
Consequently,
\[
\mathfrak I_{\mathrm{abs}}(S^I)
\cong\B^I
\]
as semirings.
\end{theorem}

\begin{proof}
Let $f:S^I\to T$ be a semiring homomorphism with ring target.  Each $\eps_A$ is additively idempotent:
\[
\eps_A+\eps_A=\eps_A.
\]
Therefore
\[
f(\eps_A)+f(\eps_A)=f(\eps_A),
\]
and subtraction in $T$ gives $f(\eps_A)=0$.

For $r=(r_i)_{i\in I}\in R^I$, let $\widetilde r\in S^I$ be its fully supported lift, so $(\widetilde r)_i=r_i\in R\subset S$.  Define
\[
g:R^I\longrightarrow T,
\qquad
g(r):=f(\widetilde r).
\]
Because sums and products of fully supported lifts remain fully supported,
\[
\widetilde{r+s}=\widetilde r+\widetilde s,
\qquad
\widetilde{rs}=\widetilde r\,\widetilde s,
\]
where the operations on $R^I$ are coordinatewise.  Hence $g$ is a ring homomorphism.

If $x\in S^I$ and $r=p_I(x)$, then $x$ and $\widetilde r$ differ only in coordinates where $r_i=0$: in such a coordinate $x_i$ is either $\tau$ or $e$, while $(\widetilde r)_i=e$.  Therefore there is a mask $\eps_A$ such that
\[
\widetilde r=x+\eps_A.
\]
Applying $f$ gives
\[
f(\widetilde r)=f(x)+f(\eps_A)=f(x).
\]
Thus $f=g\circ p_I$.  Uniqueness follows from surjectivity of $p_I$.

The kernel description is immediate coordinatewise.  The union and intersection formulas follow from
\[
e+e=e,
\quad e+\tau=e,
\quad \tau+\tau=\tau,
\]
and
\[
e^2=e,
\quad e\tau=\tau,
\quad \tau^2=\tau.
\]
These are exactly the join and meet operations on $\B^I$.
\end{proof}

\begin{example}[The mixed zero $(\tau,e)$]
In $S^2$ the four absolute infinitesimals are
\[
(\tau,\tau),
\quad
(e,\tau),
\quad
(\tau,e),
\quad
(e,e).
\]
The element $(\tau,e)$ means that the first coordinate is absent while the second coordinate is supported but arithmetically zero.  For $a\in R$,
\[
(a,\tau)+(-a,\tau)=(e,\tau),
\]
so cancellation produces supported zero rather than absence.
\end{example}

\subsection{Derivations and the failure of tangent displacement}

\begin{definition}[Semiring derivation]
Let $A$ be a semiring and let $M$ be an $A$-semimodule.  A derivation is a map $d:A\to M$ satisfying
\[
d(0)=0,
\qquad
d(1)=0,
\qquad
d(x+y)=d(x)+d(y),
\]
and
\[
d(xy)=x\,d(y)+y\,d(x).
\]
\end{definition}

\begin{proposition}[The supported zero is derivationally constant]\label{prop:derivation-kills-e}
Every derivation $d:G(R)\to M$ satisfies
\[
d(e)=0.
\]
Every derivation $\B\to M$ is zero.
\end{proposition}

\begin{proof}
Since $1+e=1$ in $G(R)$,
\[
0=d(1)=d(1+e)=d(1)+d(e)=d(e).
\]
For $\B$, the only elements are $0$ and $1$, both of which are killed by the defining axioms of a derivation.
\end{proof}

\begin{remark}
The identity $x+e=x$ for every supported $x\in R$ shows that $e$ does not move a supported point.  It is therefore an infinitesimal only in the ring-null support sense of Definition \ref{def:absolute-halo}; it is not a first-order tangent parameter.
\end{remark}

The closest ring-valued multiplicative model of the relation $e^2=e$ makes this distinction exact.

\begin{theorem}[Idempotent probes are finite-difference probes]\label{thm:idempotent-probe}
Let $A$ be a commutative ring and set
\[
C:=A[q]/(q^2-q).
\]
Then evaluation at $q=0$ and $q=1$ induces a canonical ring isomorphism
\[
C\xrightarrow{\sim}A\times A.
\]
For every polynomial $f\in A[T]$ and every $x,h\in A$,
\[
f(x+qh)
=
f(x)+q\bigl(f(x+h)-f(x)\bigr)
\]
in $C$.  Moreover,
\[
\Omega^1_{C/A}=0.
\]
\end{theorem}

\begin{proof}
The ideals $(q)$ and $(q-1)$ of $A[q]$ are comaximal because
\[
q-(q-1)=1.
\]
Their intersection equals their product $(q(q-1))=(q^2-q)$.  The Chinese remainder theorem gives
\[
A[q]/(q^2-q)
\cong
A[q]/(q)\times A[q]/(q-1)
\cong
A\times A.
\]
Under this isomorphism, $q\mapsto(0,1)$ and $x+qh\mapsto(x,x+h)$.  Both sides of the polynomial identity map to
\[
\bigl(f(x),f(x+h)\bigr),
\]
so they are equal.

For relative K\"ahler differentials, the relation $q^2-q=0$ gives
\[
(2q-1)dq=0.
\]
But
\[
(2q-1)^2=4q^2-4q+1=1
\]
in $C$, so $2q-1$ is invertible and $dq=0$.  Since $C$ is generated over $A$ by $q$, this implies $\Omega^1_{C/A}=0$.
\end{proof}

\begin{corollary}[Idempotence versus nilpotence]
An idempotent parameter records two endpoint values and exact finite differences.  A nilpotent parameter $\epsilon^2=0$ records first-order tangent data.  These structures cannot be identified in a nonzero ring.
\end{corollary}

\begin{proof}
The finite-difference statement is Theorem \ref{thm:idempotent-probe}.  If a ring element satisfies both $q^2=q$ and $q^2=0$, then $q=0$.
\end{proof}

\section{Mixed-support localization and scoped identities}

\begin{lemma}[Localization of a finite product]\label{lem:product-localization}
Let $A_i$ be commutative semirings indexed by a finite set $I$, let
\[
A:=\prod_{i\in I}A_i,
\qquad
x=(x_i)_{i\in I}\in A.
\]
Then
\[
A[x^{-1}]
\cong
\prod_{i\in I}A_i[x_i^{-1}].
\]
\end{lemma}

\begin{proof}
Use the standard fraction construction for localization.  Every element of $A[x^{-1}]$ is represented by a fraction
\[
\frac{y}{x^n},
\qquad y\in A,
\quad n\ge0.
\]
Define
\[
\Phi\left(\frac{y}{x^n}\right)
:=
\left(\frac{y_i}{x_i^n}\right)_{i\in I}.
\]
If two fractions are equal in $A[x^{-1}]$, their coordinate fractions are equal, so $\Phi$ is well defined.

For surjectivity, choose fractions $y_i/x_i^{n_i}$ in the factors and put $N=\max_i n_i$.  Replacing $y_i$ by $y_i x_i^{N-n_i}$ gives a common exponent $N$, hence a preimage in $A[x^{-1}]$.

For injectivity, suppose the coordinate fractions of $y/x^n$ and $z/x^m$ are equal.  For each $i$ there exists $k_i\ge0$ such that
\[
x_i^{k_i+m}y_i=x_i^{k_i+n}z_i.
\]
Let $k=\max_i k_i$.  Multiplying each equality by a further power of $x_i$ yields
\[
x_i^{k+m}y_i=x_i^{k+n}z_i
\]
for every $i$.  Therefore
\[
x^{k+m}y=x^{k+n}z
\]
in $A$, which is precisely the equality of the original fractions in $A[x^{-1}]$.
\end{proof}

Let $K$ be a field, $S=G(K)$, and $I$ finite.  For $x\in S^I$, define
\[
U(x):=\{i\in I:x_i\in K^\times\},
\]
\[
Z(x):=\{i\in I:x_i=e\},
\]
and
\[
N(x):=\{i\in I:x_i=\tau\}.
\]
These form a partition of $I$.

\begin{theorem}[Mixed-support localization normal form]\label{thm:mixed-localization}
For every $x\in S^I$,
\[
S^I[x^{-1}]
\cong
S^{U(x)}\times\B^{Z(x)}.
\]
The coordinates in $N(x)$ contribute one-element factors and hence disappear.  Under the displayed isomorphism,
\[
\frac{y}{x}
\longmapsto
\left(
\left(\frac{y_i}{x_i}\right)_{i\in U(x)},
\left(\chi_K(y_i)\right)_{i\in Z(x)}
\right).
\]
\end{theorem}

\begin{proof}
By Lemma \ref{lem:product-localization},
\[
S^I[x^{-1}]
\cong
\prod_{i\in I}S[x_i^{-1}].
\]
If $i\in U(x)$, then $x_i$ is already a unit of $S$, so $S[x_i^{-1}]\cong S$.  If $i\in Z(x)$, Corollary \ref{cor:two-zero-localizations} gives $S[e^{-1}]\cong\B$.  If $i\in N(x)$, the same corollary gives $S[\tau^{-1}]\cong\zero$.  Since a finite product with a one-element factor is canonically isomorphic to the product of the remaining factors, the stated formula follows.  The explicit fraction formula is the coordinatewise description of these three identifications.
\end{proof}

\begin{example}
For $x=(\tau,e)\in S^2$,
\[
S^2[x^{-1}]\cong\B,
\]
and
\[
\frac{(a,b)}{(\tau,e)}=\chi_K(b).
\]
The first coordinate is erased; the second records only support.
\end{example}

\begin{example}
For $x=(h,e)\in S^2$ with $h\in K^\times$,
\[
S^2[x^{-1}]\cong S\times\B,
\]
and
\[
\frac{(a,b)}{(h,e)}=\bigl(ah^{-1},\chi_K(b)\bigr).
\]
This is ordinary division in one coordinate and Boolean division by supported zero in the other.
\end{example}

We next formalize the effect of adjoining an equality only on selected support coordinates.

For $i\in I$ and $s\in S$, write $r_i(s)\in S^I$ for the vector with entry $s$ at $i$ and $\tau$ elsewhere.  Fix a subset $A\subseteq I$ and units $h_i\in K^\times$ for $i\in A$.  Define
\[
\eps_A:=\sum_{i\in A}r_i(e),
\qquad
h_A:=\sum_{i\in A}r_i(h_i).
\]

\begin{theorem}[Scoped identity quotient]\label{thm:scoped-identity}
Let $\sim_A$ be the smallest semiring congruence on $S^I$ satisfying
\[
\eps_A\sim_A h_A.
\]
Then
\[
S^I/{\sim_A}
\cong
\B^A\times S^{I\setminus A}.
\]
The quotient map is
\[
q_A(x)
=
\left(
(\chi_K(x_i))_{i\in A},
(x_j)_{j\notin A}
\right).
\]
\end{theorem}

\begin{proof}
The map $q_A$ is a semiring homomorphism and satisfies
\[
q_A(\eps_A)=q_A(h_A),
\]
so it factors through $S^I/{\sim_A}$.

Let $f:S^I\to T$ be any semiring homomorphism satisfying $f(\eps_A)=f(h_A)$.  For $i\in A$, multiply the relation by $r_i(1)$ to obtain
\[
f(r_i(e))=f(r_i(h_i)).
\]
Multiplying again by $r_i(h_i^{-1})$ gives
\[
f(r_i(e))=f(r_i(1)).
\]
For every supported $s\in K$,
\[
r_i(s)=r_i(s)r_i(1),
\]
and therefore
\[
f(r_i(s))
=f(r_i(s))f(r_i(1))
=f(r_i(s))f(r_i(e))
=f(r_i(se))
=f(r_i(e)).
\]
Thus on coordinate $i\in A$, the value of $f$ depends only on whether the coordinate is $\tau$ or supported.  On coordinates outside $A$, no relation has been imposed.

Every vector $x\in S^I$ decomposes as
\[
x=\sum_{i\in I}r_i(x_i).
\]
Hence $f(x)$ depends only on $q_A(x)$, which gives a unique factorization of $f$ through $q_A$.  This is the universal property of the quotient.
\end{proof}

\begin{corollary}[Modal equality without arithmetic equality]\label{cor:modal-equality}
For every $h\in K^\times$,
\[
\chi_K(e)=\chi_K(h)=1.
\]
Thus $e$ and $h$ are equal after support reflection.  Their strict arithmetic identification gives the Boolean quotient by Theorem \ref{thm:strict-identification-collapse}; their identification on selected coordinates gives the partial Booleanization of Theorem \ref{thm:scoped-identity}.
\end{corollary}


\section{Distributive-lattice support and higher zero fibres}
\label{sec:lattice-support}

The scalar split-zero semiring and the Boolean cube are the first two
instances of a construction controlled by a bounded distributive lattice.
This section gives the construction in a form which simultaneously covers
independent mixed support, iterated sub-zero layers, and nilpotent ideal
filtrations.

Let $L$ be a bounded distributive lattice with bottom $0_L$, top $1_L$,
join $\vee$, and meet $\wedge$.  Regard $L$ as an idempotent semiring by
\[
\lambda\oplus\mu:=\lambda\vee\mu,
\qquad
\lambda\odot\mu:=\lambda\wedge\mu.
\]
Its additive zero is $0_L$ and its multiplicative unit is $1_L$.

\begin{definition}[Lattice-split globalization]
\label{def:lattice-split}
For a commutative ring $R$, define
\[
G_L(R):=
\{(0,\lambda):\lambda\in L\}
\cup
\{(r,1_L):r\in R\}
\subseteq R\times L,
\]
with the componentwise operations
\begin{align*}
(r,\lambda)+(s,\mu)&=(r+s,\lambda\vee\mu),\\
(r,\lambda)(s,\mu)&=(rs,\lambda\wedge\mu).
\end{align*}
Write
\[
\tau_L:=(0,0_L),
\qquad
e_L:=(0,1_L).
\]
\end{definition}

The point $\tau_L$ is actual absence, whereas $e_L$ is the supported ring
zero.  A nonzero amplitude can occur only over maximal support $1_L$.

\begin{theorem}[Lattice normal form and universal maps]
\label{thm:lattice-normal-form}
The following statements hold.
\begin{enumerate}[label=\textup{(\arabic*)}]
\item $G_L(R)$ is a commutative unital semiring with zero $\tau_L$ and
unit $(1_R,1_L)$.
\item Projection to amplitude
\[
p_L:G_L(R)\longrightarrow R,
\qquad
p_L(r,\lambda)=r,
\]
is the universal morphism from $G_L(R)$ to a ring.
\item Projection to support
\[
\chi_L:G_L(R)\longrightarrow L,
\qquad
\chi_L(r,\lambda)=\lambda,
\]
is a semiring homomorphism.
\item The ring-null fibre is canonically the lattice semiring:
\[
p_L^{-1}(0_R)
=\{(0,\lambda):\lambda\in L\}
\cong (L,\vee,\wedge).
\]
\end{enumerate}
\end{theorem}

\begin{proof}
The subset in Definition~\ref{def:lattice-split} is closed under both
componentwise operations.  If one summand has nonzero amplitude, its
support is $1_L$, so the sum again has support $1_L$.  If one factor has
support below $1_L$, its amplitude is zero, so the product has zero
amplitude and support below or equal to that factor.  Associativity and
commutativity follow componentwise.  Distributivity follows from
\[
\lambda\wedge(\mu\vee\nu)
=(\lambda\wedge\mu)\vee(\lambda\wedge\nu)
\]
and the distributive law in $R$.  The stated zero and unit are immediate.

The maps $p_L$ and $\chi_L$ preserve the operations.  Let
$f:G_L(R)\to A$ be a semiring homomorphism into a ring $A$.  For every
$\lambda\in L$, put $z_\lambda=(0,\lambda)$.  Since
\[
e_L+z_\lambda=e_L,
\]
additive cancellation in $A$ gives $f(z_\lambda)=0$.  Hence $f$ is
determined by its restriction to the supported copy of $R$, and that
restriction is a ring homomorphism.  Conversely, every ring homomorphism
$R\to A$ extends uniquely by sending all $z_\lambda$ to zero.  This is the
universal property of $p_L$.  The description of the zero fibre is
immediate, and its inherited operations are join and meet.
\end{proof}


\begin{theorem}[Classification of Boolean characters]
\label{thm:lattice-boolean-characters}
Let $R$ be a nonzero commutative ring and let $L$ be a nontrivial bounded
distributive lattice.  Restriction to the ring-null fibre induces a canonical bijection
\[
\boxed{
\Hom_{\SemiRing}(G_L(R),\B)
\cong
\Hom_{\mathrm{DLat}_{0,1}}(L,\B).
}
\]
Explicitly, if $h:L\to\B$ is a bounded lattice homomorphism, the
corresponding semiring homomorphism $\widetilde h:G_L(R)\to\B$ is
\[
\widetilde h(0,\lambda)=h(\lambda),
\qquad
\widetilde h(r,1_L)=1\quad(r\in R).
\]
\end{theorem}

\begin{proof}
Let $u:G_L(R)\to\B$ be a semiring homomorphism.  Since
$1+(-1)=e_L=(0,1_L)$,
\[
u(e_L)=u(1)\vee u(-1)=1.
\]
For every supported amplitude $r\in R$, one has $r e_L=e_L$, hence
\[
1=u(e_L)=u(r)u(e_L)=u(r).
\]
Thus $u$ is constant equal to $1$ on the supported copy of $R$.  Its
restriction
\[
h(\lambda):=u(0,\lambda)
\]
preserves $0_L,1_L$, joins, and meets because the ring-null fibre has
exactly those operations.

Conversely, let $h:L\to\B$ be a bounded lattice homomorphism and define
$\widetilde h$ by the displayed formula.  The sum and product of two
supported elements are supported, including when their amplitude cancels
to zero.  A supported element added to $(0,\lambda)$ remains supported,
and its product with $(0,\lambda)$ equals $(0,\lambda)$.  These identities,
together with preservation of joins and meets by $h$, verify the semiring
axioms for $\widetilde h$.  The two constructions are inverse.
\end{proof}

\begin{theorem}[Finite Birkhoff support coordinates]
\label{thm:birkhoff-support-coordinates}
Let $L$ be a finite nontrivial bounded distributive lattice and let $J(L)$ be its
poset of nonzero join-irreducible elements.  For $j\in J(L)$ define
\[
\theta_j(\lambda)=
\begin{cases}
1,&j\leq\lambda,\\
0,&j\nleq\lambda.
\end{cases}
\]
Then each $\theta_j:L\to\B$ is a bounded lattice homomorphism, every
bounded lattice homomorphism $L\to\B$ is of this form, and
\[
\eta_L:L\longrightarrow\B^{J(L)},
\qquad
\eta_L(\lambda)=(\theta_j(\lambda))_{j\in J(L)},
\]
is an injective semiring homomorphism with image
\[
\operatorname{Idl}(J(L))
\subseteq\mathcal P(J(L))\cong\B^{J(L)}.
\]
Consequently,
\[
\boxed{
G_L(R)\hookrightarrow R\times\B^{J(L)},
\qquad
x\longmapsto\bigl(p_L(x),(\theta_j\chi_L(x))_j\bigr),
}
\]
is an injective semiring homomorphism with image
\[
\{(0,\eta_L(\lambda)):\lambda\in L\}
\cup
\bigl(R\times\{\mathbf1\}\bigr).
\]
\end{theorem}

\begin{proof}
In a finite distributive lattice, a join-irreducible element is join-prime:
if $j\leq x\vee y$, then
\[
j=(j\wedge x)\vee(j\wedge y),
\]
so join-irreducibility gives $j\leq x$ or $j\leq y$.  Therefore
$\theta_j$ preserves joins; preservation of meets and the bounds is
immediate.

Every nonzero element of a finite lattice is a join of join-irreducibles
below it, by induction on the length of the interval below that element.
Hence $\eta_L$ is injective.  The set
\[
I_\lambda=\{j\in J(L):j\leq\lambda\}
\]
is an order ideal.  Conversely, if $I\subseteq J(L)$ is an order ideal and
$\lambda=\bigvee I$, join-primality shows
$I_\lambda=I$.  This proves the image description.

A bounded lattice homomorphism $h:L\to\B$ has prime filter
$h^{-1}(1)$.  Since $L$ is finite, this filter has a least element $j$;
primality makes $j$ join-prime and hence join-irreducible, and
$h=\theta_j$.  The final embedding is the product of $p_L$ with all
Boolean characters from Theorem~\ref{thm:lattice-boolean-characters}.
Its injectivity follows because $p_L$ separates amplitudes and $\eta_L$
separates the zero-amplitude support strata.
\end{proof}

\begin{corollary}[Threshold reconstruction of a finite support history]
\label{cor:chain-threshold-reconstruction}
Let $C_n=\{0<1<\cdots<n\}$ with join $\max$ and meet $\min$.  Its Boolean
characters are exactly
\[
\vartheta_k(j)=\mathbf1_{\{j\geq k\}},
\qquad 1\leq k\leq n.
\]
The product map
\[
\Theta_n:C_n\longrightarrow\B^n,
\qquad
j\longmapsto
(\vartheta_1(j),\ldots,\vartheta_n(j)),
\]
is injective and has image
\[
\{(\epsilon_1,\ldots,\epsilon_n)\in\B^n:
\epsilon_1\geq\epsilon_2\geq\cdots\geq\epsilon_n\}.
\]
Thus the complete ordered history of an $n$-stage support state is recovered
from its threshold shadows.
\end{corollary}

\begin{proof}
Every element $1,\ldots,n$ is join-irreducible in $C_n$, and
Theorem~\ref{thm:birkhoff-support-coordinates} applies.  The image of $j$
is the string of $j$ initial ones followed by $n-j$ zeros.
\end{proof}

\begin{theorem}[Localization at a support stratum]
\label{thm:lattice-localization}
For $\lambda\in L$, let $z_\lambda=(0,\lambda)$.  Then
$z_\lambda^2=z_\lambda$ and
\[
G_L(R)[z_\lambda^{-1}]
\cong
\downarrow\lambda
:=\{\mu\in L:\mu\leq\lambda\},
\]
where $\downarrow\lambda$ is equipped with join and meet, zero $0_L$, and
unit $\lambda$.  In particular,
\[
G_L(R)[e_L^{-1}]\cong L,
\qquad
G_L(R)[\tau_L^{-1}]\cong\zero.
\]
\end{theorem}

\begin{proof}
The element $z_\lambda$ is multiplicatively idempotent.  By
Theorem~\ref{thm:idempotent-localization}, localization at it is the
principal semiring $z_\lambda G_L(R)$.  Multiplication gives
\[
z_\lambda(r,1_L)=z_\lambda,
\qquad
z_\lambda z_\mu=z_{\lambda\wedge\mu}.
\]
Therefore
\[
z_\lambda G_L(R)=\{z_\nu:\nu\leq\lambda\},
\]
and the induced operations are join and meet on the principal interval.
The endpoint cases are $\lambda=1_L$ and $\lambda=0_L$.
\end{proof}

\begin{example}[Boolean cubes]
Let $L=\mathcal P(I)$, ordered by inclusion.  Then the ring-null fibre of
$G_L(R)$ is the Boolean cube of support masks, addition is union, and
multiplication is intersection.  Localization at a mask $A\subseteq I$
retains exactly the subcube $\mathcal P(A)$.
\end{example}

\begin{example}[The iterated sub-zero chain]
Let
\[
C_n=\{0<1<\cdots<n\}.
\]
In $G_{C_n}(R)$ define
\[
\tau_i:=(0,n-i),
\qquad 1\leq i\leq n.
\]
Then
\[
\tau_i+\tau_j=\tau_{\min(i,j)},
\qquad
\tau_i\tau_j=\tau_{\max(i,j)},
\]
while $r+\tau_i=r$ and $r\tau_i=\tau_i$ for $r\in R$.  Thus
$G_{C_n}(R)$ is exactly the $n$-fold sub-zero globalization
\[
R\sqcup\{\tau_1,\dots,\tau_n\}.
\]
The case $n=1$ is the ordinary split-zero construction.  Corollary~\ref{cor:chain-threshold-reconstruction} gives its faithful Boolean-coordinate realization at every length.
\end{example}

\begin{theorem}[Nilpotent-thickening realization of the support chain]
\label{thm:nilpotent-chain-realization}
Let $n\geq1$.
\begin{enumerate}[label=\textup{(\arabic*)}]
\item For every prime $p$, let
\[
B_{p,n}=\Z/p^n\Z,
\qquad
I_a=p^aB_{p,n}
\quad(0\leq a\leq n).
\]
Then
\[
I_a+I_b=I_{\min(a,b)},
\qquad
I_a\cap I_b=I_{\max(a,b)}.
\]
The map
\[
C_n\longrightarrow\{I_0,\dots,I_n\},
\qquad
j\longmapsto I_{n-j},
\]
is a semiring isomorphism from $(C_n,\vee,\wedge)$ to the ideal chain
with addition given by ideal sum and multiplication by intersection.
\item For every field $k$, the same statement holds for
\[
B_{\epsilon,n}=k[\epsilon]/(\epsilon^n),
\qquad
J_a=\epsilon^aB_{\epsilon,n}.
\]
\end{enumerate}
Consequently, the zero fibre of the $n$-fold sub-zero globalization is
canonically isomorphic, after choosing its length, to the support semiring
of either a $p^n$-thickening or an infinitesimal $\epsilon^n$-thickening.
\end{theorem}

\begin{proof}
Both rings are principal ideal rings with the displayed chains.  If
$a\leq b$, then $I_b\subseteq I_a$ and $J_b\subseteq J_a$.  Hence the
sum of two ideals is the larger one and their intersection is the smaller
one, which gives the formulas.  Under $j\mapsto I_{n-j}$ or
$j\mapsto J_{n-j}$, join in $C_n$ becomes ideal sum and meet becomes
intersection.
\end{proof}

\begin{proposition}[The valuation-sensitive shadow]
\label{prop:truncated-valuation-semiring}
On the same ideal sets, use ordinary ideal multiplication instead of
intersection.  Then
\[
I_aI_b=I_{\min(n,a+b)},
\qquad
J_aJ_b=J_{\min(n,a+b)}.
\]
Thus both ideal semirings are isomorphic to
\[
V_n=\{0,1,\dots,n\},
\qquad
a\oplus b=\min(a,b),
\qquad
a\odot b=\min(n,a+b),
\]
with additive zero $n$ and multiplicative unit $0$.  The lattice semiring
of Theorem~\ref{thm:nilpotent-chain-realization} records support strata;
$V_n$ records order of vanishing.  They agree on addition and differ on
multiplication.
\end{proposition}

\begin{proof}
The product of principal ideals generated by $p^a,p^b$ is generated by
$p^{a+b}$, with $p^n=0$ in $B_{p,n}$.  The identical calculation holds
for powers of $\epsilon$.  The displayed formulas give the claimed
semiring isomorphisms.
\end{proof}

\begin{remark}
The preceding theorems identify support and valuation semirings, not the
ambient rings.  A characteristic-zero infinitesimal thickening and a
positive-characteristic $p$-power thickening need not be isomorphic.  What
is common is the finite filtration geometry seen by their ideal chains.
\end{remark}


\subsection{Ideal theory and the spectral ordinal sum}

For $\lambda\in L$ write $z_\lambda=(0,\lambda)$ and put
\[
Z_L:=\{z_\lambda:\lambda\in L\}=p_L^{-1}(0_R).
\]
A lattice ideal is a lower subset closed under finite joins.  It is prime
when it is proper and
$\lambda\wedge\mu\in\mathfrak p$ implies
$\lambda\in\mathfrak p$ or $\mu\in\mathfrak p$.

\begin{theorem}[Complete ideal classification]
\label{thm:lattice-ideal-classification}
Every ideal of $G_L(R)$ is of exactly one of the following forms:
\begin{enumerate}[label=\textup{(\roman*)}]
\item a pure-support ideal
\[
Z_{\mathfrak a}:=\{z_\lambda:\lambda\in\mathfrak a\},
\]
where $\mathfrak a\subsetneq L$ is a proper lattice ideal;
\item an amplitude ideal
\[
I_J:=Z_L\cup\{(r,1_L):r\in J\},
\]
where $J\triangleleft R$ is a ring ideal.
\end{enumerate}
The two families are disjoint, and $I_{(0)}=Z_L$.
\end{theorem}

\begin{proof}
Let $I\subseteq G_L(R)$ be an ideal.  Suppose first that $I$ contains a
supported element $(r,1_L)$.  If $r\neq0$, multiplication by $-1$ and
addition show that $e_L=(0,1_L)$ belongs to $I$; this is already true when
$r=0$.  Since
\[
e_Lz_\lambda=z_\lambda,
\]
all of $Z_L$ lies in $I$.  The supported amplitudes in $I$ form a ring
ideal $J$: multiplication by supported scalars gives absorption, and
multiplication by $-1$ gives additive inverses.  Hence $I=I_J$.

Suppose now that $I$ contains no supported element.  Then $I\subseteq Z_L$.
Put
\[
\mathfrak a=\{\lambda\in L:z_\lambda\in I\}.
\]
Closure under addition makes $\mathfrak a$ join-closed.  If
$\mu\leq\lambda\in\mathfrak a$, then
$z_\mu=z_\mu z_\lambda\in I$, so $\mathfrak a$ is lower.  It is proper,
because $1_L\in\mathfrak a$ would place $e_L$ in $I$.  Conversely, both
displayed constructions are immediately ideals.  Their forms determine
the parameters uniquely.
\end{proof}

\begin{theorem}[Prime classification]
\label{thm:lattice-prime-classification}
The prime ideals of $G_L(R)$ are exactly
\[
P_{\mathfrak p}:=Z_{\mathfrak p}
\quad(\mathfrak p\in\SpecLat(L))
\]
and
\[
Q_{\mathfrak q}:=I_{\mathfrak q}
\quad(\mathfrak q\in\Spec R).
\]
Their inclusions are
\begin{align*}
P_{\mathfrak p}\subseteq P_{\mathfrak p'}
&\iff \mathfrak p\subseteq\mathfrak p',\\
Q_{\mathfrak q}\subseteq Q_{\mathfrak q'}
&\iff \mathfrak q\subseteq\mathfrak q',\\
P_{\mathfrak p}&\subsetneq Q_{\mathfrak q}
\qquad\text{for every }\mathfrak p,\mathfrak q.
\end{align*}
\end{theorem}

\begin{proof}
For a proper lattice ideal $\mathfrak p$, a product of two supported
amplitudes has support $1_L$ and therefore cannot lie in
$P_{\mathfrak p}$.  A product involving one zero-fibre element belongs to
$P_{\mathfrak p}$ exactly when that zero-fibre factor does.  For two
zero-fibre elements, primality is precisely
\[
\lambda\wedge\mu\in\mathfrak p
\Longrightarrow
\lambda\in\mathfrak p\text{ or }\mu\in\mathfrak p.
\]
Thus $P_{\mathfrak p}$ is prime exactly when $\mathfrak p$ is.

The ideal $Q_{\mathfrak q}$ contains the full zero fibre.  Products with a
zero-fibre factor therefore satisfy primality automatically, while products
of supported amplitudes reduce to the ordinary prime-ideal condition in
$R$.  The ideal classification proves that these exhaust the prime ideals.
The inclusion statements follow directly from the definitions.
\end{proof}

For topological spaces $X$ and $Y$, write $X\blacktriangleright Y$ for the
disjoint union endowed with the topology whose open sets are the opens of
$X$ together with the sets $X\cup V$ for $V\subseteq Y$ open.  Equivalently,
every point of $X$ specializes to every point of $Y$.

\begin{theorem}[Spectral ordinal-sum theorem]
\label{thm:spectral-ordinal-sum}
There is a canonical homeomorphism
\[
\boxed{
\Spec G_L(R)
\cong
\SpecLat(L)\blacktriangleright\Spec R.
}
\]
Under this homeomorphism,
\[
D(e_L)=\SpecLat(L),
\qquad
V(e_L)=\Spec R.
\]
Thus the lattice spectrum is an open stratum, the ring spectrum is a closed
stratum, and every lattice prime is a generization of every ring prime.
\end{theorem}

\begin{proof}
Theorem~\ref{thm:lattice-prime-classification} gives the underlying
bijection.  For $\lambda\in L$,
\[
D(z_\lambda)
=\{P_{\mathfrak p}:\lambda\notin\mathfrak p\},
\]
because every amplitude prime contains the full zero fibre.  These are
exactly the basic opens of $\SpecLat(L)$.  For a supported element
$r\in R$,
\[
D((r,1_L))
=\SpecLat(L)\cup D_R(r),
\]
because no pure-support prime contains a supported element.  These two
families generate precisely the ordinal-sum topology.  The identities for
$e_L=z_{1_L}$ follow because a proper lattice ideal omits $1_L$, whereas
every amplitude prime contains $Z_L$.
\end{proof}

\begin{corollary}[Krull dimension]
\label{cor:lattice-split-dimension}
Assume that $L$ is finite and that $\Spec R$ has finite Krull dimension.
Then
\[
\boxed{
\dim G_L(R)=\dim\SpecLat(L)+\dim R+1.
}
\]
In particular, the ordinary split-zero case $L=\B$ satisfies
$\dim G(R)=\dim R+1$.
\end{corollary}

\begin{proof}
Every chain of prime ideals is an initial chain of lattice primes followed
by a terminal chain of ring primes, by
Theorem~\ref{thm:lattice-prime-classification}.  Conversely, any maximal
chain in $\SpecLat(L)$ concatenates with any maximal chain in $\Spec R$,
with one additional strict inclusion between the two strata.  Taking
maximal lengths gives the formula.
\end{proof}


\section{Pre-Boolean support and the \texorpdfstring{$\Fnot$}{F-not-2} principle}
\label{sec:fnot2-principle}

The notation $\Fone$ is useful when the multiplicative monoid and its
Frobenius endomorphisms are the load-bearing data.  The support theory of
this paper uses a different coefficient layer.  Before support is
collapsed to the two values $0$ and $1$, it is a lattice, frame, locale, or
measure algebra.  This section isolates that layer.

\begin{definition}[Pre-Boolean support coefficient]
\label{def:fnot2-coefficient}
For a bounded distributive lattice $L$, define
\[
\boxed{\Fnot(L):=(L,\vee,\wedge).}
\]
This is an idempotent semiring with additive zero $0_L$ and multiplicative
unit $1_L$.  A \emph{Boolean branch} of $L$ is a bounded lattice
homomorphism
\[
\epsilon:L\longrightarrow\B.
\]
The notation $\Fnot$ is a mnemonic for support before the two-valued
quotient.  It is not a field notation and no theorem below uses it as one.
\end{definition}

\begin{remark}[Terminology]
The phrase ``not two'' is used only to name the formal condition that no
Boolean branch has yet been selected.  Informally one may call this an
unbranched or undifferentiated support state.  The formal object is the
lattice or frame $L$, together with its spectrum of branches.
\end{remark}

\begin{theorem}[Booleanization is branch selection]
\label{thm:booleanization-prime-branches}
For a bounded distributive lattice $L$, the following sets are naturally
identified:
\[
\Hom_{\mathrm{DLat}_{0,1}}(L,\B),
\qquad
\SpecLat(L),
\qquad
\{\text{prime filters of }L\}.
\]
Under this identification, a branch $\epsilon:L\to\B$ corresponds to the
prime ideal
\[
\mathfrak p_\epsilon=\epsilon^{-1}(0)
\]
and to the complementary prime filter
\[
F_\epsilon=\epsilon^{-1}(1).
\]
\end{theorem}

\begin{proof}
If $\epsilon:L\to\B$ is a bounded lattice homomorphism, then
$\mathfrak p_\epsilon$ is a proper lower set closed under joins.  If
$\lambda\wedge\mu\in\mathfrak p_\epsilon$, then
\[
0=\epsilon(\lambda\wedge\mu)=\epsilon(\lambda)\wedge\epsilon(\mu),
\]
so at least one of $\epsilon(\lambda),\epsilon(\mu)$ is $0$.  Thus
$\mathfrak p_\epsilon$ is prime.  Its complement is a prime filter.

Conversely, if $\mathfrak p\subset L$ is a prime lattice ideal, define
\[
\epsilon_{\mathfrak p}(\lambda)=
\begin{cases}
0,&\lambda\in\mathfrak p,\\
1,&\lambda\notin\mathfrak p.
\end{cases}
\]
The ideal axioms give preservation of $0,1$ and joins, and primality gives
preservation of meets.  These two constructions are inverse.
\end{proof}

Theorem~\ref{thm:lattice-boolean-characters} can now be read as follows:
Boolean-valued split-support characters do not choose an amplitude; they
choose a support branch.

\begin{corollary}[Boolean split characters are support branches]
\label{cor:boolean-split-characters-branches}
For every nonzero commutative ring $R$ and nontrivial bounded distributive
lattice $L$,
\[
\boxed{
\Hom_{\SemiRing}(G_L(R),\B)
\cong
\SpecLat(L).
}
\]
Every such homomorphism is the composite
\[
G_L(R)\xrightarrow{\chi_L}L\xrightarrow{\epsilon}\B
\]
on the zero fibre, together with the rule that every supported amplitude
maps to $1$.
\end{corollary}

\begin{proof}
Combine Theorem~\ref{thm:lattice-boolean-characters} with
Theorem~\ref{thm:booleanization-prime-branches}.
\end{proof}

\begin{definition}[Boolean branch base change]
\label{def:boolean-branch-base-change}
Let $\epsilon:L\to\B$ be a Boolean branch.  Define
\[
\beta_\epsilon:G_L(R)\longrightarrow G_{\B}(R)
\]
by
\[
\beta_\epsilon(0,\lambda)=(0,\epsilon(\lambda)),
\qquad
\beta_\epsilon(r,1_L)=(r,1)
\quad(r\in R).
\]
\end{definition}

\begin{theorem}[Universal property of Boolean branch base change]
\label{thm:boolean-branch-base-change}
The map $\beta_\epsilon$ is a semiring homomorphism.  It is universal
among semiring homomorphisms $f:G_L(R)\to T$ equipped with a factorization
of the zero-fibre support through $\epsilon$ and an amplitude-preserving
map from the supported copy of $R$.

Equivalently, $G_{\B}(R)$ is obtained from $G_L(R)$ by imposing exactly the
relations
\[
z_\lambda\sim z_\mu
\quad\text{whenever}\quad
\epsilon(\lambda)=\epsilon(\mu),
\]
while leaving the supported copy of $R$ unchanged.
\end{theorem}

\begin{proof}
The formula is compatible with addition because
\[
\epsilon(\lambda\vee\mu)=\epsilon(\lambda)\vee\epsilon(\mu),
\]
and compatible with multiplication because
\[
\epsilon(\lambda\wedge\mu)=\epsilon(\lambda)\wedge\epsilon(\mu).
\]
The interactions with fully supported amplitudes follow from
$\epsilon(1_L)=1$.  This proves that $\beta_\epsilon$ is a semiring
homomorphism.

Let $q:G_L(R)\to G_\B(R)$ be the displayed quotient.  Any homomorphism
which factors the zero fibre through $\epsilon$ must identify precisely the
zero-fibre elements with the same $\epsilon$-value; amplitude preservation
forces the same value on every supported element.  Hence it factors
uniquely through $\beta_\epsilon$.
\end{proof}

\begin{theorem}[The branch-selected arithmetic spectrum]
\label{thm:branch-selected-spectrum}
Let $\epsilon:L\to\B$ be a Boolean branch and set
$\mathfrak p_\epsilon=\epsilon^{-1}(0)$.  The induced spectral map
\[
\Spec G_{\B}(R)\longrightarrow\Spec G_L(R)
\]
has image
\[
\boxed{
\{P_{\mathfrak p_\epsilon}\}\blacktriangleright\Spec R.
}
\]
More explicitly,
\[
P_\tau\longmapsto P_{\mathfrak p_\epsilon},
\qquad
Q_{\mathfrak q}\longmapsto Q_{\mathfrak q}
\quad(\mathfrak q\in\Spec R).
\]
For $R=\Z$, every branch therefore determines the chain
\[
P_{\mathfrak p_\epsilon}
\subset
Q_{(0)}
\subset
Q_{(p)}
\qquad(p\text{ prime}),
\]
which is the Boolean branch followed by the ordinary arithmetic
specialization chain.
\end{theorem}

\begin{proof}
The prime $P_\tau$ of $G_\B(R)$ is the zero-fibre ideal consisting only of
its actual zero.  Its preimage under $\beta_\epsilon$ is
\[
\{z_\lambda:\epsilon(\lambda)=0\}=P_{\mathfrak p_\epsilon}.
\]
For a ring prime $\mathfrak q\subset R$, the amplitude prime of
$G_\B(R)$ contains the whole Boolean zero fibre and the supported
amplitudes with values in $\mathfrak q$.  Its preimage is therefore the
whole zero fibre $Z_L$ together with the supported amplitudes in
$\mathfrak q$, namely $Q_{\mathfrak q}$.  The specialization statements are
Theorem~\ref{thm:lattice-prime-classification}.
\end{proof}

\begin{corollary}[The three coefficient levels]
\label{cor:three-coefficient-levels}
The constructions in this paper separate three coefficient levels:
\[
\boxed{
\Fone
\quad\longrightarrow\quad
\Fnot(L)
\quad\xrightarrow{\epsilon}\quad
\B.
}
\]
Here $\Fone$ records multiplicative generators and Frobenius maps,
$\Fnot(L)$ records unresolved support, and $\B$ records a selected
two-valued support branch.  The scalar split-zero semiring $G(R)$ is the
case $L=\B$, where support has already been Booleanized.
\end{corollary}

\begin{proof}
The first level is the monoidal or blueprint coefficient layer used for
absolute and spherical algebra.  The second level is Definition
\ref{def:fnot2-coefficient}.  The third is obtained from the second by a
branch $\epsilon$ and Theorem~\ref{thm:boolean-branch-base-change}.  For
$L=\B$, the only nontrivial branch is the identity, so the pre-Boolean and
Boolean support layers coincide.
\end{proof}

\begin{proposition}[Locale-valued support]
\label{prop:locale-valued-support}
If $L=\mathcal O(X)$ is the frame of open sets of a topological space and
branches are required to preserve arbitrary joins, then branches
$\mathcal O(X)\to\B$ are exactly completely prime filters of opens.  If
$X$ is sober, these are the neighbourhood filters of points of $X$.
Thus $\Fnot(\mathcal O(X))$ is the support locale before evaluation at a
point.
\end{proposition}

\begin{proof}
A frame homomorphism to $\B$ has inverse image of $1$ a filter closed under
finite meets and completely prime with respect to arbitrary joins.  The
converse is the characteristic map of a completely prime filter.  The final
statement is the standard definition of sobriety.
\end{proof}

\begin{theorem}[Unsplit and branch states for a compact symmetry]
\label{thm:haar-dirac-branch}
Let $G$ be a compact Hausdorff group and let $G$ act on the compact convex
space $P(G)$ of Borel probability measures by left translation.  In the
action groupoid $P(G)//G$, Haar measure $m_G$ has stabilizer $G$, while
each Dirac measure $\delta_g$ has trivial stabilizer.  Hence Haar measure
is an unsplit invariant state and the Dirac measures are branch states.
\end{theorem}

\begin{proof}
For every $h\in G$, left invariance gives $h_*m_G=m_G$, so the stabilizer
of $m_G$ is all of $G$.  If $h_*\delta_g=\delta_g$, then
$\delta_{hg}=\delta_g$, hence $hg=g$ and $h=1$.
\end{proof}

\begin{corollary}[Bost--Connes branch interpretation]
\label{cor:bost-connes-unsplit-branch}
For the cyclotomic symmetry group $\widehat\Z^\times$, the invariant KMS
state in the unbroken Bost--Connes regime is the dynamical analogue of an
unsplit support state, while the low-temperature extremal KMS states are
branch states carrying the cyclotomic Galois action.  The pole of
$\zeta(\beta)$ at $\beta=1$ is the analytic boundary between these two
regimes.
\end{corollary}

\begin{proof}
The Bost--Connes classification gives a unique symmetry-invariant KMS
state for $0<\beta\leq1$ and a simplex of extremal KMS states with
cyclotomic symmetry action for $\beta>1$.  Apply
Theorem~\ref{thm:haar-dirac-branch} to the compact cyclotomic symmetry
acting on the branch space.
\end{proof}

\section{Nonassociative coefficients and cocycle-controlled associators}
\label{sec:nonassociative-support}

The support construction does not require associative coefficients.  This
provides a precise interface with Cayley--Dickson algebras and with
nonassociative algebras internal to tensor categories.

\begin{definition}[Lattice-supported algebra]
Let $k$ be a commutative ring and let $A$ be a unital, not necessarily
associative, $k$-algebra.  Define
\[
G_L(A)=
\{(0,\lambda):\lambda\in L\}
\cup
\{(a,1_L):a\in A\}
\]
with componentwise addition and
\[
(a,\lambda)(b,\mu)=(ab,\lambda\wedge\mu).
\]
\end{definition}

\begin{theorem}[Separation of support and associator]
\label{thm:support-associator-separation}
For $x=(a,\lambda)$, $y=(b,\mu)$, and $z=(c,\nu)$ in $G_L(A)$,
\begin{align*}
(xy)z&=((ab)c,\lambda\wedge\mu\wedge\nu),\\
x(yz)&=(a(bc),\lambda\wedge\mu\wedge\nu).
\end{align*}
Therefore:
\begin{enumerate}[label=\textup{(\arabic*)}]
\item $G_L(A)$ is associative if and only if $A$ is associative;
\item alternativity, flexibility, and power associativity pass from $A$
to $G_L(A)$ and are reflected by the fully supported copy of $A$;
\item if $a,b\neq0$ and $ab=0$ in $A$, then
\[
(a,1_L)(b,1_L)=e_L=(0,1_L)\neq\tau_L;
\]
thus intrinsic zero-divisor cancellation is supported, whereas multiplication
by an absent factor produces $\tau_L$.
\end{enumerate}
\end{theorem}

\begin{proof}
The two displayed formulas are immediate from componentwise multiplication
and associativity of the meet operation.  Every polynomial identity in the
multiplication having the same multiset of variables on both sides has the
same support on both sides, so any such identity valid in $A$ remains valid
in $G_L(A)$.  The fully supported embedding $a\mapsto(a,1_L)$ reflects
failures of these identities.  The final assertion follows directly from
the multiplication law and from $e_L\neq\tau_L$.
\end{proof}

\begin{definition}[Twisted graded coefficient algebra]
Let $\Gamma$ be a group, let $k$ be a field, and let
\[
F:\Gamma\times\Gamma\longrightarrow k^\times
\]
be a normalized $2$-cochain.  On the vector space with basis
$\{u_g:g\in\Gamma\}$ define
\[
u_gu_h=F(g,h)u_{gh}.
\]
Set
\[
\phi_F(g,h,k)
=
\frac{F(g,h)F(gh,k)}{F(h,k)F(g,hk)}.
\]
\end{definition}

\begin{theorem}[Associator cocycle and internal associativity]
\label{thm:associator-cocycle}
The function $\phi_F$ is a normalized $3$-cocycle; with the displayed
convention it is the coboundary of $F$.  On homogeneous basis vectors,
\[
(u_gu_h)u_k
=\phi_F(g,h,k)\,u_g(u_hu_k).
\]
Consequently, the twisted algebra is associative precisely when
$\phi_F=1$.  For arbitrary $F$ it is an associative algebra object in the
monoidal category of $\Gamma$-graded vector spaces whose associator on
degrees $(g,h,k)$ is multiplication by $\phi_F(g,h,k)$.

Moreover, the lattice-supported algebra $G_L(k_F[\Gamma])$ has the same
basis associator cocycle $\phi_F$; the support meet contributes no
additional associator.
\end{theorem}

\begin{proof}
The displayed associator identity follows by comparing
\[
F(g,h)F(gh,k)u_{ghk}
\quad\text{and}\quad
F(h,k)F(g,hk)u_{ghk}.
\]
The pentagon identity for the monoidal associator is exactly the
$3$-cocycle equation for $\phi_F$; equivalently, $\phi_F$ is the
cochain coboundary of $F$ in the stated convention.  The algebra-object
associativity diagram is the displayed identity.  In the
lattice-supported extension, both parenthesizations have support
$\lambda\wedge\mu\wedge\nu$, so the scalar associator remains
$\phi_F$.
\end{proof}

\begin{corollary}[A three-axis coefficient package]
\label{cor:support-phase-associator-package}
The object
\[
\mathscr A(L,\Gamma,F):=G_L(k_F[\Gamma])
\]
separates three kinds of information:
\[
\begin{array}{c|c}
L&\text{support and definedness},\\
\Gamma&\text{phase or grading sector},\\
\phi_F\in Z^3(\Gamma;k^\times)&\text{pointwise associator data}.
\end{array}
\]
Ring reflection forgets only the support lattice, localization restricts
support to a principal interval, and neither operation changes the
supported associator cochain.
\end{corollary}

\begin{remark}[Cohomology class versus chosen multiplication]
Because a global multiplication cochain $F$ has been chosen,
$\phi_F$ is exact in ordinary group cohomology.  Its pointwise values can
nevertheless be nontrivial and encode the failure of the chosen basis
multiplication to associate.  A genuinely nonzero class in
$H^3(\Gamma;k^\times)$ arises when the monoidal category carries a
background associator which cannot be globally strictified by such an
$F$.
\end{remark}

\begin{theorem}[Cayley--Dickson support threshold]
\label{thm:cayley-dickson-support-threshold}
Let $\CD_r$ denote the standard real Cayley--Dickson algebra of dimension
$2^r$, so that
\[
\CD_0=\R,
\quad
\CD_1=\C,
\quad
\CD_2=\mathbb H,
\quad
\CD_3=\mathbb O,
\quad
\CD_4=\mathbb S
\]
with $\mathbb S$ the sedenions.  Then:
\begin{enumerate}[label=\textup{(\arabic*)}]
\item for $r\leq3$, no two nonzero fully supported elements of
$G_L(\CD_r)$ multiply to $e_L$;
\item for every $r\geq4$, there exist nonzero fully supported elements
$x,y\in G_L(\CD_r)$ with $xy=e_L$;
\item in a signed graded Cayley--Dickson basis indexed by
$(\Z/2\Z)^r$, the support completion preserves the basis associator
$3$-cocycle.
\end{enumerate}
Thus the octonionic loss of associativity and the sedenionic appearance of
zero divisors are distinct events, and split support records the second
without conflating it with absence.
\end{theorem}

\begin{proof}
The algebras $\R,\C,\mathbb H,\mathbb O$ are division algebras, so the
first assertion follows from
Theorem~\ref{thm:support-associator-separation}.  The sedenions possess
nontrivial zero divisors, and every later Cayley--Dickson algebra contains
the preceding one through the standard embedding $a\mapsto(a,0)$.  Hence
sedenion zero divisors give the second assertion at every later stage.
A signed graded basis has multiplication controlled by a
$\{\pm1\}$-valued $2$-cochain, so the third assertion is
Theorem~\ref{thm:associator-cocycle}.
\end{proof}

\begin{remark}[Two independent defect coordinates]
Theorem~\ref{thm:cayley-dickson-support-threshold} isolates two independent
invariants: the support coordinate in $L$ and zero-divisibility in the
coefficient algebra.  An intrinsic coefficient zero-divisor pair lands at
$e_L$, while combinatorial absence lands at $\tau_L$.
\end{remark}

\section{The split spectrum and the \texorpdfstring{$\mathbb F_1$}{F1} support skeleton}

\begin{theorem}[Ideals and prime ideals]\label{thm:prime-classification}
Let $R$ be a nonzero commutative ring and $S=G(R)$.
\begin{enumerate}[label=\textup{(\alph*)}]
\item Every ideal of $S$ is either $\{\tau\}$ or
\[
I_J:=J\cup\{\tau\}
\]
for a unique ideal $J\triangleleft R$.
\item The prime ideals are
\[
P_\tau:=\{\tau\}
\]
and
\[
P_{\mathfrak p}:=\mathfrak p\cup\{\tau\},
\qquad
\mathfrak p\in\Spec R.
\]
\end{enumerate}
\end{theorem}

\begin{proof}
Every ideal contains $\tau$.  If $I\cap R=\varnothing$, then $I=\{\tau\}$.  Otherwise $J:=I\cap R$ is closed under addition, multiplication by arbitrary elements of $R$, and additive inverses because $-a=(-1)a$.  Hence $J\triangleleft R$ and $I=J\cup\{\tau\}$.  Conversely every such set is an ideal.

The ideal $P_\tau$ is prime because a product in $S$ equals $\tau$ only when one factor equals $\tau$.  For $I_J$, the quotient by the corresponding ring-reflection congruence has supported part $R/J$; equivalently, for supported $a,b$,
\[
ab\in I_J
\iff
ab\in J.
\]
Thus $I_J$ is prime exactly when $J$ is prime.
\end{proof}

\begin{theorem}[Generic-point suspension]\label{thm:generic-suspension}
The map
\[
i:\Spec R\longrightarrow\Spec G(R),
\qquad
\mathfrak p\longmapsto P_{\mathfrak p},
\]
is a homeomorphism onto the closed subspace
\[
V(e)=\{P_{\mathfrak p}:\mathfrak p\in\Spec R\}.
\]
Furthermore,
\[
D(e)=\{P_\tau\},
\]
and $P_\tau$ is open, dense, and contained in every prime ideal.  Hence $\Spec G(R)$ is obtained from $\Spec R$ by adjoining a new open dense generic point below every classical prime in the specialization order.
\end{theorem}

\begin{proof}
Since $e=0_R$ belongs to every ring prime $\mathfrak p$ but does not belong to $P_\tau$, the displayed descriptions of $V(e)$ and $D(e)$ follow from Theorem \ref{thm:prime-classification}.  For an ideal $J\triangleleft R$,
\[
i(V_R(J))
=
V_S(I_J)\cap V_S(e),
\]
so $i$ identifies the Zariski closed sets and is a homeomorphism onto $V(e)$.

The containment $P_\tau\subseteq P$ holds for every prime $P$ because every ideal contains the semiring zero $\tau$.  Therefore the closure of $\{P_\tau\}$ is all of $\Spec S$.  Since $D(e)=\{P_\tau\}$, this point is also open.
\end{proof}

\begin{corollary}[Local stalks]\label{cor:split-stalks}
For $\mathfrak p\in\Spec R$,
\[
G(R)_{P_\tau}\cong\B,
\qquad
G(R)_{P_{\mathfrak p}}\cong G(R_{\mathfrak p}).
\]
\end{corollary}

\begin{proof}
The complement of $P_\tau$ is the whole supported part $R$, which contains $e$, so Corollary \ref{cor:fraction-comparison} gives $\B$.  The complement of $P_{\mathfrak p}$ is $R\setminus\mathfrak p$, which does not contain $0$; Theorem \ref{thm:ring-localization} gives $G(R_{\mathfrak p})$.
\end{proof}

\begin{example}[The shifted generic point of $\Spec\Z$]
For $S=G(\Z)$,
\[
P_\tau=\{\tau\}
\subsetneq
P_0=\{\tau,0\}
\subsetneq
P_p=p\Z\cup\{\tau\}
\]
for every rational prime $p$.  Thus the classical zero ideal $(0)$ becomes the nonzero prime $P_0$, while $P_\tau$ is the new zero ideal and new generic point.
\end{example}

\subsection{The exact blueprint presentation}

Let $M_S$ be the multiplicative monoid of $S$, including its absorbing element $\tau$.  Let $\N_0[M_S]$ denote the contracted monoid semiring in which the basis element $[\tau]$ is identified with the additive zero.

\begin{theorem}[Multiplicative skeleton plus preaddition]\label{thm:blueprint-presentation}
The canonical map
\[
\N_0[M_S]\longrightarrow S,
\qquad
[x]\longmapsto x,
\]
induces an isomorphism
\[
\N_0[M_S]\Big/\Big\langle [a]+[b]\sim[a+b]
\;:\;a,b\in R\Big\rangle
\xrightarrow{\sim}
G(R).
\]
Consequently the associated blueprint consists of the multiplicative monoid $M_S$ together with precisely the additive relations inherited from $R$ and the contraction of $\tau$ to additive zero.
\end{theorem}

\begin{proof}
The canonical map is surjective.  The imposed relations hold in $G(R)$, so the map factors through the quotient.

Every element of the quotient is represented by a finite sum of monoid generators.  Terms $[\tau]$ vanish in the contracted monoid semiring.  If at least one supported term remains, repeated use of
\[
[a]+[b]\sim[a+b]
\]
reduces the sum to a single generator $[r]$ with $r\in R$.  If no supported term remains, the sum is zero, represented by $[\tau]$.  Thus every congruence class has a representative in $M_S$.  Two such representatives have the same image in $G(R)$ only when they are equal, because the reduction relations reproduce the actual ring sum and never identify $\tau$ with a supported element.  Hence the induced map is bijective.
\end{proof}

\begin{remark}[The role of $\mathbb F_1$]
The monoid $M_S$ is the multiplicative $\mathbb F_1$ shadow in the sense of monoidal geometry, while Theorem \ref{thm:blueprint-presentation} gives the exact blueprint that restores additive cancellation data; see \cite{Deitmar,Lorscheid}.  The Boolean quotient $\chi_R:S\to\B$ is the pure support completion.  The projective geometry below makes this skeleton visible as a finite Boolean face poset whose fibres carry ordinary projective geometry.
\end{remark}



\part{Projective and invariant geometry}

\section{Conservative ideal arithmetic and the separated zeta sheet}
\label{sec:ideal-zeta}

Let $R$ be a commutative ring.  For an ideal $J\triangleleft R$, write
\[
I_J:=J\sqcup\{\tau\}\subseteq G(R).
\]
The baseline ideal classification in \cite{Globalization,Secondary} shows
that these are exactly the nontrivial supported ideals.

\begin{theorem}[Ideal multiplication is classical]
\label{thm:ideal-multiplication-classical}
For ideals $J,K\triangleleft R$,
\[
I_JI_K=I_{JK}.
\]
\end{theorem}

\begin{proof}
Every product of an element of $I_J$ with an element of $I_K$ is either
$\tau$ or a supported product $jk$, so $I_JI_K\subseteq I_{JK}$.  Conversely,
$\tau\in I_JI_K$, and every element of $JK$ is a finite sum of products
$j_ak_a$; the ideal product $I_JI_K$ contains each such product and is
closed under addition.  Hence $I_{JK}\subseteq I_JI_K$.
\end{proof}

Let $A=\mathcal O_K$ be the ring of integers of a number field.  Define the
\emph{normed split ideal class group} to be the group of nonzero ideals
$I_J$ modulo the principal supported ideals $I_{(a)}$.

\begin{corollary}[Class group and Dedekind zeta preservation]
\label{cor:dedekind-preservation}
The map $J\mapsto I_J$ induces canonical isomorphisms
\[
\operatorname{Cl}^{\mathrm{sp}}(G(A))\cong\operatorname{Cl}(A)
\]
and, on the ordinary quotient-norm sheet,
\[
\zeta^{\mathrm{ar}}_{G(A)}(s)
:=\sum_{0\ne J\triangleleft A}N(J)^{-s}
=\zeta_K(s).
\]
\end{corollary}

\begin{proof}
Theorem~\ref{thm:ideal-multiplication-classical} makes $J\mapsto I_J$ a
multiplicative bijection between nonzero ideals of $A$ and nontrivial
finite-norm supported ideals of $G(A)$; it carries $(a)$ to $I_{(a)}$.
This proves the class-group assertion.  The ordinary congruence quotient
collapsing $\tau$ with the supported zero is $A/J$, so its cardinality is
$N(J)$, giving the zeta identity term by term.
\end{proof}

The split quotient retaining the distinction between $\tau$ and $0_{A/J}$
is $G(A/J)$ and has cardinality $N(J)+1$.

\begin{theorem}[The separated weight is nonmultiplicative]
\label{thm:separated-weight}
For nonzero finite-norm ideals $J,K$ put
\[
w(J):=N(J)+1.
\]
Then
\[
w(JK)=N(J)N(K)+1
\]
while
\[
w(J)w(K)=N(J)N(K)+N(J)+N(K)+1.
\]
In particular $w$ is not multiplicative.
\end{theorem}

\begin{proof}
Use multiplicativity of the ordinary ideal norm and expand the product.
Since $N(J),N(K)>0$, the two expressions differ.
\end{proof}

\begin{corollary}[The support-separated sheet is non-Eulerian]
\label{cor:separated-non-eulerian}
The Dirichlet series
\[
\zeta^{\mathrm{sep}}_{G(A)}(s)
:=\sum_{0\ne J\triangleleft A}(N(J)+1)^{-s}
\]
has no ordinary Euler product indexed by the prime ideals of $A$.
For $A=\Z$,
\[
\zeta^{\mathrm{sep}}_{G(\Z)}(s)=\sum_{n\geq1}(n+1)^{-s}=\zeta(s)-1.
\]
\end{corollary}

\begin{proof}
An ordinary ideal Euler product gives multiplicative Dirichlet
coefficients on coprime ideals.  Theorem~\ref{thm:separated-weight} rules
this out.  The final identity is the change of index $m=n+1$.
\end{proof}

\begin{remark}
The two series answer different questions.  The normed sheet counts the
classical arithmetic quotient and preserves the Euler product.  The
separated sheet counts the extra unsupported class as a genuine point and
therefore destroys multiplicativity.  The absence of an Euler product is
not a defect in the calculation; it is the exact analytic cost of retaining
split support in quotient cardinalities.
\end{remark}

\section{Projective semimodules and finite distributive support}
\label{sec:projective-semimodules}

The semimodule equivalence of \cite{Secondary} sends a $G(R)$-semimodule
$M$ to its support join-semilattice $L_M=eM$, its $R$-module fibres
$M_\ell$, and the transport maps $T_{\ell,\ell'}$ for
$\ell\leq\ell'$.  The following result is the projective consequence that
can be obtained without imposing an unproved normal form on all transport
data.

\begin{theorem}[Necessary structure of a finitely generated projective]
\label{thm:projective-support-structure}
Let $P$ be a finitely generated projective $G(R)$-semimodule.  Then:
\begin{enumerate}[label=\textup{(\roman*)}]
\item $L_P=eP$ is a finite distributive lattice;
\item every fibre $P_\ell$ is a finitely generated projective $R$-module;
\item every transport map
\[
T_{\ell,\ell'}:P_\ell\longrightarrow P_{\ell'}
\qquad(\ell\leq\ell')
\]
is a split monomorphism of $R$-modules.
\end{enumerate}
If $R$ is a field, the function
$\ell\mapsto\dim_R P_\ell$ is therefore order-preserving.
\end{theorem}

\begin{proof}
Choose $G(R)$-linear maps
\[
i:P\longrightarrow G(R)^n,
\qquad
r:G(R)^n\longrightarrow P,
\qquad
r i=\id_P.
\]
Applying multiplication by $e$ gives join-homomorphisms
\[
i_s:L_P\longrightarrow\mathcal P([n]),
\qquad
r_s:\mathcal P([n])\longrightarrow L_P,
\qquad
r_s i_s=\id_{L_P}.
\]
Thus $L_P$ is a finite join-retract of a Boolean lattice.  For
$x,y\in L_P$, define
\[
x\wedge y:=r_s\bigl(i_s(x)\cap i_s(y)\bigr).
\]
If $z\leq x,y$, then $i_s(z)\subseteq i_s(x)\cap i_s(y)$, hence
$z\leq x\wedge y$; conversely $x\wedge y\leq x,y$ by monotonicity of
$r_s$.  This is the meet in $L_P$.  Since $i_s$ and $r_s$ preserve joins,
distributivity follows from distributivity in the Boolean lattice:
\[
\begin{aligned}
x\wedge(y\vee z)
&=r_s\bigl(i_s(x)\cap(i_s(y)\cup i_s(z))\bigr)\\
&=(x\wedge y)\vee(x\wedge z).
\end{aligned}
\]
This proves (i).

Fix $\ell\in L_P$ and write $I=i_s(\ell)$.  The maps $i$ and $r$
restrict to $R$-linear maps
\[
P_\ell\longrightarrow R^I\longrightarrow P_\ell
\]
whose composite is the identity.  Hence $P_\ell$ is a direct summand of a
finite free $R$-module, proving (ii).

For $\ell\leq\ell'$, put $I=i_s(\ell)$ and $I'=i_s(\ell')$; then
$I\subseteq I'$.  Let
\[
t_{I,I'}:R^I\hookrightarrow R^{I'}
\]
be extension by zero and let $\pi_{I',I}$ be coordinate projection.  Define
\[
s_{\ell',\ell}
:=r_\ell\circ\pi_{I',I}\circ i_{\ell'}:P_{\ell'}\to P_\ell.
\]
Naturality of $i$ with respect to transport gives
$i_{\ell'}T_{\ell,\ell'}=t_{I,I'}i_\ell$.  Therefore
\[
s_{\ell',\ell}T_{\ell,\ell'}
=r_\ell\pi_{I',I}t_{I,I'}i_\ell
=r_\ell i_\ell
=\id_{P_\ell}.
\]
This proves (iii).  The dimension assertion is immediate.
\end{proof}

\begin{remark}[Projective support invariant]
\label{rem:projective-classification-boundary}
Theorem~\ref{thm:projective-support-structure} associates to a finitely
generated projective $G(R)$-semimodule a finite distributive support
lattice, projective $R$-module fibres, and split-monic transport maps.
These data form an isomorphism invariant and provide the exact projective
input used by the later support and localization constructions.
\end{remark}

\section{Unimodular split projective geometry}

For a semiring $A$ and $n\ge0$, a vector
\[
x=(x_0,\dots,x_n)\in A^{n+1}
\]
is called \emph{unimodular} if the ideal generated by its coordinates is the unit ideal.  Equivalently, there exist $a_0,\dots,a_n\in A$ such that
\[
a_0x_0+\cdots+a_nx_n=1.
\]

\begin{definition}[Unimodular projective point set]\label{def:unimodular-projective}
Define
\[
\Pum^n(A)
:=
\{x\in A^{n+1}:x\text{ is unimodular}\}/A^\times,
\]
where $A^\times$ acts by simultaneous scalar multiplication.
\end{definition}

This is the usual set of projective points over a field.  Over a general semiring it is deliberately denoted by $\Pum$ rather than identified without qualification with a scheme-theoretic $\Proj$.

Let $K$ be a field and set $S=G(K)$.  For $x\in S^{n+1}$ define
\[
\supp(x):=\{i:x_i\in K\}.
\]
This records supported coordinates, including coordinates equal to $e=0_K$.

\begin{lemma}[Unimodularity criterion]\label{lem:unimodularity-criterion}
A vector $x\in S^{n+1}$ is unimodular if and only if at least one coordinate belongs to $K^\times$.
\end{lemma}

\begin{proof}
If $x_i\in K^\times$, then $x_i^{-1}x_i=1$, so the coordinates generate the unit ideal.

Conversely, suppose every coordinate belongs to $\{\tau,e\}$.  Multiplication of such a coordinate by any scalar remains in $\{\tau,e\}$, and a finite sum of elements of $\{\tau,e\}$ again lies in $\{\tau,e\}$.  It therefore cannot equal $1$.  Hence the vector is not unimodular.
\end{proof}

The support character induces a map
\[
\Supp:\Pum^n(S)\longrightarrow\Pum^n(\B),
\qquad
[x_0:\cdots:x_n]
\longmapsto
[\chi_K(x_0):\cdots:\chi_K(x_n)].
\]
Since $\B^\times=\{1\}$,
\[
\Pum^n(\B)=\B^{n+1}\setminus\{0\}
\]
is canonically the set of nonempty subsets of $[n]_0=\{0,\dots,n\}$.

\begin{theorem}[Boolean-face decomposition]\label{thm:projective-face-decomposition}
For every field $K$ and $n\ge0$, there is a canonical decomposition
\[
\boxed{
\Pum^n(G(K))
\cong
\bigsqcup_{\varnothing\neq A\subseteq[n]_0}
\mathbf P^{|A|-1}(K).}
\]
Under this identification, the fibre of $\Supp$ over $A$ is $\mathbf P^{|A|-1}(K)$.
\end{theorem}

\begin{proof}
Fix a nonempty subset $A\subseteq[n]_0$.  A vector of support exactly $A$ has the form
\[
x_i\in K\quad(i\in A),
\qquad
x_i=\tau\quad(i\notin A).
\]
By Lemma \ref{lem:unimodularity-criterion}, it is unimodular exactly when the amplitude vector $(x_i)_{i\in A}\in K^A$ is nonzero.  The unit group of $S$ is $K^\times$, and simultaneous multiplication by a supported unit acts on the amplitude vector by the usual scalar action.  Hence the set of projective classes of support $A$ is
\[
(K^A\setminus\{0\})/K^\times
=\mathbf P^{|A|-1}(K).
\]
Distinct support masks cannot be related by multiplication by a supported unit, so the union is disjoint.  The description of the support fibre is immediate.
\end{proof}


\subsection{Lattice-valued support profiles}

For a bounded distributive lattice $L$ and a field $K$, every coordinate of
$G_L(K)$ has a support value in $L$; a coordinate of support below $1_L$
has zero amplitude.  For
$\boldsymbol\lambda=(\lambda_0,\dots,\lambda_n)\in L^{n+1}$ put
\[
T(\boldsymbol\lambda)=\{i:\lambda_i=1_L\}.
\]

\begin{theorem}[Projective profile decomposition]
\label{thm:lattice-projective-profile}
For every field $K$, nontrivial bounded distributive lattice $L$, and
$n\ge0$, there is a canonical decomposition
\[
\boxed{
\Pum^n(G_L(K))
\cong
\bigsqcup_{\varnothing\neq A\subseteq[n]_0}
\mathbf P^{|A|-1}(K)\times
(L\setminus\{1_L\})^{[n]_0\setminus A}.
}
\]
The component indexed by $A$ consists of support profiles whose top-support
set is $A$.
\end{theorem}

\begin{proof}
A vector in $G_L(K)^{n+1}$ is unimodular exactly when at least one
coordinate has nonzero amplitude: such a coordinate is a unit, while a
linear combination of zero-amplitude coordinates remains in the zero
fibre.  Fix a support profile with top-support set $A$.  The coordinates in
$A$ form a nonzero vector in $K^A$; the coordinates outside $A$ have
arbitrary support values in $L\setminus\{1_L\}$ and zero amplitude.  The
unit group of $G_L(K)$ is $K^\times$, acting by the ordinary scalar action
on $K^A$ and fixing the support profile.  Quotienting therefore gives
$\mathbf P^{|A|-1}(K)$ for each lower-support decoration.  Distinct profiles
cannot lie in the same orbit.
\end{proof}

The Boolean-face decomposition is the case $L=\B$, because then
$L\setminus\{1\}=\{0\}$.

\begin{corollary}[Amplitude and support projections]\label{cor:projective-projections}
There are canonical maps
\[
\Pum^n(G(K))\xrightarrow{\;p_*\;}\mathbf P^n(K)
\]
and
\[
\Pum^n(G(K))\xrightarrow{\;\Supp\;}\Pum^n(\B).
\]
The first forgets whether a zero coordinate was absent or supported; the second forgets all amplitudes and retains only the support face.
\end{corollary}

\begin{proof}
A unimodular split vector has at least one nonzero amplitude coordinate, so its coordinatewise ring reflection is a nonzero vector in $K^{n+1}$.  This defines $p_*$.  The support map was defined above.  Both are invariant under multiplication by $K^\times$.
\end{proof}

\subsection{The split projective line}

For $n=1$, Theorem \ref{thm:projective-face-decomposition} gives
\[
\Pum^1(G(K))
\cong
\mathbf P^1(K)\sqcup\{0_\tau,\infty_\tau\},
\]
where
\[
0_\tau:=[\tau:1],
\qquad
\infty_\tau:=[1:\tau].
\]
Inside the full-support stratum $\mathbf P^1(K)$ lie the supported zero and supported infinity
\[
0_e:=[e:1],
\qquad
\infty_e:=[1:e].
\]

\begin{proposition}[Four zero-pole states]\label{prop:four-zero-pole-states}
The ring-reflection map satisfies
\[
p_*(0_\tau)=p_*(0_e)=0,
\qquad
p_*(\infty_\tau)=p_*(\infty_e)=\infty.
\]
The support map satisfies
\[
\Supp(0_\tau)=[0:1],
\qquad
\Supp(\infty_\tau)=[1:0],
\]
while every point of the full-support stratum, including $0_e$ and $\infty_e$, maps to
\[
[1:1].
\]
\end{proposition}

\begin{proof}
This follows directly from
\[
p_K(\tau)=p_K(e)=0,
\qquad
\chi_K(\tau)=0,
\quad
\chi_K(e)=1.
\]
\end{proof}

Thus the Boolean projective line
\[
\Pum^1(\B)=\{[0:1],[1:1],[1:0]\}
\]
is a three-point incidence skeleton: absent zero, full-support torus, and absent infinity.

\section{Finite-field counts and split projective symmetry}

\subsection{Point counts and the support-projective zeta function}


\begin{theorem}[Lattice-profile finite-field count]
\label{thm:lattice-finite-field-count}
Let $L$ be a finite nontrivial bounded distributive lattice of cardinality
$\ell$, let $q$ be a prime power, and put $N=n+1$.  Then
\[
\boxed{
\#\Pum^n(G_L(\F_q))
=\sum_{r=1}^{N}\binom Nr(\ell-1)^{N-r}\frac{q^r-1}{q-1}
=\frac{(q+\ell-1)^N-\ell^N}{q-1}.
}
\]
If
\[
P_{N,\ell}(X)
:=\frac{(X+\ell-1)^N-\ell^N}{X-1}
=\sum_{j=0}^{N-1}a_{N,\ell,j}X^j,
\]
then
\[
a_{N,\ell,j}
=\sum_{r=j+1}^{N}\binom Nr(\ell-1)^{N-r}
\]
and the associated point-counting zeta function is
\[
\boxed{
Z_{L,n}(\F_q,T)
=\prod_{j=0}^{N-1}(1-q^jT)^{-a_{N,\ell,j}}.
}
\]
\end{theorem}

\begin{proof}
Theorem~\ref{thm:lattice-projective-profile} gives
$\binom Nr$ choices of top-support set of size $r$,
$(\ell-1)^{N-r}$ lower-support decorations, and
$(q^r-1)/(q-1)$ amplitude classes.  This proves the first sum.  Writing
$c=\ell-1$ and using
\[
\sum_{r=1}^{N}\binom Nr c^{N-r}(q^r-1)
=(q+c)^N-(c+1)^N
\]
gives the closed count.  Finally,
\[
P_{N,\ell}(X)
=\sum_{r=1}^{N}\binom Nr(\ell-1)^{N-r}
(1+X+\cdots+X^{r-1}),
\]
which yields the coefficient formula.  Substitution of $q^m$ and the
identity
$\sum_{m\ge1}(q^jT)^m/m=-\log(1-q^jT)$ give the zeta product.
\end{proof}

\begin{theorem}[Finite-field point count]\label{thm:finite-field-point-count}
For every prime power $q$ and every $n\ge0$,
\[
\#\Pum^n(G(\F_q))
=
\sum_{r=1}^{n+1}
\binom{n+1}{r}\frac{q^r-1}{q-1}
=
\frac{(q+1)^{n+1}-2^{n+1}}{q-1}.
\]
\end{theorem}

\begin{proof}
There are $\binom{n+1}{r}$ support faces of cardinality $r$, and the fibre over each such face is $\mathbf P^{r-1}(\F_q)$, which has $(q^r-1)/(q-1)$ points.  This gives the first formula.  For the second,
\[
\sum_{r=1}^{n+1}\binom{n+1}{r}(q^r-1)
=
\bigl((q+1)^{n+1}-1\bigr)-\bigl(2^{n+1}-1\bigr),
\]
which equals $(q+1)^{n+1}-2^{n+1}$.
\end{proof}

\begin{definition}[Support-projective zeta function]\label{def:support-zeta}
Define
\[
Z_{\mathrm{sp}}(\Pum^n/\F_q,T)
:=
\exp\left(
\sum_{m\ge1}
\#\Pum^n(G(\F_{q^m}))\frac{T^m}{m}
\right).
\]
This is the Hasse--Weil type point-counting series attached to the split projective functor $q\mapsto\Pum^n(G(\F_q))$.
\end{definition}

\begin{theorem}[Closed zeta product]\label{thm:support-zeta-product}
Put
\[
c_{n,j}:=\sum_{r=j+1}^{n+1}\binom{n+1}{r}
\qquad(0\le j\le n).
\]
Then
\[
\#\Pum^n(G(\F_q))
=
\sum_{j=0}^{n}c_{n,j}q^j,
\]
and
\[
\boxed{
Z_{\mathrm{sp}}(\Pum^n/\F_q,T)
=
\prod_{j=0}^{n}(1-q^jT)^{-c_{n,j}}.}
\]
In particular,
\[
Z_{\mathrm{sp}}(\Pum^1/\F_q,T)
=
\frac{1}{(1-T)^3(1-qT)}.
\]
\end{theorem}

\begin{proof}
Since
\[
\frac{q^r-1}{q-1}=1+q+\cdots+q^{r-1},
\]
the coefficient of $q^j$ in the point count of Theorem \ref{thm:finite-field-point-count} is the sum of $\binom{n+1}{r}$ over $r\ge j+1$, which is $c_{n,j}$.

Substituting $q^m$ for $q$ gives
\[
\#\Pum^n(G(\F_{q^m}))
=
\sum_{j=0}^{n}c_{n,j}q^{mj}.
\]
Therefore
\[
\begin{aligned}
\log Z_{\mathrm{sp}}
&=
\sum_{m\ge1}\sum_{j=0}^{n}c_{n,j}q^{mj}\frac{T^m}{m}\\
&=
\sum_{j=0}^{n}c_{n,j}
\sum_{m\ge1}\frac{(q^jT)^m}{m}\\
&=
-\sum_{j=0}^{n}c_{n,j}\log(1-q^jT),
\end{aligned}
\]
which exponentiates to the product formula.  For $n=1$,
\[
c_{1,0}=\binom21+\binom22=3,
\qquad
c_{1,1}=\binom22=1.
\]
\end{proof}

\begin{remark}
The ordinary projective-line zeta function is
\[
Z(\mathbf P^1/\F_q,T)=\frac{1}{(1-T)(1-qT)}.
\]
The split support completion contributes two additional degree-zero factors, corresponding to $0_\tau$ and $\infty_\tau$.
\end{remark}

\subsection{Invertible matrices}

\begin{lemma}[Boolean invertible matrices]\label{lem:boolean-gl}
Every invertible $m\times m$ matrix over $\B$ is a permutation matrix.  Hence
\[
\GL_m(\B)\cong S_m.
\]
\end{lemma}

\begin{proof}
The free $\B$-semimodule $\B^m$ is the Boolean lattice of subsets of $\{1,\dots,m\}$, with addition given by union.  A $\B$-linear automorphism preserves joins and zero, hence is an automorphism of this finite Boolean lattice.  The atoms are exactly the standard basis vectors.  Every lattice automorphism permutes the atoms, so the corresponding matrix is a permutation matrix.  Conversely every permutation matrix is invertible.
\end{proof}


\begin{theorem}[Lattice-independence of split linear symmetry]
\label{thm:lattice-monomial-gl}
Let $K$ be a field and let $L$ be any nontrivial bounded distributive
lattice.  Every invertible matrix over $G_L(K)$ has exactly one nonzero
supported entry in each row and column and $\tau_L$ in all remaining
positions.  Consequently
\[
\boxed{
\GL_m(G_L(K))\cong(K^\times)^m\rtimes S_m.
}
\]
\end{theorem}

\begin{proof}
Let $A B=B A=I_m$.  For each row $i$, the diagonal entry
$(AB)_{ii}=1$ has nonzero amplitude, so some product
$A_{ik}B_{ki}$ has nonzero amplitude.  Hence both factors are supported
units.  For $j\neq i$, the entry $(AB)_{ij}=\tau_L$.  A finite sum in
$G_L(K)$ equals $\tau_L$ only when every summand equals $\tau_L$; a
supported cancellation gives $e_L$, not $\tau_L$.  Therefore
$A_{ik}B_{kj}=\tau_L$, and the unit $A_{ik}$ forces
$B_{kj}=\tau_L$.  Applying the same argument to $(BA)_{k\ell}$ shows
$A_{i\ell}=\tau_L$ for $\ell\neq k$.

Thus row $i$ of $A$ has a unique non-$\tau_L$ entry, and it is a unit.
Two different rows cannot select the same column, because the corresponding
row of $B$ would then contain two non-$\tau_L$ entries.  The selected
columns therefore form a permutation.  The converse is immediate from the
monomial inverse.
\end{proof}

\begin{theorem}[Monomial general linear group]\label{thm:monomial-gl}
Let $K$ be a field and $S=G(K)$.  Then every matrix in $\GL_m(S)$ is monomial with one entry from $K^\times$ in each row and column and $\tau$ in every other position.  Consequently,
\[
\boxed{
\GL_m(G(K))\cong(K^\times)^m\rtimes S_m.}
\]
\end{theorem}

\begin{proof}
Let $U\in\GL_m(S)$ and let $V$ be its inverse.  Applying the support character entrywise gives
\[
\chi(U)\chi(V)=I_m,
\qquad
\chi(V)\chi(U)=I_m
\]
over $\B$.  By Lemma \ref{lem:boolean-gl}, $\chi(U)$ is a permutation matrix.  Therefore $U$ has exactly one supported entry in each row and column and has $\tau$ elsewhere.

Applying the ring reflection entrywise gives
\[
p(U)p(V)=I_m
\]
over $K$.  Thus $p(U)$ is an invertible monomial matrix, so each of its monomial entries is nonzero.  The corresponding supported entries of $U$ therefore lie in $K^\times$.

Conversely, a monomial matrix with entries in $K^\times$ has the evident monomial inverse.  The semidirect-product description is the standard decomposition into diagonal and permutation matrices.
\end{proof}

\begin{corollary}[Split projective linear group]\label{cor:split-pgl}
The group of projective transformations induced by invertible split matrices is
\[
\operatorname{PGL}^{\mathrm{lin}}_m(G(K))
\cong
\bigl((K^\times)^m/K^\times\bigr)\rtimes S_m,
\]
where the denominator is the diagonal scalar subgroup.  For $m=2$,
\[
\operatorname{PGL}^{\mathrm{lin}}_2(G(K))
\cong
K^\times\rtimes C_2.
\]
On the full-support affine coordinate $z$, the transformations are exactly
\[
z\longmapsto az
\qquad\text{and}\qquad
z\longmapsto\frac{a}{z},
\qquad a\in K^\times.
\]
\end{corollary}

\begin{proof}
Projectivization quotients $\GL_m(G(K))$ by the central diagonal copy of $K^\times$.  For $m=2$, a diagonal monomial matrix acts by scaling and an antidiagonal monomial matrix acts by scaling followed by inversion.
\end{proof}

\begin{corollary}[Obstruction to global translation]\label{cor:no-global-translation}
For $b\neq0$, the classical translation
\[
z\longmapsto z+b
\]
does not arise from a global split-projective linear automorphism of $\Pum^1(G(K))$.  More generally, a classical M\"obius transformation extends linearly to the mixed-support projective line only when its matrix is monomial after replacing structural zeros by $\tau$.
\end{corollary}

\begin{proof}
A matrix for translation is
\[
\begin{pmatrix}1&b\\0&1\end{pmatrix}.
\]
Its support pattern is not a permutation matrix, regardless of whether the classical zero entry is represented by $e$ or by $\tau$.  Theorem \ref{thm:monomial-gl} therefore excludes invertibility over $G(K)$.  The same argument applies to a general M\"obius matrix.
\end{proof}

\begin{remark}[Support-homogeneous affine maps]
Corollary \ref{cor:no-global-translation} characterizes the global split-linear group.  Affine maps are represented on every support fibre by the support-idempotent homogenization of \cite{Homogenization}, which replaces $(x,1)$ by $(x,h)$ and sends
\[
x\longmapsto ax+b
\]
by
\[
(x,h)\longmapsto(ax+bh,h).
\]
This is the correct fibrewise extension across the Boolean support diagram.
\end{remark}


\section{Classical invariant coordinates for split matrices}
\label{sec:split-matrix-invariants}


\subsection{Lattice-valued determinant support}

Let $L$ be a nontrivial bounded distributive lattice, let $K$ be a field,
and let $A=(a_{ij})\in M_n(G_L(K))$.  Define
\[
\det^{\mathrm{sp}}_L(A)
=\bigoplus_{\sigma\in S_n}
\widetilde{\operatorname{sgn}(\sigma)}
\prod_{i=1}^n a_{i,\sigma(i)}
\]
and
\[
\perL(\Lambda)
=\bigvee_{\sigma\in S_n}\bigwedge_{i=1}^n\lambda_{i,\sigma(i)}
\qquad(\Lambda=(\lambda_{ij})\in M_n(L)).
\]

\begin{theorem}[Lattice split determinant]
\label{thm:lattice-split-determinant}
For every $A\in M_n(G_L(K))$,
\[
\boxed{
p_L(\det^{\mathrm{sp}}_L A)=\det(p_LA),
\qquad
\chi_L(\det^{\mathrm{sp}}_L A)=\perL(\chi_LA).
}
\]
Hence
\[
\det^{\mathrm{sp}}_L A=
\begin{cases}
(\det(p_LA),1_L),&\det(p_LA)\neq0,\\
z_{\perL(\chi_LA)},&\det(p_LA)=0.
\end{cases}
\]
The support value therefore records the greatest lattice level at which a
perfect matching survives, even when the signed amplitudes cancel.
\end{theorem}

\begin{proof}
The amplitude projection preserves the supported signs, sums, and products,
so it sends the Leibniz formula to the ordinary determinant.  The support
of a permutation term is the meet of the selected entry supports, and the
support of the total sum is the join of the term supports.  This is exactly
the lattice permanent.  A nonzero amplitude necessarily lies at support
$1_L$; when the amplitude vanishes, the resulting element is the unique
zero-fibre element with the computed support.
\end{proof}

\begin{corollary}[Lattice-valued characteristic support]
\label{cor:lattice-characteristic-support}
Define the $k$th split characteristic coefficient by summing
$\det^{\mathrm{sp}}_L(A_I)$ over principal $k\times k$ submatrices.  Its
amplitude is the ordinary characteristic coefficient, and its support is
\[
\bigvee_{|I|=k}\ \bigvee_{\sigma\in S_I}
\bigwedge_{i\in I}\chi_L(a_{i,\sigma(i)}).
\]
Thus each coefficient records the maximal support level of a directed
cycle cover on a $k$-vertex induced subgraph.
\end{corollary}

\begin{proof}
Apply Theorem~\ref{thm:lattice-split-determinant} to every principal
submatrix and use that addition in the support fibre is join.
\end{proof}

The remaining subsections specialize to $L=\B$, where lattice support is a
single incidence bit and the formulas become the matching and closed-walk
trichotomies.

Let $K$ be a field and $S=G(K)$.  Write
\[
p:M_n(S)\longrightarrow M_n(K),
\qquad
\chi:M_n(S)\longrightarrow M_n(\B)
\]
for the entrywise amplitude and support maps.  For
$A=(a_{ij})\in M_n(S)$, its support digraph $\Gamma_A$ has vertex set
$\{1,\dots,n\}$ and an arrow $i\to j$ precisely when
$\chi(a_{ij})=1$.

\subsection{Split determinant and matching trichotomy}

\begin{definition}[Split determinant and Boolean permanent]
For $A\in M_n(S)$ define
\[
\detsp(A)
:=
\bigoplus_{\sigma\in S_n}
\widetilde{\operatorname{sgn}(\sigma)}
\prod_{i=1}^n a_{i,\sigma(i)},
\]
where $\widetilde{\pm1}$ denotes the supported scalar $\pm1\in K\subset S$.
For $B=(b_{ij})\in M_n(\B)$ define
\[
\operatorname{per}_{\B}(B)
:=
\bigvee_{\sigma\in S_n}
\bigwedge_{i=1}^n b_{i,\sigma(i)}.
\]
\end{definition}

\begin{theorem}[Split determinant theorem]
\label{thm:split-determinant}
For every $A\in M_n(S)$,
\[
p(\detsp A)=\det(pA),
\qquad
\chi(\detsp A)=\operatorname{per}_{\B}(\chi A).
\]
Consequently,
\[
\detsp A=
\begin{cases}
\tau,&\Gamma_A\text{ has no directed perfect matching},\\
e,&\Gamma_A\text{ has a directed perfect matching and }\det(pA)=0,\\
\det(pA)\in K^\times,&\det(pA)\neq0.
\end{cases}
\]
\end{theorem}

\begin{proof}
The amplitude map preserves finite sums, products, and the supported signs,
so applying $p$ to the Leibniz formula gives the ordinary determinant.
A permutation term is supported exactly when every selected entry is
supported.  Since Boolean addition and multiplication are join and meet,
applying $\chi$ gives the Boolean permanent.  If no term is supported, the
sum is the empty supported sum $\tau$.  If at least one term is supported,
the entire sum lies in the supported copy of $K$ and equals $e$ exactly when
its ordinary amplitude vanishes.
\end{proof}

\begin{example}[The two-by-two determinant]
For
\[
A=\begin{pmatrix}a&b\\c&d\end{pmatrix}
\]
one has
\[
\detsp(A)=ad-bc,
\qquad
\chi(\detsp A)
=(\chi a\wedge\chi d)\vee(\chi b\wedge\chi c).
\]
Thus $\detsp(A)=e$ records an actual diagonal or off-diagonal matching whose
two signed amplitudes cancel; $\detsp(A)=\tau$ records absence of both
matchings.
\end{example}

\subsection{Characteristic coefficients and cycle covers}

For $I\subseteq\{1,\dots,n\}$ let $A_I$ denote the principal submatrix on
$I$.  Define
\[
c_k^{\mathrm{sp}}(A)
:=
\bigoplus_{\substack{I\subseteq\{1,\dots,n\}\\|I|=k}}
\detsp(A_I),
\qquad 1\le k\le n,
\]
and
\[
\charsp(A;t)
:=t^n\oplus
\bigoplus_{k=1}^n(-1)^k c_k^{\mathrm{sp}}(A)t^{n-k}
\in S[t].
\]
Let $c_k(B)\in\B$ be the Boolean statement that some induced $k$-vertex
subgraph of the digraph of $B$ admits a directed cycle cover.

\begin{theorem}[Support-refined characteristic invariants]
\label{thm:split-characteristic}
For every $A\in M_n(S)$ and every $k$,
\[
p(c_k^{\mathrm{sp}}(A))=c_k(pA),
\qquad
\chi(c_k^{\mathrm{sp}}(A))=c_k(\chi A),
\]
where the first $c_k$ is the ordinary $k$th characteristic coefficient.
Equivalently,
\[
p(\charsp(A;t))=\det(tI-pA).
\]
Every $c_k^{\mathrm{sp}}$ is invariant under conjugation by
$\GL_n(S)$.
\end{theorem}

\begin{proof}
The amplitude identity follows by summing the ordinary determinants of all
principal $k\times k$ submatrices.  The support identity follows from
Theorem~\ref{thm:split-determinant}: a principal determinant is supported
exactly when the corresponding induced digraph has a cycle cover.

By Theorem~\ref{thm:monomial-gl}, an element of $\GL_n(S)$ is a product
$DP$ of an invertible diagonal matrix and a permutation matrix.  Diagonal
conjugation multiplies row $i$ by $d_i$ and column $i$ by $d_i^{-1}$, so
all factors cancel in every principal determinant.  Permutation conjugation
permutes the principal submatrices.  Their total sum is therefore invariant.
\end{proof}

Over an algebraically closed field of characteristic zero, the classical
one-matrix invariant theorem gives
\[
K[M_n]^{\GL_n}=K[c_1,\dots,c_n].
\]
Theorem~\ref{thm:split-characteristic} therefore supplies a canonical
support lift of every fundamental polynomial conjugation invariant.

\subsection{Trace words and closed labelled walks}

Let $A_1,\dots,A_r\in M_n(S)$, and let
$w=x_{i_1}\cdots x_{i_\ell}$ be a word in $r$ noncommuting letters.  Define
\[
T_w^{\mathrm{sp}}(A_1,\dots,A_r)
:=\tr(A_{i_1}\cdots A_{i_\ell})\in S,
\]
where the trace is the split sum of the diagonal entries.

\begin{theorem}[Trace-word support theorem]
\label{thm:split-trace-word}
For every word $w=x_{i_1}\cdots x_{i_\ell}$,
\[
p(T_w^{\mathrm{sp}})=
\tr\bigl(pA_{i_1}\cdots pA_{i_\ell}\bigr).
\]
Moreover, $\chi(T_w^{\mathrm{sp}})=1$ if and only if there exist vertices
\[
v_0,v_1,\dots,v_\ell=v_0
\]
such that $v_{j-1}\to v_j$ is supported in $A_{i_j}$ for every $j$.
The quantities $T_w^{\mathrm{sp}}$ are invariant under simultaneous
conjugation by $\GL_n(S)$.
\end{theorem}

\begin{proof}
Expanding the diagonal entries of the matrix product expresses the trace as
the sum over all closed walks with the prescribed sequence of matrix
colours.  Applying $p$ gives the usual trace expansion.  Applying $\chi$
gives the Boolean disjunction of the support conjunctions along those
walks.  For simultaneous monomial conjugation, the product transforms as
$U(A_{i_1}\cdots A_{i_\ell})U^{-1}$.  Diagonal conjugation fixes each
diagonal entry and permutation conjugation permutes them, so the split
trace is unchanged.
\end{proof}

\subsection{Cycle generators for the monomial quotient}

Let $V=M_n(K)^r$ and let $T_n=(\Gm)^n$ act by simultaneous diagonal
conjugation.  Write $x_{ij}^{(a)}$ for the coordinate in row $i$, column
$j$, and colour $a$.  A coloured directed cycle monomial is
\[
x_{i_0i_1}^{(a_1)}x_{i_1i_2}^{(a_2)}\cdots
x_{i_{m-1}i_0}^{(a_m)}.
\]

\begin{theorem}[Cycle invariant theorem]
\label{thm:cycle-invariant-ring}
The invariant ring $K[V]^{T_n}$ is generated by coloured directed-cycle
monomials.  If $N_n=T_n\rtimes S_n$ is the monomial normalizer, then
\[
K[V]^{N_n}=(K[V]^{T_n})^{S_n}
\]
is generated by the $S_n$-orbit sums of products of coloured cycle
monomials.
\end{theorem}

\begin{proof}
Under diagonal conjugation, a monomial
$M=\prod_{a,i,j}(x_{ij}^{(a)})^{m_{aij}}$ is multiplied by
\[
\prod_i t_i^{\sum_{a,j}m_{aij}-\sum_{a,j}m_{aji}}.
\]
It has weight zero exactly when the coloured directed multigraph encoded by
its exponents has equal indegree and outdegree at every vertex.  Every finite
balanced directed multigraph is a disjoint union of directed cycles: start
with any edge, follow outgoing edges until a vertex repeats, remove the
resulting cycle, and iterate.  Hence every invariant monomial factors into
cycle monomials.  Since a split torus acts diagonally on the monomial basis,
its invariant ring is spanned by its weight-zero monomials, proving the first
statement.

The normalizer is the semidirect product $T_n\rtimes S_n$, so its invariants
are the $S_n$-fixed part of the torus invariant ring.  An $S_n$-invariant
polynomial has constant coefficient on each permutation orbit of monomials;
it is therefore a linear combination of orbit sums of products of cycle
monomials.
\end{proof}

In characteristic zero, the Procesi theorem states that
$K[M_n(K)^r]^{\GL_n}$ is generated by the ordinary trace words.  The
inclusion $N_n\hookrightarrow\GL_n$ consequently induces the quotient
morphism
\[
M_n(K)^r/\!/N_n\longrightarrow M_n(K)^r/\!/\GL_n,
\]
and the pullback of every Procesi generator is the sum of the based coloured
cycle monomials with the corresponding word.  Theorem
\ref{thm:split-trace-word} refines this pullback by retaining its Boolean
closed-walk support.

\begin{corollary}[Monomial character morphism]
\label{cor:monomial-character-morphism}
For the free group $F_r$, inclusion of the monomial normalizer induces
\[
\Hom(F_r,N_n)/\!/N_n
\longrightarrow
\Hom(F_r,\GL_n)/\!/\GL_n.
\]
On coordinate rings, every trace-character function pulls back to an
explicit orbit sum of coloured cycles.  A split $F_r$-representation carries
in addition the Boolean statement that at least one such cycle is supported.
\end{corollary}


\subsection{Weighted-cycle normal form and Euler factors}

Let $M\in\GL_n(G_L(K))$.  By
Theorem~\ref{thm:lattice-monomial-gl}, there are a permutation
$\sigma\in S_n$ and weights $d_i\in K^\times$ such that
\[
M e_i=d_i e_{\sigma(i)}.
\]
For a cycle $C=(i_1\,i_2\,\dots\,i_r)$ of $\sigma$, define its cycle
weight
\[
w_C(M)=d_{i_1}d_{i_2}\cdots d_{i_r}.
\]

\begin{theorem}[Normalizer-conjugacy classification]
\label{thm:weighted-cycle-classification}
Two monomial matrices $M,M'\in\GL_n(G_L(K))$ are conjugate by
$\GL_n(G_L(K))$ if and only if their permutation cycles can be bijected so
that corresponding cycles have the same length and the same cycle weight.
Equivalently, the multiset
\[
\mathcal W(M)=\{(|C|,w_C(M)):C\text{ a cycle of }\sigma\}
\]
is a complete conjugacy invariant.
\end{theorem}

\begin{proof}
Diagonal conjugation by $D=\operatorname{diag}(c_1,\dots,c_n)$ changes the
weight on $i\to\sigma(i)$ to
$d_i c_{\sigma(i)}/c_i$, so the product around a cycle is unchanged.
Permutation conjugation merely relabels cycles.  Hence $\mathcal W(M)$ is
invariant.

Conversely, first use a permutation conjugation to identify corresponding
cycles.  On one cycle, choose $c_{i_1}=1$ and define the remaining
$c_{i_j}$ recursively so that diagonal conjugation changes the first
$r-1$ weights of $M$ to the corresponding weights of $M'$.  Equality of
cycle products forces equality of the final weight.  Carrying this out on
every cycle gives a diagonal conjugation from $M$ to $M'$.
\end{proof}

\begin{corollary}[Characteristic polynomial and traces]
\label{cor:weighted-cycle-characteristic}
For every monomial $M$,
\[
\boxed{
\det(tI-M)=\prod_C\bigl(t^{|C|}-w_C(M)\bigr).
}
\]
Moreover, for $k\ge1$,
\[
\boxed{
\Tr(M^k)=
\sum_{C:\,|C|\mid k}|C|\,w_C(M)^{k/|C|}.
}
\]
Consequently,
\[
\det(I-uM)^{-1}
=\prod_C\bigl(1-w_C(M)u^{|C|}\bigr)^{-1}.
\]
\end{corollary}

\begin{proof}
After ordering the basis by permutation cycles, $M$ is block diagonal.  A
cycle block of length $r$ has characteristic polynomial $t^r-w_C$.  Its
$k$th power has zero diagonal unless $r\mid k$; when $k=rs$, every one of
its $r$ diagonal entries equals $w_C^s$.  Multiplication over blocks gives
the formulas.
\end{proof}

\begin{theorem}[Periodic-orbit factorization for monomial local systems]
\label{thm:monomial-periodic-euler-factor}
Let $X$ be connected, let $\pi:Y\to X$ be a finite $n$-sheeted cover, let
$\mathcal L$ be a rank-one local system on $Y$, and put
$\mathcal V=\pi_*\mathcal L$.  For a closed loop $\gamma$ in $X$, let the
monodromy permutation on the fibre of $\pi$ decompose into cycles $C$.
If $r_C=|C|$, the lift of $\gamma^{r_C}$ through any sheet in $C$ is a
closed loop $\widetilde\gamma_C$ in $Y$; put
\[
w_C=\Hol_{\mathcal L}(\widetilde\gamma_C).
\]
Then
\[
\boxed{
\det(I-u\,\Hol_{\mathcal V}(\gamma))
=\prod_C(1-w_Cu^{r_C}).
}
\]
The right-hand side is independent of the chosen sheet in each cycle.
\end{theorem}

\begin{proof}
Choose a basis of the fibre of $\mathcal V$ indexed by the sheets over the
basepoint.  The holonomy of $\mathcal V$ is monomial: it permutes the sheets
and multiplies each basis vector by the rank-one transport in $\mathcal L$.
The product of these transport factors around a permutation cycle is
exactly the holonomy of $\mathcal L$ along the closed lifted loop
$\widetilde\gamma_C$.  A different starting sheet cyclically permutes the
factors and leaves their product unchanged.  Apply
Corollary~\ref{cor:weighted-cycle-characteristic}.
\end{proof}

\begin{corollary}[Frobenius-cycle factorization on the scaling site]
\label{cor:scaling-frobenius-cycle-factor}
For a finite quotient of the Connes--Consani abelian cover of the
scaling-site sector and a periodic orbit $C_p$, the lifted orbit is the
mapping torus of the Frobenius
monodromy at $p$.  Any rank-one local system on the cover therefore gives a
local determinant whose factors are indexed by the Frobenius cycles:
\[
\det(I-uM_p)=\prod_C(1-w_Cu^{|C|}).
\]
Here $|C|$ is the orbit length of the Frobenius action on the covering
fibre and $w_C$ is the rank-one holonomy along the corresponding closed
lift.
\end{corollary}

\begin{proof}
The mapping-torus description of the lifted periodic orbit is the main
monodromy theorem of Connes--Consani \cite{CCKnotsPrimes}.  Apply
Theorem~\ref{thm:monomial-periodic-euler-factor} to that covering
monodromy.
\end{proof}

\subsection{Compatibility with denominator matrix embeddings}

Let $n\mid m$, put $q=m/n$, and identify
$M_m(S)\cong M_n(S)\otimes M_q(S)$.  The denominator embedding is
\[
\iota_{n,m}(A)=A\otimes I_q.
\]

\begin{proposition}[Coherent classical invariants along the UHF tower]
\label{prop:uhf-invariant-compatibility}
For every $A\in M_n(S)$ and every positive integer $d$,
\begin{align*}
\charsp(\iota_{n,m}(A);t)&=\charsp(A;t)^q,\\
\detsp(\iota_{n,m}(A))&=\detsp(A)^q,\\
\tr\bigl(\iota_{n,m}(A)^d\bigr)&=q\,\tr(A^d).
\end{align*}
Over $K=\C$, normalized traces are preserved:
\[
\frac1m\tr\bigl(\iota_{n,m}(A)^d\bigr)
=
\frac1n\tr(A^d).
\]
\end{proposition}

\begin{proof}
A permutation similarity identifies $A\otimes I_q$ with a block diagonal
matrix consisting of $q$ copies of $A$.  The split Leibniz formula factors
over block diagonal matrices, giving the characteristic-polynomial and
determinant identities.  The power remains block diagonal with $q$ copies
of $A^d$, giving the trace identity and its normalized form.
\end{proof}

\part{Infinitesimal and projective analysis}

\section{Residue monads and Boolean-cubical infinitesimal geometry}

The support residue $\B$ is not itself Leibniz's infinitesimal monad.  A monad requires a residue map with a nontrivial kernel of infinitesimal amplitudes.  Split zero supplies the support geometry in which such kernels are organized.

\begin{definition}[Split residue datum]\label{def:split-residue-datum}
A \emph{split residue datum} is a surjective ring homomorphism
\[
q:A\twoheadrightarrow k
\]
with kernel
\[
\mathfrak m:=\ker q.
\]
It induces a semiring homomorphism
\[
G(q):G(A)\longrightarrow G(k)
\]
by
\[
G(q)(\tau)=\tau,
\qquad
G(q)(a)=q(a)
\quad(a\in A).
\]
\end{definition}

\begin{definition}[Supported monad]
For a supported element $a\in k\subset G(k)$ and a chosen lift $\widetilde a\in A$, define
\[
\mu_q(a):=G(q)^{-1}(a).
\]
For the absent state define
\[
\mu_q(\tau):=G(q)^{-1}(\tau).
\]
\end{definition}

\begin{theorem}[Monad fibre formula]\label{thm:monad-fibre}
Let $q:A\twoheadrightarrow k$ be a split residue datum.  Then
\[
\mu_q(a)=\widetilde a+\mathfrak m
\qquad(a\in k),
\]
and
\[
\mu_q(\tau)=\{\tau\}.
\]
In particular,
\[
\mu_q(e_k)=\mathfrak m,
\]
where $e_k=0_k$ is the supported zero in $G(k)$.
\end{theorem}

\begin{proof}
A supported element $x\in A\subset G(A)$ maps to $a$ exactly when $q(x)=a$, which is equivalent to $x\in\widetilde a+\mathfrak m$.  No supported element maps to $\tau$, because $G(q)$ preserves support.  Hence the only preimage of $\tau$ is $\tau$ itself.
\end{proof}

\begin{corollary}[Centre and halo]\label{cor:centre-and-halo}
The supported zero $e_k$ is the centre of the zero monad, the ideal $\mathfrak m$ is its infinitesimal halo, and $\tau$ is the absent state outside that halo.
\end{corollary}

\begin{proof}
The fibre over $e_k$ is $\mathfrak m$ by Theorem \ref{thm:monad-fibre}.  Its distinguished zero element is $0_A=e_A$, which maps to $e_k$.  The absent state has the singleton fibre $\{\tau\}$.
\end{proof}

The higher-dimensional construction is cubical.  Let $I$ be finite and let
\[
a=(a_i)_{i\in I}\in G(k)^I.
\]
Write
\[
J:=\supp(a)=\{i:a_i\in k\}.
\]
For $i\in J$, choose a lift $\widetilde a_i\in A$.

\begin{definition}[Support star and cubical monad]\label{def:cubical-monad}
The \emph{support star} of $a$ is
\[
\operatorname{Star}(a)
:=
\left\{
 b\in G(k)^I:
 b_i=a_i\ (i\in J),
\quad
b_i\in\{\tau,e_k\}\ (i\notin J)
\right\}.
\]
The \emph{upward cubical monad} of $a$ is
\[
\mu_q^\square(a)
:=
(G(q)^I)^{-1}(\operatorname{Star}(a)).
\]
\end{definition}

\begin{theorem}[Boolean-cubical monad decomposition]\label{thm:cubical-monad}
There is a canonical disjoint decomposition
\[
\boxed{
\mu_q^\square(a)
=
\bigsqcup_{L\subseteq I\setminus J}
\left(
\prod_{i\in J}(\widetilde a_i+\mathfrak m)
\right)
\times
\mathfrak m^L
\times
\{\tau\}^{I\setminus(J\cup L)}.}
\]
The subset $L$ records exactly the formerly absent coordinates that have been activated with residue zero.
\end{theorem}

\begin{proof}
Every $b\in\operatorname{Star}(a)$ is determined by the subset
\[
L:=\{i\in I\setminus J:b_i=e_k\}.
\]
For $i\in J$, Theorem \ref{thm:monad-fibre} gives the fibre $\widetilde a_i+\mathfrak m$.  For $i\in L$, the fibre of $e_k$ is $\mathfrak m$.  For the remaining coordinates, the fibre of $\tau$ is $\{\tau\}$.  Different subsets $L$ produce disjoint support patterns, so the union is disjoint.
\end{proof}

\begin{example}[The local path from absence to infinitesimal amplitude]
Take one coordinate and the absent base point $a=\tau$.  Then
\[
\mu_q^\square(\tau)
=
\{\tau\}\sqcup\mathfrak m.
\]
Inside the second component lie
\[
e_A=0_A
\]
and all nonzero infinitesimals $h\in\mathfrak m\setminus\{0\}$.  Thus the exact local sequence is
\[
\tau
\longrightarrow
 e_A
\longrightarrow
 h\in\mathfrak m\setminus\{0\},
\]
meaning absence, supported exact zero, and supported infinitesimal amplitude.
\end{example}

\subsection{First-order cubical tangent data}

Suppose now that $k$ is a field.  The ordinary first-order tangent space of the residue datum is the $k$-vector space
\[
T_q:=\mathfrak m/\mathfrak m^2.
\]
The cubical monad carries one copy of this vector space for every active coordinate.

\begin{theorem}[Boolean cube of first-order tangent fibres]\label{thm:first-order-cubical}
Let $q:A\twoheadrightarrow k$ be a surjection onto a field, let $a\in G(k)^I$ have support $J$, and fix a stratum indexed by $L\subseteq I\setminus J$ in Theorem \ref{thm:cubical-monad}.  Modulo coordinatewise congruence by $\mathfrak m^2$, that stratum is an affine torsor under
\[
T_q^{J\cup L}.
\]
As $L$ varies, the first-order split neighbourhood is a Boolean cube of affine tangent spaces indexed by $\mathcal P(I\setminus J)$.
\end{theorem}

\begin{proof}
For each $i\in J$, the coset $\widetilde a_i+\mathfrak m$ modulo $\mathfrak m^2$ is an affine torsor under $\mathfrak m/\mathfrak m^2$.  For each newly activated coordinate $i\in L$, the fibre $\mathfrak m$ modulo $\mathfrak m^2$ is the vector space $T_q$ itself.  Coordinates outside $J\cup L$ are fixed at $\tau$.  Taking products gives an affine torsor under $T_q^{J\cup L}$.  The indexing subsets $L$ form the Boolean cube of possible support activations.
\end{proof}

\begin{remark}
Theorem \ref{thm:first-order-cubical} separates two independent infinitesimal directions:
\[
\text{Boolean activation between support faces}
\]
and
\[
\text{linear tangent displacement inside an active face}.
\]
The element $e$ controls the first, while $\mathfrak m/\mathfrak m^2$ controls the second.
\end{remark}

\subsection{Formal completion}

Let $A$ be a ring and $I\triangleleft A$ an ideal.  The quotient maps $A/I^{r+1}\to A/I^r$ induce an inverse system of split-zero semirings.

\begin{theorem}[Split zero commutes with adic completion]\label{thm:completion-commutes}
There is a canonical semiring isomorphism
\[
\varprojlim_{r\ge1}G(A/I^r)
\xrightarrow{\sim}
G\left(\varprojlim_{r\ge1}A/I^r\right).
\]
In particular, if $A$ is $I$-adically complete, then
\[
\varprojlim_{r\ge1}G(A/I^r)
\cong
G(A).
\]
\end{theorem}

\begin{proof}
An element of the inverse limit on the left is a coherent sequence $(x_r)_{r\ge1}$.  Every transition map sends $\tau$ to $\tau$ and sends supported elements to supported elements.  Hence either all $x_r$ equal $\tau$, or every $x_r$ is supported.  In the first case the sequence corresponds to the external zero of the right-hand side.  In the second case the sequence is exactly a coherent sequence in the ring quotients $A/I^r$, hence an element of $\varprojlim A/I^r$, regarded as supported.  This gives a bijection, and the operations are preserved coordinatewise.
\end{proof}

\begin{definition}[Split formal monad]
The split formal neighbourhood associated with $(A,I)$ is the pro-semiring
\[
\operatorname{Spf}_{\mathrm{sp}}(A,I)
:=
\{G(A/I^r)\}_{r\ge1}.
\]
Its supported sector is the ordinary formal neighbourhood; its support masks form the Boolean incidence skeleton.
\end{definition}

\section{Hyperreal and surcomplex realizations}

\subsection{Finite hyperreals}

Let ${}^{\ast}\R$ be a nonstandard elementary extension of $\R$.  Let
\[
\cO_{\mathrm{hyp}}
:=
\{x\in{}^{\ast}\R:|x|<n\text{ for some }n\in\N\}
\]
be the ring of finite hyperreals and let
\[
\mathfrak m_{\mathrm{hyp}}
:=
\{x\in{}^{\ast}\R:|x|<1/n\text{ for every }n\ge1\}
\]
be its ideal of infinitesimals.  The standard-part map is a surjective ring homomorphism
\[
\st:\cO_{\mathrm{hyp}}\twoheadrightarrow\R
\]
with kernel $\mathfrak m_{\mathrm{hyp}}$; see \cite{Robinson,Goldblatt}.

\begin{corollary}[Split hyperreal monads]\label{cor:hyperreal-monads}
For every $a\in\R$,
\[
\mu_{\mathrm{hyp}}(a)
=
a+\mathfrak m_{\mathrm{hyp}}.
\]
For a mixed-support point $a\in G(\R)^I$, the upward monad is the Boolean cube of ordinary hyperreal monads described by Theorem \ref{thm:cubical-monad}.
\end{corollary}

\begin{proof}
Apply Theorems \ref{thm:monad-fibre} and \ref{thm:cubical-monad} to the standard-part map.
\end{proof}

\subsection{Surreal and surcomplex normal forms}

We work in von Neumann--Bernays--G\"odel set theory so that the proper class $\No$ may be treated as a class-sized ordered field.  Every support occurring in a Conway normal form is a set.  We use Conway's normal-form theorem and the standard monomial map
\[
x\longmapsto\omega^x
\]
from \cite{Conway}.  The uploaded expository notes \cite{Simons} give a concise account of birthdays, ordinal embeddings, and the first infinitesimal $\omega^{-1}$.

Set
\[
K:=\No[i]
=
\{x+iy:x,y\in\No,\ i^2=-1\}.
\]
Conway proved that $K$ is algebraically closed \cite{Conway}.

\begin{theorem}[Surcomplex Conway normal form]\label{thm:surcomplex-normal-form}
Every nonzero $z\in K$ has a unique expansion
\[
z=
\sum_{\rho<\lambda}
 c_\rho\omega^{\gamma_\rho},
\]
where
\[
c_\rho\in\C^\times,
\qquad
\gamma_0>\gamma_1>\cdots,
\]
and the support $\{\gamma_\rho:\rho<\lambda\}$ is reverse well ordered.  Equivalently, $K$ is the class-sized Hahn field with coefficient field $\C$ and value group $\No$, with set-sized reverse-well-ordered support.
\end{theorem}

\begin{proof}
Write $z=x+iy$ with $x,y\in\No$.  By Conway's normal-form theorem,
\[
x=\sum_{\gamma\in A}a_\gamma\omega^\gamma,
\qquad
 y=\sum_{\gamma\in B}b_\gamma\omega^\gamma,
\]
where $A$ and $B$ are reverse well ordered and the coefficients are real.  The finite union $A\cup B$ is reverse well ordered.  Extending missing coefficients by zero gives
\[
z=
\sum_{\gamma\in A\cup B}
(a_\gamma+ib_\gamma)\omega^\gamma.
\]
Delete zero coefficients.  Uniqueness follows by comparing real and imaginary parts and using uniqueness of the real Conway normal form.
\end{proof}

For nonzero $z$, let $\ell(z)$ be its leading exponent and define the natural valuation
\[
v(z):=-\ell(z),
\qquad
v(0):=\infty.
\]

\begin{proposition}[Valuation laws]\label{prop:surcomplex-valuation}
For $z,w\in K$,
\[
v(zw)=v(z)+v(w),
\]
and
\[
v(z+w)\ge\min\{v(z),v(w)\},
\]
with equality unless the leading terms cancel.  If $z=x+iy$, then
\[
v(z)=\min\{v(x),v(y)\}.
\]
\end{proposition}

\begin{proof}
If
\[
z=c\omega^\alpha+\text{lower terms},
\qquad
w=d\omega^\beta+\text{lower terms},
\]
then
\[
zw=cd\,\omega^{\alpha+\beta}+\text{lower terms},
\]
and $cd\neq0$.  Thus $\ell(zw)=\alpha+\beta$, proving multiplicativity of $v$.  The addition inequality is the standard leading-term argument.

For $z=x+iy$, let the leading real exponents of $x$ and $y$ be $\alpha$ and $\beta$.  If $\alpha\neq\beta$, the larger exponent survives in $z$.  If $\alpha=\beta$, the leading complex coefficient is $a+ib$ with $a,b\in\R$ not both zero, so there is no cancellation between real and imaginary components.  Hence $\ell(z)=\max\{\alpha,\beta\}$, which is equivalent to the displayed valuation formula.
\end{proof}

Define the finite valuation ring and its maximal ideal by
\[
\cO_K:=\{z\in K:v(z)\ge0\},
\qquad
\mathfrak m_K:=\{z\in K:v(z)>0\}.
\]

\begin{theorem}[Surcomplex standard part]\label{thm:surcomplex-standard-part}
The coefficient of $\omega^0$ defines a surjective ring homomorphism
\[
\st:\cO_K\twoheadrightarrow\C
\]
with kernel $\mathfrak m_K$.  Hence
\[
\cO_K/\mathfrak m_K\cong\C.
\]
\end{theorem}

\begin{proof}
An element of $\cO_K$ has no positive exponent in its normal form.  Define $\st(z)$ to be the coefficient of $\omega^0$, taking that coefficient to be zero if the exponent $0$ is absent.  Addition preserves constant coefficients.  In a product of two finite series, an exponent sum can equal zero only when both exponents are zero, because all exponents are nonpositive.  Therefore the constant coefficient of a product is the product of the constant coefficients, so $\st$ is a ring homomorphism.  It is surjective because every $c\in\C$ is the constant series $c\omega^0$.  Its kernel consists exactly of series whose leading exponent is negative, equivalently those of positive valuation.
\end{proof}

\begin{corollary}[The surcomplex Leibniz monad]\label{cor:surcomplex-monad}
For every $c\in\C$,
\[
\mu_{\mathrm{sur}}(c)
=
c+\mathfrak m_K.
\]
In particular, the zero monad is
\[
\mu_{\mathrm{sur}}(0)=\mathfrak m_K.
\]
After split-zero globalization, the upward mixed-support monad is the Boolean cube of these surcomplex fibres.
\end{corollary}

\begin{proof}
Apply Theorems \ref{thm:monad-fibre} and \ref{thm:cubical-monad} to the standard-part map of Theorem \ref{thm:surcomplex-standard-part}.
\end{proof}

\subsection{Ordinal scales and reciprocal infinitesimals}

Every ordinal embeds canonically in $\No$.  Inside the ordered field $\No$, its arithmetic is the natural or Hessenberg arithmetic induced by Conway normal forms, not the noncommutative ordinal addition and multiplication.

\begin{proposition}[Reciprocals of infinite ordinals]\label{prop:ordinal-reciprocals}
If $\alpha$ is an infinite ordinal, regarded as a positive surreal number, then
\[
\alpha^{-1}
\]
is a positive surreal infinitesimal.  The map
\[
\alpha\longmapsto\alpha^{-1}
\]
is order reversing on positive ordinals.  Since $\omega$ is the least infinite ordinal, $\omega^{-1}$ is the largest reciprocal of an infinite ordinal, although it is not the largest positive surreal infinitesimal.
\end{proposition}

\begin{proof}
For every positive integer $n$,
\[
\alpha>n.
\]
Inversion in an ordered field reverses inequalities among positive elements, so
\[
0<\alpha^{-1}<n^{-1}.
\]
This is the definition of a positive infinitesimal.  Order reversal is the same argument.  The final assertion follows from minimality of $\omega$ among infinite ordinals.  There is no largest positive infinitesimal because $2\epsilon$ is infinitesimal and larger than $\epsilon$ whenever $\epsilon>0$ is infinitesimal.
\end{proof}

\begin{theorem}[Reciprocal ordinals are coinitial]\label{thm:reciprocal-coinitial}
For every positive surreal infinitesimal $\epsilon$, there exists an infinite ordinal $\alpha$ such that
\[
0<\alpha^{-1}<\epsilon.
\]
\end{theorem}

\begin{proof}
The reciprocal $x:=\epsilon^{-1}$ is positive infinite.  The ordinals are cofinal in $\No$; equivalently, every surreal number is bounded above by an ordinal \cite[Chap.~5]{Conway}.  Choose an ordinal $\alpha>x$.  Then $\alpha$ is infinite and inversion gives
\[
0<\alpha^{-1}<x^{-1}=\epsilon.
\]
\end{proof}

\begin{lemma}[Geometric inversion in the Hahn field]\label{lem:hahn-geometric}
If $u\in K$ satisfies $v(u)>0$, then the family $\{(-u)^n:n\ge0\}$ is Hahn-summable and
\[
(1+u)^{-1}=\sum_{n=0}^{\infty}(-u)^n.
\]
\end{lemma}

\begin{proof}
Since $v(u^n)=nv(u)$ tends strictly upward, for every valuation bound only finitely many terms have valuation below that bound.  This is precisely the Hahn summability condition.  Multiplying partial sums gives
\[
(1+u)\sum_{n=0}^{N}(-u)^n
=1+(-1)^Nu^{N+1}.
\]
The remainder tends to arbitrarily high valuation, so the Hahn sum is an inverse of $1+u$.
\end{proof}

\begin{theorem}[Reciprocal from Cantor normal form]\label{thm:ordinal-inverse-expansion}
Let an infinite ordinal have Cantor normal form
\[
\alpha
=
n_0\omega^{\beta_0}
+n_1\omega^{\beta_1}
+\cdots
+n_m\omega^{\beta_m},
\]
where
\[
\beta_0>\beta_1>\cdots>\beta_m,
\qquad
n_j\in\N_{>0}.
\]
Set
\[
u:=
\sum_{j=1}^{m}
\frac{n_j}{n_0}\omega^{\beta_j-\beta_0}.
\]
Then $u$ is infinitesimal and
\[
\boxed{
\alpha^{-1}
=
\frac1{n_0}\omega^{-\beta_0}
\sum_{k=0}^{\infty}(-u)^k.}
\]
Its leading term is
\[
\frac1{n_0}\omega^{-\beta_0},
\]
so
\[
v(\alpha^{-1})=\beta_0.
\]
\end{theorem}

\begin{proof}
Inside $\No$, Cantor normal form is Conway normal form for ordinals.  Factor
\[
\alpha
=
n_0\omega^{\beta_0}(1+u).
\]
Each exponent $\beta_j-\beta_0$ is negative, so $v(u)>0$.  Lemma \ref{lem:hahn-geometric} gives the inverse of $1+u$, and multiplication by $(n_0\omega^{\beta_0})^{-1}$ gives the formula.  Since the geometric factor has leading term $1$, the leading term of $\alpha^{-1}$ is the displayed monomial.
\end{proof}

\begin{example}
For $\alpha=\omega+1$,
\[
\frac1{\omega+1}
=
\omega^{-1}
-\omega^{-2}
+\omega^{-3}
-\omega^{-4}
+\cdots.
\]
This is an exact surreal equality.
\end{example}

\subsection{Complex phase at every infinitesimal order}

For $\lambda\in\No$, define
\[
F_{\ge\lambda}:=\{z\in K:v(z)\ge\lambda\},
\qquad
F_{>\lambda}:=\{z\in K:v(z)>\lambda\}.
\]

\begin{theorem}[Complex graded layers]\label{thm:complex-graded-layers}
For every $\lambda\in\No$, coefficient extraction at the monomial $\omega^{-\lambda}$ induces a canonical $\C$-linear isomorphism
\[
F_{\ge\lambda}/F_{>\lambda}
\xrightarrow{\sim}
\C.
\]
Multiplication of initial terms gives bilinear maps
\[
\gr_\lambda(K)\times\gr_\mu(K)
\longrightarrow
\gr_{\lambda+\mu}(K)
\]
corresponding to multiplication in $\C$.
\end{theorem}

\begin{proof}
For $z\in F_{\ge\lambda}$, let $c_{-\lambda}(z)$ be the coefficient of $\omega^{-\lambda}$, taking it to be zero when that monomial is absent.  This is a $\C$-linear map to $\C$.  Its kernel consists exactly of elements with valuation strictly greater than $\lambda$, namely $F_{>\lambda}$.  It is surjective because $c\omega^{-\lambda}$ maps to $c$.  The multiplication statement follows from
\[
(c\omega^{-\lambda})(d\omega^{-\mu})
=cd\,\omega^{-(\lambda+\mu)}.
\]
\end{proof}

\begin{definition}[Support--order--phase initial datum]\label{def:support-order-phase}
For $x\in G(K)$ define
\[
\operatorname{in}_{\mathrm{sp}}(x)
=
\begin{cases}
(0,\bot,0),&x=\tau,\\
(1,\infty,0),&x=e,\\
(1,v(x),c),&x\in K^\times,
\end{cases}
\]
where $c\in\C^\times$ is the leading coefficient of $x$.  Here $\bot$ denotes absence and $\infty$ denotes supported exact zero.
\end{definition}

\begin{remark}
The three coordinates have independent meanings:
\[
\text{support}\in\B,
\qquad
\text{order}\in\No\cup\{\infty\},
\qquad
\text{phase}\in\C.
\]
The split-zero residue retains the first, the valuation retains the second, and the associated graded retains the third.
\end{remark}

\begin{corollary}[Mixed surreal division]\label{cor:mixed-surreal-division}
Let $\alpha>0$ be an ordinal and let $S=G(\No[i])$.  Then
\[
S^2[(\omega^\alpha,e)^{-1}]
\cong
S\times\B,
\]
and
\[
\frac{(a,b)}{(\omega^\alpha,e)}
=
\bigl(a\omega^{-\alpha},\chi(b)\bigr).
\]
Similarly,
\[
\frac{(a,b)}{(\omega^{-\alpha},e)}
=
\bigl(a\omega^{\alpha},\chi(b)\bigr).
\]
\end{corollary}

\begin{proof}
Every nonzero surreal or surcomplex number is a unit of $K$.  Apply Theorem \ref{thm:mixed-localization}.
\end{proof}

\begin{corollary}[Strictly identifying $e$ with a surreal infinitesimal]
For every nonzero $h\in\No[i]$, including $h=\omega^{-\alpha}$,
\[
G(\No[i])/(e=h)\cong\B.
\]
On a selected coordinate face $A$, the corresponding quotient is the partial Booleanization of Theorem \ref{thm:scoped-identity}.
\end{corollary}

\begin{proof}
Apply Theorems \ref{thm:strict-identification-collapse} and \ref{thm:scoped-identity}.
\end{proof}

\subsection{Differential separation of \texorpdfstring{$e$ and $\omega^{-1}$}{e and omega inverse}}

Berarducci and Mantova construct a simplest surreal derivation
\[
\partial:\No\longrightarrow\No
\]
with constant field $\R$, strong additivity, compatibility with exponentiation, and
\[
\partial(\omega)=1;
\]
see \cite{BerarducciMantova}.  Extend it $\C$-linearly to $K=\No[i]$ by $\partial(i)=0$.

\begin{proposition}[Differential obstruction to identification]\label{prop:differential-obstruction}
One has
\[
\partial(\omega^{-1})=-\omega^{-2}\neq0,
\]
whereas every semiring derivation on $G(K)$ kills the supported zero $e$.  Hence no derivation-preserving identification can send $e$ to $\omega^{-1}$.
\end{proposition}

\begin{proof}
Differentiate $\omega\omega^{-1}=1$:
\[
0
=
\partial(1)
=
\partial(\omega)\omega^{-1}
+
\omega\partial(\omega^{-1})
=
\omega^{-1}
+
\omega\partial(\omega^{-1}).
\]
Multiplication by $\omega^{-1}$ gives
\[
\partial(\omega^{-1})=-\omega^{-2}.
\]
The statement for $e$ is Proposition \ref{prop:derivation-kills-e}.
\end{proof}

\section{Projective zero, projective infinity, and the surcomplex circle}

\subsection{Inversion and the two poles}

Let $K$ be a field and $S=G(K)$.  Consider the matrix
\[
J:=
\begin{pmatrix}
\tau&1\\
1&\tau
\end{pmatrix}.
\]

\begin{proposition}[Global split inversion]\label{prop:global-inversion}
The matrix $J$ belongs to $\GL_2(S)$ and satisfies $J^2=I$.  Its action on $\Pum^1(S)$ is
\[
[x:y]\longmapsto[y:x].
\]
In particular,
\[
0_\tau\longleftrightarrow\infty_\tau,
\qquad
0_e\longleftrightarrow\infty_e.
\]
\end{proposition}

\begin{proof}
A direct multiplication gives
\[
J^2
=
\begin{pmatrix}
\tau^2+1&\tau+\tau\\
\tau+\tau&1+\tau^2
\end{pmatrix}
=
\begin{pmatrix}
1&\tau\\
\tau&1
\end{pmatrix}
=I.
\]
The coordinate action and the stated exchanges are immediate.
\end{proof}

\begin{theorem}[Equivalence of zero and infinity monads]\label{thm:zero-infinity-monads}
Let $q:A\twoheadrightarrow k$ be a split residue datum and let $S_A=G(A)$.  In the affine chart $[x:1]$, the upward zero monad is
\[
\{[\tau:1]\}
\sqcup
\{[h:1]:h\in\mathfrak m\}.
\]
Projective inversion restricts to a bijection onto the upward infinity monad
\[
\{[1:\tau]\}
\sqcup
\{[1:h]:h\in\mathfrak m\}.
\]
It sends
\[
[\tau:1]\mapsto[1:\tau],
\qquad
[e:1]\mapsto[1:e],
\]
and, for $h\neq0$,
\[
[h:1]\mapsto[1:h]=[h^{-1}:1]
\]
whenever $A$ is contained in a field in which $h$ is invertible.
\end{theorem}

\begin{proof}
The description of the upward zero monad is Theorem \ref{thm:cubical-monad} in one coordinate, written in the affine projective chart.  Proposition \ref{prop:global-inversion} swaps the homogeneous coordinates and gives the displayed image.  If $h$ is nonzero in an ambient field, multiplication by the unit $h^{-1}$ identifies $[1:h]$ with $[h^{-1}:1]$.
\end{proof}

\begin{corollary}[Projective realization of infinity]\label{cor:no-x-infinity}
The supported projective infinity
\[
\infty_e=[1:e]
\]
and the absent-denominator infinity
\[
\infty_\tau=[1:\tau]
\]
are points of $\Pum^1(G(K))$.  Under inversion, the punctured valuative monad of $0_e$ is carried to the punctured monad of $\infty_e$, whose elements are precisely the corresponding infinite field elements.
\end{corollary}

\begin{proof}
The two points are distinct by their support masks.  The punctured-monad assertion is Theorem \ref{thm:zero-infinity-monads} applied to a valued field: nonzero infinitesimals invert to infinite elements.
\end{proof}

\subsection{The Cayley circle}

Let $F$ be a real closed field, let $F[i]$ be its algebraic closure, and define
\[
U_F:=\{z\in F[i]:z\overline z=1\}.
\]

\begin{theorem}[Algebraic Cayley circle]\label{thm:cayley-circle}
The Cayley transform
\[
C:\mathbf P^1(F)\longrightarrow U_F
\]
defined by
\[
C(x)=\frac{x-i}{x+i}
\qquad(x\in F),
\qquad
C(\infty)=1,
\]
is a bijection.  Its inverse is
\[
C^{-1}(z)=i\frac{1+z}{1-z}
\qquad(z\neq1),
\qquad
C^{-1}(1)=\infty.
\]
In particular,
\[
\mathbf P^1(\No)
\cong
\{z\in\No[i]:z\overline z=1\}.
\]
\end{theorem}

\begin{proof}
For $x\in F$,
\[
C(x)\overline{C(x)}
=
\frac{x-i}{x+i}\frac{x+i}{x-i}=1.
\]
Conversely, let $z\in U_F$ with $z\neq1$.  Since $\overline z=z^{-1}$,
\[
\overline{\left(i\frac{1+z}{1-z}\right)}
=
-i\frac{1+z^{-1}}{1-z^{-1}}
=
-i\frac{z+1}{z-1}
=
i\frac{1+z}{1-z}.
\]
Thus $C^{-1}(z)$ is fixed by conjugation and lies in $F$.  Direct substitution shows that the formulas are inverse.  The final statement uses that $\No$ is real closed.
\end{proof}

\begin{proposition}[Cayley obstruction on mixed support]\label{prop:cayley-obstruction}
The Cayley transform acts on the full-support classical stratum but is not induced by a global split-projective linear automorphism of $\Pum^1(G(F[i]))$.
\end{proposition}

\begin{proof}
A matrix for the Cayley transform is
\[
\begin{pmatrix}
1&-i\\
1&i
\end{pmatrix}.
\]
Every entry is supported, so its Boolean support matrix is
\[
\begin{pmatrix}
1&1\\
1&1
\end{pmatrix},
\]
which is not invertible over $\B$.  Theorem \ref{thm:monomial-gl} therefore excludes invertibility over the split semiring.  Over the field $F[i]$, the determinant is nonzero and the classical action is valid on the full-support stratum.
\end{proof}

\subsection{Placing \texorpdfstring{$\omega$}{omega} at a projective pole}

In the classical surcomplex projective line, the M\"obius transformation
\[
u=\frac{x}{\omega-x}
\]
is induced by
\[
[X:Y]\longmapsto[X:\omega Y-X].
\]
It sends $x=0$ to $u=0$ and $x=\omega$ to $u=\infty$.  This realizes $\omega$ as a projective pole rather than as a new arithmetic zero.

\begin{proposition}[The $\omega$-pole chart is not globally split-linear]\label{prop:omega-pole-obstruction}
The transformation $x\mapsto x/(\omega-x)$ is a classical projective automorphism of $\mathbf P^1(\No[i])$, but it does not extend as a global linear automorphism of $\Pum^1(G(\No[i]))$.
\end{proposition}

\begin{proof}
Its classical determinant is $\omega\neq0$, so it is a field-projective automorphism.  Any representing split matrix has a support pattern with at least three supported entries, hence its Boolean support matrix is not a permutation matrix.  Theorem \ref{thm:monomial-gl} excludes global split invertibility.
\end{proof}

\begin{remark}[Birthdays]
In Conway's construction, $0=\{\mid\}$ is born on day $0$, while $\omega^{-1}$ is born at a transfinite stage; see \cite{Conway,Simons}.  The external state $\tau$ is not an earlier surreal number.  One may assign it a separate support grade $\bot$ below all birthdays, but this is an enrichment of the birthday filtration, not a modification of Conway's recursive order.
\end{remark}

\section{Limits, poles, and removable indeterminacies}

The split formalism does not replace all limit theories by a single operation.  It separates a limit problem into an amplitude condition inside a monad and a support condition across the Boolean incidence skeleton.

\begin{definition}[Residue-monadic limit]\label{def:monadic-limit}
Let $q:A\twoheadrightarrow k$ be a split residue datum, let $a,L\in k$, choose a lift $\widetilde a\in A$, and let
\[
f:(\widetilde a+\mathfrak m)\setminus\{\widetilde a\}\longrightarrow A
\]
be a function.  We say that $f$ has \emph{residue-monadic limit} $L$ at $a$ if
\[
q(f(\widetilde a+h))=L
\]
for every nonzero $h\in\mathfrak m$ for which $f$ is defined.
\end{definition}

For finite hyperreals and an internal extension of a real function, this is the standard nonstandard criterion for a limit; see \cite{Robinson,Goldblatt}.  For the surcomplex valuation ring it states that all infinitesimal perturbations have the same complex standard part.

\begin{proposition}[Projective reciprocal extension]\label{prop:reciprocal-extension}
For every field $K$, the map
\[
\mathcal R:G(K)\longrightarrow\Pum^1(G(K)),
\qquad
x\longmapsto[1:x],
\]
extends ordinary reciprocal projectively.  It satisfies
\[
\mathcal R(h)=[h^{-1}:1]
\qquad(h\in K^\times),
\]
\[
\mathcal R(e)=\infty_e,
\qquad
\mathcal R(\tau)=\infty_\tau.
\]
\end{proposition}

\begin{proof}
If $h\in K^\times$, simultaneous multiplication by $h^{-1}$ gives
\[
[1:h]=[h^{-1}:1].
\]
The two zero cases are the definitions of $\infty_e$ and $\infty_\tau$.
\end{proof}

Thus a pole at supported zero and an absent denominator have distinct projective values.

\begin{theorem}[Resolution of a common-factor indeterminacy]\label{thm:common-factor-resolution}
Let $K$ be a field, let $a\in K$, and let $P,Q\in K[T]$.  Suppose
\[
P(T)=(T-a)^rP_0(T),
\qquad
Q(T)=(T-a)^rQ_0(T)
\]
for some $r\ge1$, with
\[
Q_0(a)\neq0.
\]
Then the rational projective map
\[
\Phi(x):=[P(x):Q(x)]
\]
agrees for $x\neq a$ with
\[
\widetilde\Phi(x):=[P_0(x):Q_0(x)],
\]
and $\widetilde\Phi$ extends to $x=a$ with value
\[
\widetilde\Phi(a)=[P_0(a):Q_0(a)].
\]
In $G(K)$, the unresolved pair at $x=a$ is
\[
[e:e],
\]
which has full support but is non-unimodular.  It is distinct from the absent pair
\[
[\tau:\tau].
\]
\end{theorem}

\begin{proof}
For $x\neq a$, the scalar $(x-a)^r$ is a nonzero unit of $K$, so
\[
[P(x):Q(x)]
=
[(x-a)^rP_0(x):(x-a)^rQ_0(x)]
=
[P_0(x):Q_0(x)].
\]
At $x=a$, the second coordinate $Q_0(a)$ is a unit, so $[P_0(a):Q_0(a)]$ is unimodular and defines a projective point.  Before cancellation, both $P(a)$ and $Q(a)$ are the supported ring zero $e$, giving the full-support but non-unimodular pair $(e,e)$.  The pair $(\tau,\tau)$ has empty support, so the two are distinct in $G(K)^2$.
\end{proof}

\begin{example}[The form $0/0$ at $x=1$]
For
\[
f(x)=\frac{x^2-1}{x-1},
\]
Theorem \ref{thm:common-factor-resolution} gives
\[
[x^2-1:x-1]=[x+1:1]
\]
away from $x=1$, and the resolved value is
\[
[2:1].
\]
If $h$ is a nonzero hyperreal or surcomplex infinitesimal, then
\[
f(1+h)=\frac{2h+h^2}{h}=2+h,
\]
so its standard part is $2$.  The Newtonian limit, the Leibniz infinitesimal computation, and the projective common-factor resolution therefore give the same amplitude result, while split support additionally distinguishes $(e,e)$ from $(\tau,\tau)$.
\end{example}

\begin{theorem}[Two-component split limit datum]\label{thm:two-component-limit}
For a map defined on a mixed-support punctured neighbourhood, a split limit consists of two independent conditions:
\begin{enumerate}[label=\textup{(\alph*)}]
\item constancy of the residue or standard part on each active infinitesimal fibre;
\item compatibility of the induced support map on the Boolean star of the base point.
\end{enumerate}
Neither condition implies the other.
\end{theorem}

\begin{proof}
The cubical decomposition of Theorem \ref{thm:cubical-monad} separates a neighbourhood into disjoint strata indexed by support activation masks.  Within each stratum, the residue map forgets infinitesimal amplitudes and tests condition (a).  The support character forgets all amplitudes and tests condition (b).  A function can preserve residues while changing a support mask, or preserve support while changing residues, so the conditions are independent.
\end{proof}

\section{Identity levels and univalence}

There are three mathematically distinct senses in which two infinitesimal objects may be called identical.

\begin{definition}[Three identity levels]\label{def:identity-levels}
Let $S=G(K)$.
\begin{enumerate}[label=\textup{(\roman*)}]
\item \emph{Arithmetic identity} is equality in $S$ or in a specified semiring quotient.
\item \emph{Support-modal identity} is equality after applying $\chi_K:S\to\B$.
\item \emph{Structural identity} is isomorphism or equivalence of pointed monads, formal neighbourhoods, or projective local geometries.
\end{enumerate}
\end{definition}

\begin{proposition}[Arithmetic and modal identities]\label{prop:arithmetic-modal}
For every $h\in K^\times$,
\[
e\neq h
\]
in $S$, but
\[
\chi_K(e)=\chi_K(h)=1.
\]
Freely adjoining the arithmetic identity $e=h$ yields $\B$, while adjoining it only on coordinates in $A$ yields
\[
\B^A\times S^{I\setminus A}.
\]
\end{proposition}

\begin{proof}
The inequality is set-theoretic because $e=0_K$ and $h\neq0$.  The support equality follows from the definition of $\chi_K$.  The quotient statements are Theorems \ref{thm:strict-identification-collapse} and \ref{thm:scoped-identity}.
\end{proof}

\begin{proposition}[Structural identity of zero and infinity monads]\label{prop:structural-identity}
Projective inversion gives an isomorphism of pointed split monads
\[
(\mu^\square(0_e),0_e)
\xrightarrow{\sim}
(\mu^\square(\infty_e),\infty_e)
\]
and likewise for the unsupported points $0_\tau$ and $\infty_\tau$.
\end{proposition}

\begin{proof}
This is Theorem \ref{thm:zero-infinity-monads}, with the distinguished centres included in the data.
\end{proof}

\begin{theorem}[Identity in the univalent completion]\label{thm:univalent-interpretation}
Let $\widehat{\mathcal M}_{\mathrm{sp}}$ be the univalent completion of
the groupoid of pointed split monads and pointed equivalences.  The
projective inversion equivalence of Proposition
\ref{prop:structural-identity} determines the identity path
\[
[(\mu^\square(0_e),0_e)]
=
[(\mu^\square(\infty_e),\infty_e)]
\]
in $\widehat{\mathcal M}_{\mathrm{sp}}$.
The arithmetic elements remain separated by their multiplicative laws:
\[
e^2=e,
\qquad
(\omega^{-1})^2=\omega^{-2}\neq\omega^{-1}.
\]
\end{theorem}

\begin{proof}
The Rezk--univalent completion sends equivalences in the original groupoid
to identity paths; see \cite{HoTT}.  Apply this universal property to the
projective inversion equivalence.  The final formulas are computed in
$G(\No[i])$ and show that the two amplitude elements occupy different
multiplicative structures.  The quotient obtained by imposing their
arithmetic equality is the Boolean quotient of
Theorem~\ref{thm:strict-identification-collapse}.
\end{proof}

\begin{remark}[Identity is indexed by structure]
The same pair of objects can be equal in the univalent completion of
pointed monads, equal after support truncation, and unequal as semiring
elements.  These are three functorially related identity types rather than
competing uses of one equality symbol.
\end{remark}

\part{Arithmetic curves, adelic balance, and Frobenius perfection}

\section{The absolute arithmetic curve and the split Abel--Jacobi fibre}
\label{sec:absolute-curve}

Three coefficient geometries must be distinguished.
\begin{enumerate}[label=\textup{(\roman*)}]
\item The original arithmetic site of Connes--Consani is
\[
\mathfrak A=(\widehat{\N^\times},\Nbar),
\qquad
\Nbar=(\N\cup\{\infty\},\min,+),
\]
a semiringed topos of characteristic one \cite{CCArithmeticSite}.
\item The absolute $\Fone$-site is
\[
\operatorname{Pic}_*
 =\bigl(\widehat{\N_0^\times},\cO_{\Fone}\bigr),
\qquad
\cO_{\Fone}=\Fone[T],
\]
where $\Fone[T]$ is a spherical algebra rather than an ordinary semiring
object \cite{CCAbsolute}.
\item Extension of scalars of the arithmetic site produces the scaling
site, whose structure sheaf is idempotent and encodes valuation profiles
\cite{CCScalingSite}.
\end{enumerate}
The point and stalk results below concern the absolute $\Fone$-site.  The
sheafwise split construction of Section~\ref{sec:split-scaling-site}
applies directly to the ordinary semiring objects $\Nbar$ and the scaling
sheaf; it does not silently identify spherical $\Fone$-algebras with
Boolean semirings.  The facts recalled in the next theorem are imported
from \cite{CCAbsolute,CCArithmeticSite}; they fix the precise interface
with split support.

\begin{theorem}[Connes--Consani point and stalk package]
\label{thm:cc-package}
The following hold.
\begin{enumerate}[label=\textup{(\roman*)}]
\item A point of $\widehat{\N_0^\times}$ is represented by a totally ordered
rank-one abelian group $H$ isomorphic to a subgroup of $\Q$, together with
the distinguished trivial point $\{0\}$.
\item The stalk of $\cO_{\Fone}$ at $H$ is
\[
A_H=\Fone[T^{H_+}].
\]
\item There is a geometric morphism
\[
\Theta:\Spec\Z\longrightarrow\widehat{\N_0^\times}
\]
whose map on points satisfies
\[
\Theta(\eta)=[\{0\}],
\qquad
\Theta(p)=[\Z[1/p]]
\]
for the generic point $\eta$ and every rational prime $p$.
\item There is a surjective map on point classes
\[
\kappa:\operatorname{Pt}(\widehat{\N_0^\times})
\longrightarrow \operatorname{Pic}(\Spec\Z)
\]
such that
\[
\kappa\circ\Theta=\theta,
\]
where $\theta$ is the arithmetic Abel--Jacobi map.  The two point objects
$\{0\}$ and $\Z$ are nonisomorphic, but both are sent by $\kappa$ to the
zero divisor class.
\end{enumerate}
\end{theorem}

Let
\[
S=G(\Z)=\Z\sqcup\{\tau\},
\qquad e=0_{\Z}\in S.
\]
Its prime spectrum is
\[
P_\tau=\{\tau\},
\qquad
P_0=\{\tau,e\},
\qquad
P_p=p\Z\cup\{\tau\}
\]
for rational primes $p$.  Recall that
\[
P_\tau\subsetneq P_0\subsetneq P_p.
\]

\begin{proposition}[Classical collapse of the split spectrum]
\label{prop:classical-collapse}
The assignment
\[
q:\Spec S\longrightarrow\Spec\Z,
\qquad
q(P_\tau)=q(P_0)=\eta,
\qquad
q(P_p)=p,
\]
is continuous.
\end{proposition}

\begin{proof}
For every nonzero $n\in\Z$, the usual basic open $D_{\Z}(n)$ has inverse
image
\[
q^{-1}(D_{\Z}(n))=D_S(n).
\]
Indeed, $P_\tau$ and $P_0$ both map to the generic point and neither
contains a nonzero integer, while $P_p$ belongs to $D_S(n)$ precisely when
$p\nmid n$.  The opens $D_{\Z}(n)$ with $n\neq0$, together with the empty
open, form a basis.  Hence $q$ is continuous.
\end{proof}

\begin{definition}[Split absolute point map]
\label{def:split-absolute-point-map}
Define a map on point classes
\[
\widetilde\Theta_{\mathrm{pt}}:\Spec S
\longrightarrow \operatorname{Pt}(\widehat{\N_0^\times})
\]
by
\[
\widetilde\Theta_{\mathrm{pt}}(P_\tau)=\{0\},
\qquad
\widetilde\Theta_{\mathrm{pt}}(P_0)=\Z,
\qquad
\widetilde\Theta_{\mathrm{pt}}(P_p)=\Z[1/p].
\]
\end{definition}

\begin{theorem}[Resolution of the zero Abel--Jacobi fibre]
\label{thm:abel-jacobi-resolution}
At the level of point classes one has a commutative diagram
\[
\begin{tikzcd}[column sep=large,row sep=large]
\Spec G(\Z)
  \arrow[r,"\widetilde\Theta_{\mathrm{pt}}"]
  \arrow[d,"q"']
& \operatorname{Pt}(\widehat{\N_0^\times})
  \arrow[d,"\kappa"]\\
\Spec\Z \arrow[r,"\theta"']
& \operatorname{Pic}(\Spec\Z).
\end{tikzcd}
\]
Moreover,
\[
\widetilde\Theta_{\mathrm{pt}}(P_\tau)
 \not\cong
\widetilde\Theta_{\mathrm{pt}}(P_0),
\]
whereas
\[
\kappa\bigl(\widetilde\Theta_{\mathrm{pt}}(P_\tau)\bigr)
=
\kappa\bigl(\widetilde\Theta_{\mathrm{pt}}(P_0)\bigr).
\]
Thus the split spectrum separates two absolute point objects lying over the
classical zero divisor class.
\end{theorem}

\begin{proof}
At $P_p$, both composites give the divisor class represented by
$\Z[1/p]$.  At $P_\tau$ and $P_0$, the lower composite is $\theta(\eta)$.
By Theorem \ref{thm:cc-package},
\[
\kappa(\{0\})=[\Z]=\kappa(\Z)=\theta(\eta).
\]
This proves commutativity.  The trivial group $\{0\}$ and the infinite
cyclic ordered group $\Z$ are nonisomorphic, proving the strict
separation.
\end{proof}

\begin{corollary}[The resolved stalk chain]
\label{cor:resolved-stalk-chain}
The point chain
\[
P_\tau\subset P_0\subset P_p
\]
is sent to the ordered-group chain
\[
\{0\}\hookrightarrow\Z\hookrightarrow\Z[1/p]
\]
and to the spherical-algebra chain
\[
\Fone
\longrightarrow
\Fone[T]
\longrightarrow
\Fone[T^{\Z[1/p]_+}].
\]
The last term is perfect for the $p$-th Frobenius and generally imperfect
for the other prime Frobenii.
\end{corollary}

\begin{proof}
The ordered-group arrows are the unique map from the trivial group and the
standard inclusion $\Z\subseteq\Z[1/p]$.  Applying
$H\mapsto\Fone[T^{H_+}]$ gives the displayed algebra maps.  Multiplication
by $p$ is bijective on $\Z[1/p]_+$, with inverse multiplication by $1/p$.
\end{proof}

\begin{remark}[Point-and-stalk realization]
Theorem~\ref{thm:abel-jacobi-resolution} gives a morphism of specialization
posets together with the compatible stalk diagram.  This point-and-stalk
realization is the functorial level used in the denominator and Frobenius
constructions below.
\end{remark}


\section{The arithmetic Jacobian as an unreduced split-support monoid}
\label{sec:arithmetic-jacobian-split}

We now use the completed arithmetic Picard and Jacobian monoids of
Connes--Consani \cite{CCJacobian}.  To avoid notational ambiguity, let
$\Picf$ denote the monoid of isomorphism classes of torsion-free rank-one
abelian groups, let $\Picar$ denote the monoid of such groups equipped with
a possibly zero archimedean seminorm, and let
\[
\Jacar:=\Picar/\R_+^\times.
\]

\begin{theorem}[Connes--Consani configuration and Jacobian package]
\label{thm:cc-jacobian-package}
The following statements hold.
\begin{enumerate}[label=\textup{(\roman*)}]
\item The map
\[
\Theta_{\mathrm{conf}}:\mathcal P(\mathbb P_f)\longrightarrow\Picf,
\qquad
S\longmapsto\Z[S^{-1}],
\]
where $\mathbb P_f$ is the set of finite rational primes, is an isomorphism
from the union semilattice onto $\Idem(\Picf)$.
\item There is a canonical monoid isomorphism
\[
\Jacar\cong\Picf\times\{0,\infty\},
\]
where $0$ denotes a nondegenerate metric orbit, $\infty$ the zero-seminorm
fixed stratum, and the second factor is multiplied by
$0+0=0$ and $0+\infty=\infty+\infty=\infty$.
\item The extended Abel--Jacobi map sends
\[
\eta\longmapsto(\Z,0),
\qquad
p\longmapsto(\Z[1/p],0),
\qquad
\infty\longmapsto(\Z,\infty).
\]
\end{enumerate}
\end{theorem}

\begin{proof}
These are Theorems 2.15--2.16, Definition 4.1 and the structural
identification of Section 4, and the extended Abel--Jacobi formula of
\cite{CCJacobian}.  They are collected here with notation separating the
finite Picard monoid from its completed metric version.
\end{proof}

Identify $\{0,\infty\}$ with the multiplicative Boolean monoid by
\[
0\longmapsto1,
\qquad
\infty\longmapsto0.
\]
The rule for the archimedean coordinate then becomes Boolean
multiplication.

\begin{theorem}[Exact split-support form of the arithmetic Jacobian]
\label{thm:jacobian-boolean-product}
There is a canonical commutative-monoid isomorphism
\[
\boxed{\Jacar\cong\Picf\times\B.}
\]
Consequently,
\[
\boxed{
\Idem(\Jacar)
\cong
\mathcal P(\mathbb P_f)\times\B
\cong
\mathcal P(\mathbb P_f\sqcup\{\infty\})
}
\]
as join-semilattices.  Under the last isomorphism,
$(S,1)$ corresponds to $S$ and $(S,0)$ to $S\cup\{\infty\}$.
\end{theorem}

\begin{proof}
The first assertion is Theorem~\ref{thm:cc-jacobian-package}(ii) after the
stated identification of the two-element factors.  An element of a product
monoid is idempotent exactly when both components are idempotent.  Every
Boolean element is idempotent, and
Theorem~\ref{thm:cc-jacobian-package}(i) identifies the idempotents of
$\Picf$ with prime configurations.  Finally, Boolean multiplication in the
second coordinate becomes union of the distinguished symbol $\infty$ under
$0\leftrightarrow\{\infty\}$ and $1\leftrightarrow\varnothing$.
\end{proof}

\begin{proposition}[Boolean truncation of a degenerate Arakelov metric]
\label{prop:arakelov-support-truncation}
Let
\[
\sigma:\R_{\geq0}\longrightarrow\B,
\qquad
\sigma(t)=
\begin{cases}
0,&t=0,\\
1,&t>0.
\end{cases}
\]
Then $\sigma$ is a semiring homomorphism from
$(\R_{\geq0},+,\cdot)$ to $(\B,\vee,\wedge)$.  If $\ell$ is a real
one-dimensional vector space with a possibly zero seminorm $\|\cdot\|$,
the value
\[
\epsilon(\ell,\|\cdot\|):=\sigma(\|v\|),
\qquad 0\neq v\in\ell,
\]
is independent of $v$, invariant under positive rescaling of the
seminorm, and satisfies
\[
\epsilon(\ell\otimes m)=\epsilon(\ell)\epsilon(m).
\]
Under Theorem~\ref{thm:jacobian-boolean-product}, this is precisely the
projection $\Jacar\to\B$.
\end{proposition}

\begin{proof}
For nonnegative reals, $x+y=0$ exactly when $x=y=0$, and $xy>0$ exactly
when both factors are positive; hence $\sigma$ preserves addition and
multiplication.  A seminorm on a one-dimensional real vector space is
either identically zero or positive on every nonzero vector, so
$\epsilon$ is well defined.  Positive rescaling does not change its
support.  The tensor-product seminorm is zero exactly when at least one
factor seminorm is zero, proving multiplicativity.  Connes--Consani's two
archimedean orbits are precisely the nondegenerate and zero-seminorm
orbits, which are identified with $1$ and $0$ in
Theorem~\ref{thm:jacobian-boolean-product}.
\end{proof}

\begin{remark}[Metric data and its Boolean shadow]
The map $\sigma$ is the idempotent support truncation of the archimedean
metric orbit.  The logarithmic metric, Arakelov degree, and height live in
the quantitative fibre over the nondegenerate Boolean value $1$.
\end{remark}

\begin{corollary}[The Abel--Jacobi support vertices]
\label{cor:abel-jacobi-support-vertices}
Under the idempotent identification of
Theorem~\ref{thm:jacobian-boolean-product}, the extended Abel--Jacobi map
sends
\[
\eta\longmapsto\varnothing,
\qquad
p\longmapsto\{p\},
\qquad
\infty\longmapsto\{\infty\}.
\]
For any finite set $J$ of places, its configuration submonoid is the
literal Boolean cube $\mathcal P(J)$.
\end{corollary}

\begin{proof}
The group $\Z$ is the empty localization, $\Z[1/p]$ is the singleton
finite-prime localization, and the zero seminorm sets precisely the
archimedean Boolean coordinate to $0$.
\end{proof}

Let $\pi_{\infty}(S)=(S\cap\mathbb P_f,\epsilon_S)$, where $\epsilon_S=0$ exactly when $\infty\in S$.  The two support coordinates and the ambient Jacobian are summarized by
\[
\begin{tikzcd}[column sep=large,row sep=large]
\mathcal P(\mathbb P_f\sqcup\{\infty\})
  \arrow[r,"\sim"] \arrow[d,"\pi_{\infty}"']
& \Idem(\Jacar) \arrow[d,hook]\\
\mathcal P(\mathbb P_f)\times\B
  \arrow[r,"\sim"']
& \Picf\times\B\cong\Jacar.
\end{tikzcd}
\]

The preceding product retains the group label even on the degenerate
stratum.  This is an unreduced split completion, whereas $G(R)$ has only
one new unsupported zero.

\begin{definition}[Unreduced and reduced multiplicative split]
\label{def:unreduced-multiplicative-split}
For a commutative monoid $M$, define
\[
\operatorname{USpl}(M):=M\times\B.
\]
Define the pointed split monoid
\[
M^\tau:=M\sqcup\{\tau\},
\]
where $\tau$ is absorbing.  Let
\[
c_M:\operatorname{USpl}(M)\to M^\tau,
\qquad
c_M(m,1)=m,
\quad
c_M(m,0)=\tau.
\]
\end{definition}

\begin{proposition}[Reduction of the degenerate stratum]
\label{prop:jacobian-reduction}
Let $\sim_0$ be the Rees congruence on $\operatorname{USpl}(M)$ whose only
non-singleton class is $M\times\{0\}$.  The map $c_M$ is a surjective
monoid homomorphism with kernel congruence $\sim_0$.  Hence
\[
\operatorname{USpl}(M)/{\sim_0}\cong M^\tau.
\]
For $M=\Picf$, the arithmetic Jacobian is the unreduced split completion
and its quotient by the whole zero-seminorm stratum is the reduced pointed
split monoid $\Picf^\tau$.
\end{proposition}

\begin{proof}
For $(m,\epsilon),(n,\delta)\in M\times\B$,
\[
c_M((m,\epsilon)(n,\delta))
=c_M(mn,\epsilon\delta).
\]
If $\epsilon\delta=1$, both factors have Boolean value $1$ and the image is
$mn$; if $\epsilon\delta=0$, at least one image is $\tau$ and both sides
are $\tau$.  Thus $c_M$ is multiplicative.  The remaining assertions are
immediate from its fibres.
\end{proof}

\begin{remark}[Exact comparison morphism]
Theorem~\ref{thm:jacobian-boolean-product} is the Connes--Consani Jacobian
theorem in Boolean coordinates.  Its relation to split zero is the explicit
multiplicative quotient
\[
\Jacar\cong\operatorname{USpl}(\Picf)
\xrightarrow{\ c_{\Picf}\ }
\Picf^\tau,
\]
obtained after forgetting the additive structures of the coefficient
semirings.
\end{remark}

\section{Adelic curves, measure-algebra support, and Arakelov defects}
\label{sec:adelic-support-defect}

The split-support lattice has a canonical quantitative realization on an
adelic curve.  The appropriate support object is not the set of places
itself but its measure algebra, and the product formula defines a vector
measure on that Boolean algebra.  This section makes that statement precise
and separates three notions which must not be conflated: an absent local
datum, a zero seminorm, and a genuine trivial absolute value.

Let
\[
S=(K,(\Omega,\mathcal A,\nu),\phi)
\]
be an adelic curve in the sense of Chen--Moriwaki
\cite{ChenMoriwakiBook,ChenMoriwakiIntersection}.  Thus
$\phi(\omega)=|\cdot|_\omega$ is an absolute value of $K$ for every
$\omega\in\Omega$, and $\omega\mapsto\log|a|_\omega$ is integrable for
every $a\in K^\times$.  The curve is proper when
\[
\int_\Omega\log|a|_\omega\,\nu(d\omega)=0
\qquad(a\in K^\times).
\]

\subsection{The measure support lattice and its defect valuation}

Let
\[
\mathcal N_\nu:=\{A\in\mathcal A:\nu(A)=0\}
\]
and write
\[
\mathscr B_S:=\mathcal A/\!\sim_\nu,
\qquad
A\sim_\nu B\iff\nu(A\mathbin\triangle B)=0.
\]
With union, intersection, and complement, $\mathscr B_S$ is a Boolean
$\sigma$-algebra.  We write $[A]$ for the class of a measurable subset
$A\subseteq\Omega$.

\begin{definition}[Adelic support defect]
\label{def:adelic-support-defect}
For $[A]\in\mathscr B_S$, define
\[
\Delta_S([A]):K^\times\longrightarrow\R,
\qquad
\Delta_S([A])(a)
:=\int_A\log|a|_\omega\,\nu(d\omega).
\]
We call $[A]$ \emph{balanced} if $\Delta_S([A])=0$ as a homomorphism
$K^\times\to\R$.
\end{definition}

\begin{theorem}[The support-defect valuation]
\label{thm:adelic-support-defect}
For every adelic curve $S$, the following statements hold.
\begin{enumerate}[label=\textup{(\arabic*)}]
\item $\Delta_S([A])$ is a well-defined group homomorphism
\[
K^\times\longrightarrow(\R,+).
\]
\item The map
\[
\Delta_S:\mathscr B_S\longrightarrow
\operatorname{Hom}_{\mathrm{gp}}(K^\times,\R)
\]
is a valuation of Boolean lattices:
\[
\Delta_S(A\cup B)+\Delta_S(A\cap B)
=\Delta_S(A)+\Delta_S(B).
\]
In particular, it is finitely additive on disjoint supports.
\item If $S$ is proper, then
\[
\Delta_S(\Omega)=0,
\qquad
\Delta_S(A^c)=-\Delta_S(A).
\]
\item The restricted measured family
\[
S|_A=(K,(A,\mathcal A|_A,\nu|_A),\phi|_A)
\]
is a proper adelic curve if and only if $[A]$ is balanced.
\end{enumerate}
\end{theorem}

\begin{proof}
If $A\sim_\nu B$, then the integrals over $A$ and $B$ agree because the
integrand is integrable and $A\mathbin\triangle B$ is null.  Multiplicativity
of an absolute value gives
\[
\log|ab|_\omega=\log|a|_\omega+\log|b|_\omega,
\]
so $\Delta_S(A)$ is a group homomorphism.  The pointwise identity
\[
\mathbf 1_{A\cup B}+\mathbf 1_{A\cap B}
=\mathbf 1_A+\mathbf 1_B
\]
gives the valuation formula after multiplication by
$\log|a|_\omega$ and integration.  Properness gives
$\Delta_S(\Omega)=0$, and the complement formula follows from
$\mathbf 1_{A^c}=1-\mathbf 1_A$.  Finally, the restricted logarithms are
integrable because the original logarithms are, and the product formula on
$S|_A$ is exactly the condition $\Delta_S(A)=0$.
\end{proof}


\begin{theorem}[Linearization of adelic support]
\label{thm:adelic-defect-linearization}
Let $L^\infty_S=L^\infty(\Omega,\mathcal A,\nu;\R)$.  Define
\[
\mathscr D_S:L^\infty_S\longrightarrow
\Hom_{\mathrm{gp}}(K^\times,\R),
\qquad
\mathscr D_S(f)(a)
=\int_\Omega f(\omega)\log|a|_\omega\,\nu(d\omega).
\]
Then:
\begin{enumerate}[label=\textup{(\arabic*)}]
\item $\mathscr D_S$ is a well-defined real-linear map;
\item the idempotents of the commutative algebra $L^\infty_S$ are exactly
the characteristic functions $\mathbf1_A$ modulo null sets, and
\[
\mathscr D_S(\mathbf1_A)=\Delta_S([A]);
\]
\item if $S$ is proper, then $\R\mathbf1_\Omega\subseteq\ker\mathscr D_S$;
\item the induced map
\[
\overline{\mathscr D}_S:
\mathfrak D_S:=L^\infty_S/\ker\mathscr D_S
\longrightarrow
\Hom_{\mathrm{gp}}(K^\times,\R)
\]
is injective with image $\operatorname{im}(\mathscr D_S)$.
\end{enumerate}
The vector space $\mathfrak D_S$ is the quantitative defect space of the
adelic curve, and $\mathscr B_S$ is its Boolean idempotent skeleton.
\end{theorem}

\begin{proof}
For $f\in L^\infty_S$ and $a\in K^\times$,
\[
\int_\Omega |f(\omega)\log|a|_\omega|\,d\nu
\leq
\|f\|_\infty
\int_\Omega|\log|a|_\omega|\,d\nu<\infty,
\]
so the integral is defined.  Linearity in $f$ and the identity
$\log|ab|_\omega=\log|a|_\omega+\log|b|_\omega$ prove (1).
An idempotent real-valued measurable function satisfies $f^2=f$ almost
everywhere and therefore takes values in $\{0,1\}$ almost everywhere;
this gives (2).  Properness is precisely
$\mathscr D_S(\mathbf1_\Omega)=0$, which proves (3).  Statement (4) is the
first isomorphism theorem for vector spaces.
\end{proof}

\begin{theorem}[Principal-logarithm duality]
\label{thm:principal-logarithm-duality}
For $a\in K^\times$, put
\[
\ell_a(\omega)=\log|a|_\omega
\]
and let
\[
\mathscr L_S
=\operatorname{span}_{\R}\{\ell_a:a\in K^\times\}
\subseteq L^1(\Omega,\nu).
\]
The formula
\[
\langle[f],g\rangle_S=\int_\Omega f(\omega)g(\omega)\,d\nu
\]
defines a nondegenerate bilinear pairing
\[
\boxed{
\mathfrak D_S\times\mathscr L_S\longrightarrow\R.
}
\]
Under this pairing,
\[
\langle[f],\ell_a\rangle_S
=\overline{\mathscr D}_S([f])(a).
\]
Thus the quantitative defect space is the separating dual object of the
space generated by principal logarithms.
\end{theorem}

\begin{proof}
If $f$ is changed by an element of $\ker\mathscr D_S$, then its integral
against every generator $\ell_a$, and hence against every element of
$\mathscr L_S$, is unchanged.  The pairing is therefore well defined.
If $[f]$ pairs trivially with $\mathscr L_S$, then
$\mathscr D_S(f)(a)=0$ for all $a$, so $[f]=0$ in $\mathfrak D_S$.
Conversely, if $0\neq g\in\mathscr L_S\subseteq L^1$, the bounded
measurable function $\operatorname{sgn}(g)$ satisfies
\[
\int_\Omega \operatorname{sgn}(g)g\,d\nu
=\int_\Omega|g|\,d\nu>0.
\]
Hence the pairing is nondegenerate in the second variable as well.
\end{proof}

\begin{definition}[Bounded weighted adelic height]
\label{def:weighted-adelic-height}
Let $n\geq0$, let
$\mathbf x=(x_0,\dots,x_n)\in K^{n+1}\setminus\{0\}$, and let
$f\in L^\infty(\Omega,\nu)$.  Define
\[
h_{S,f}^{(n)}(\mathbf x)
:=\int_\Omega f(\omega)
\log\!\left(\max_{0\leq i\leq n}|x_i|_\omega\right)d\nu.
\]
For $f=\mathbf1_A$ we write $h_{S,A}^{(n)}$.
\end{definition}

\begin{theorem}[Projective descent is controlled by the defect transform]
\label{thm:weighted-height-projective-descent}
For every $f\in L^\infty(\Omega,\nu)$ and every
$\mathbf x\in K^{n+1}\setminus\{0\}$, the weighted height is finite and
satisfies
\[
\boxed{
h_{S,f}^{(n)}(a\mathbf x)
=h_{S,f}^{(n)}(\mathbf x)+\mathscr D_S(f)(a)
\qquad(a\in K^\times).
}
\]
Consequently, the following conditions are equivalent:
\begin{enumerate}[label=\textup{(\roman*)}]
\item $h_{S,f}^{(n)}$ descends to a function on $\mathbf P^n(K)$;
\item $f\in\ker\mathscr D_S$;
\item $[f]=0$ in $\mathfrak D_S$.
\end{enumerate}
For measurable supports $A,B$ one also has
\[
h_{S,A\cup B}^{(n)}+h_{S,A\cap B}^{(n)}
=h_{S,A}^{(n)}+h_{S,B}^{(n)}.
\]
Hence a support face $A$ carries a projectively well-defined partial
height exactly when $A$ is balanced.
\end{theorem}

\begin{proof}
Discard the zero coordinates of $\mathbf x$.  The remaining finitely many
functions $\log|x_i|_\omega$ lie in $L^1(\Omega,\nu)$, and
\[
\left|\max_i\log|x_i|_\omega\right|
\leq\sum_i\left|\log|x_i|_\omega\right|.
\]
Multiplication by the bounded function $f$ therefore gives an integrable
height integrand.  Since
\[
\max_i|a x_i|_\omega
=|a|_\omega\max_i|x_i|_\omega,
\]
the displayed transformation law follows after taking logarithms and
integrating.  It is invariant under every projective rescaling precisely
when $\mathscr D_S(f)=0$.  The Boolean valuation formula follows from
$\mathbf1_{A\cup B}+\mathbf1_{A\cap B}
=\mathbf1_A+\mathbf1_B$.
\end{proof}

\begin{corollary}[Weighted Arakelov degree descent]
\label{cor:weighted-degree-descent}
Let $(E,\xi)$ be a nonzero strongly adelic vector bundle on $S$.  For
$f\in L^\infty(\Omega,\nu)$ and $0\neq s\in E$, define
\[
\widehat{\deg}_{\xi,f}(s)
:=-\int_\Omega f(\omega)\log\|s\|_\omega\,d\nu.
\]
Then
\[
\widehat{\deg}_{\xi,f}(as)
=\widehat{\deg}_{\xi,f}(s)-\mathscr D_S(f)(a).
\]
Therefore $\widehat{\deg}_{\xi,f}$ descends to the projective space
$\mathbf P(E)(K)$ exactly when $f\in\ker\mathscr D_S$.
\end{corollary}

\begin{proof}
Strong adelicity makes $\log\|s\|_\omega$ integrable.  The norm identity
$\|as\|_\omega=|a|_\omega\|s\|_\omega$ gives the transformation formula,
and projective descent follows exactly as in
Theorem~\ref{thm:weighted-height-projective-descent}.
\end{proof}

\begin{definition}[Measure-preserving refinement of place families]
\label{def:adelic-place-refinement}
Let $S'=(K,(\Omega',\mathcal A',\nu'),\phi')$ and
$S=(K,(\Omega,\mathcal A,\nu),\phi)$ be adelic curves over the same field.
A \emph{measure-preserving refinement} is a measurable map
$\pi:\Omega'\to\Omega$ such that
\[
\pi_*\nu'=\nu,
\qquad
|a|'_{\omega'}=|a|_{\pi(\omega')}
\quad(a\in K,\ \omega'\in\Omega').
\]
\end{definition}

\begin{theorem}[Naturality of the defect transform]
\label{thm:adelic-defect-naturality}
For a measure-preserving refinement $\pi:S'\to S$, pullback defines
\[
\pi^*:L^\infty_S\longrightarrow L^\infty_{S'},
\qquad
\pi^*f=f\circ\pi,
\]
and the diagram
\[
\begin{tikzcd}[column sep=large]
L^\infty_S \arrow[r,"\mathscr D_S"] \arrow[d,"\pi^*"']
& \Hom_{\mathrm{gp}}(K^\times,\R) \arrow[d,equal]\\
L^\infty_{S'} \arrow[r,"\mathscr D_{S'}"']
& \Hom_{\mathrm{gp}}(K^\times,\R)
\end{tikzcd}
\]
commutes.  In particular,
\[
\Delta_{S'}(\pi^{-1}A)=\Delta_S(A),
\]
and balanced supports pull back to balanced supports.
\end{theorem}

\begin{proof}
For $a\in K^\times$, the change-of-variables formula for the pushforward
measure gives
\[
\begin{aligned}
\mathscr D_{S'}(f\circ\pi)(a)
&=\int_{\Omega'}f(\pi(\omega'))
   \log|a|'_{\omega'}\,d\nu'(\omega')\\
&=\int_{\Omega'}f(\pi(\omega'))
   \log|a|_{\pi(\omega')}\,d\nu'(\omega')\\
&=\int_\Omega f(\omega)\log|a|_\omega\,d\nu(\omega)
 =\mathscr D_S(f)(a).
\end{aligned}
\]
Apply this to $f=\mathbf1_A$.
\end{proof}

\begin{corollary}[Adelic realization of lattice-valued split support]
\label{cor:adelic-lattice-support}
Let $R$ be a commutative ring and form the lattice-split semiring
$G_{\mathscr B_S}(R)$.  Its ring-null fibre is the Boolean measure algebra
$\mathscr B_S$, and for $z_A=(0,[A])$ one has
\[
G_{\mathscr B_S}(R)[z_A^{-1}]\cong\downarrow[A].
\]
The composite support datum is therefore
\[
\begin{tikzcd}[column sep=large]
G_{\mathscr B_S}(R)
  \arrow[r,"\chi_S"] \arrow[dr,"p_S"']
& \mathscr B_S \arrow[r,"\Delta_S"]
& \operatorname{Hom}_{\mathrm{gp}}(K^\times,\R)\\
& R &
\end{tikzcd}
\]
where $\chi_S$ records the measurable support face and $\Delta_S$ records
its exact product-formula defect.
\end{corollary}

\begin{proof}
Apply Theorems~\ref{thm:lattice-normal-form} and
\ref{thm:lattice-localization} to $L=\mathscr B_S$, and then apply
Theorem~\ref{thm:adelic-support-defect}.
\end{proof}


\begin{theorem}[Balanced support localization]
\label{thm:balanced-support-localization}
Let $S=(K,(\Omega,\mathcal A,\nu),\phi)$ be a proper adelic curve and let
$[A]\in\mathscr B_S$.  Then the map
\[
\downarrow[A]\longrightarrow\mathscr B_{S|_A},
\qquad
[B]\longmapsto[B\cap A],
\]
is a Boolean-algebra isomorphism.  Consequently,
\[
\boxed{
G_{\mathscr B_S}(K)[z_A^{-1}]
\cong\mathscr B_{S|_A}.
}
\]
Under this identification, the following conditions are equivalent:
\begin{enumerate}[label=\textup{(\roman*)}]
\item $[A]$ is balanced;
\item the restricted adelic curve $S|_A$ is proper;
\item the weighted height $h_{S,A}^{(n)}$ descends to $\mathbf P^n(K)$ for
every $n\ge0$.
\end{enumerate}
If $\pi:S'\to S$ is a measure-preserving refinement, pullback of supports
induces a commutative localization diagram
\[
\begin{tikzcd}[column sep=large]
G_{\mathscr B_S}(K)[z_A^{-1}]
  \arrow[r] \arrow[d,"\sim"']
& G_{\mathscr B_{S'}}(K)[z_{\pi^{-1}A}^{-1}]
  \arrow[d,"\sim"]\\
\mathscr B_{S|_A} \arrow[r,"\pi^{-1}"']
& \mathscr B_{S'|_{\pi^{-1}A}}.
\end{tikzcd}
\]
\end{theorem}

\begin{proof}
Every measurable subset of the restricted space $A$ is represented by
$B\cap A$ for some $B\in\mathcal A$, and two such subsets agree modulo the
restricted null ideal exactly when their classes agree below $[A]$ in
$\mathscr B_S$.  This proves the Boolean-algebra isomorphism.  The
localization formula follows from
Theorem~\ref{thm:lattice-localization}.  The equivalence of (i) and (ii) is
Theorem~\ref{thm:adelic-support-defect}; the equivalence with (iii) is
Theorem~\ref{thm:weighted-height-projective-descent}.  Naturality follows
from pullback of measurable subsets and
Theorem~\ref{thm:adelic-defect-naturality}.
\end{proof}

\begin{corollary}[Rigidity of proper support localizations for global fields]
\label{cor:global-balanced-localization-rigidity}
For the standard proper adelic curve of a global field, the only support
localizations yielding proper restricted adelic curves are the empty and
full localizations.
\end{corollary}

\begin{proof}
Combine Theorems~\ref{thm:balanced-support-localization} and
\ref{thm:global-field-balanced-supports}.
\end{proof}

\begin{theorem}[Global-field rigidity of balanced supports]
\label{thm:global-field-balanced-supports}
Let $K$ be a global field and let $S_K$ be its standard proper adelic
curve with normalized local absolute values and their standard weights.
A subset $A$ of the place set of $K$ is balanced if and only if
\[
A=\varnothing
\qquad\text{or}\qquad
A=\Omega_K.
\]
\end{theorem}

\begin{proof}
We treat number fields and global function fields separately.

\emph{Number-field case.}
For each place $v$, let $c_v=\mathbf1_A(v)$.  At the archimedean places,
the balance relation applied to units $u\in\mathcal O_K^\times$ is
\[
\sum_{v\mid\infty}c_v n_v\log|u|_v=0.
\]
By the Dirichlet unit theorem, the vectors
$(n_v\log|u|_v)_{v\mid\infty}$ span the hyperplane whose coordinate sum
is zero.  Its annihilator is the line of constant vectors.  Hence the
$c_v$ are constant on the archimedean places; write their common value as
$c\in\{0,1\}$.

Let $\mathfrak p$ be a finite prime.  Finiteness of the ideal class group
gives an integer $h>0$ and $\alpha\in K^\times$ such that
$(\alpha)=\mathfrak p^h$.  Every finite local term of $\alpha$ vanishes
except the $\mathfrak p$-term.  The full product formula therefore gives
\[
\sum_{v\mid\infty}n_v\log|\alpha|_v
=-n_{\mathfrak p}\log|\alpha|_{\mathfrak p},
\]
and the right-hand side is nonzero.  The balanced relation restricted to
$A$ becomes
\[
(c_{\mathfrak p}-c)
 n_{\mathfrak p}\log|\alpha|_{\mathfrak p}=0.
\]
Thus $c_{\mathfrak p}=c$.  Since $\mathfrak p$ was arbitrary, all place
indicators equal $c$.

\emph{Function-field case.}
Write $K=\mathbf F_q(C)$ for a smooth projective geometrically integral
curve $C$.  Put $c_x=\mathbf1_A(x)$ for a closed point $x$ and define
\[
\ell_A:\operatorname{Div}(C)\longrightarrow\R,
\qquad
\ell_A\!\left(\sum_xm_x[x]\right)
=\sum_xm_xc_x\deg(x).
\]
Balance says exactly that $\ell_A$ vanishes on principal divisors, so it
descends to $\operatorname{Pic}(C)$.  The group
$\operatorname{Pic}^0(C)(\mathbf F_q)$ is finite, hence every homomorphism
from it to $\R$ is zero.  The degree map has finite kernel and image a
nonzero subgroup of $\Z$, so the descended homomorphism is
$\lambda\deg$ for some $\lambda\in\R$.  Evaluating on each prime divisor
$[x]$ gives
\[
c_x\deg(x)=\lambda\deg(x),
\]
so every $c_x$ equals the same value $\lambda\in\{0,1\}$.

The converse in both cases is the empty integral or the product formula.
\end{proof}

\begin{corollary}[The standard adelic curve of $\Q$]
\label{thm:q-balanced-supports}
For
\[
\Omega_\Q=\{\infty\}\sqcup\{p:p\text{ prime}\},
\]
the balanced supports are exactly $\varnothing$ and $\Omega_\Q$.
\end{corollary}

\begin{proof}
Apply Theorem~\ref{thm:global-field-balanced-supports} to $K=\Q$.
\end{proof}

\subsection{Height-neutral extensions by trivial places}

Let $(T,\mathcal T,\mu)$ be a measure space.  Denote by
$|\cdot|_{\mathrm{triv}}$ the trivial absolute value of $K$:
\[
|0|_{\mathrm{triv}}=0,
\qquad
|a|_{\mathrm{triv}}=1\quad(a\neq0).
\]

\begin{definition}[Trivial-place extension]
\label{def:trivial-place-extension}
Define
\[
S\boxplus T
:=
\bigl(K,(\Omega\sqcup T,\mathcal A\sqcup\mathcal T,
\nu\sqcup\mu),\phi\sqcup\phi_{\mathrm{triv}}\bigr),
\]
where $\phi_{\mathrm{triv}}(t)=|\cdot|_{\mathrm{triv}}$ for every
$t\in T$.
\end{definition}

\begin{theorem}[Neutrality of trivial-place extension]
\label{thm:trivial-place-neutrality}
The following statements hold.
\begin{enumerate}[label=\textup{(\arabic*)}]
\item $S\boxplus T$ is an adelic curve.  It is proper if and only if $S$ is
proper.
\item For every nonzero tuple $(a_0,\dots,a_n)\in K^{n+1}$,
\[
h_{S\boxplus T}(a_0,\dots,a_n)
=h_S(a_0,\dots,a_n).
\]
\item There is a canonical Boolean-$\sigma$-algebra isomorphism
\[
\mathscr B_{S\boxplus T}
\cong
\mathscr B_S\times\mathscr B_T,
\]
and under it
\[
\Delta_{S\boxplus T}(A,B)=\Delta_S(A).
\]
Hence every support contained in $T$ is balanced.
\item Under the canonical identification
\[
L^\infty(\Omega\sqcup T)\cong
L^\infty(\Omega)\oplus_\infty L^\infty(T),
\]
one has
\[
\mathscr D_{S\boxplus T}(f,g)=\mathscr D_S(f),
\]
and therefore
\[
\ker\mathscr D_{S\boxplus T}
=\ker\mathscr D_S\oplus_\infty L^\infty(T),
\qquad
\mathfrak D_{S\boxplus T}\cong\mathfrak D_S.
\]
\item Let $(E,\xi)$ be a strongly adelic vector bundle on $S$.  At every
$t\in T$, equip $E$ with the canonical trivial norm
\[
\|x\|_{\mathrm{triv}}
=
\begin{cases}
0,&x=0,\\
1,&x\neq0.
\end{cases}
\]
The extended norm family $\xi\boxplus\xi_{\mathrm{triv}}$ is strongly
adelic, and for every $0\neq s\in E$,
\[
\widehat{\deg}_{\xi\boxplus\xi_{\mathrm{triv}}}(s)
=
\widehat{\deg}_{\xi}(s).
\]
\end{enumerate}
\end{theorem}

\begin{proof}
For $a\in K^\times$, the logarithm of the trivial absolute value is the
zero function on $T$.  It is integrable for every measure $\mu$, and it
contributes zero to the product-formula integral.  This proves (1).
For a nonzero tuple, the maximum of the trivial absolute values of its
coordinates is $1$, so the additional height integrand is again zero,
which proves (2).

Every measurable subset of the disjoint union has a unique form
$A\sqcup B$, modulo the respective null ideals.  This gives the product
description of the measure algebra.  Since all logarithmic functions
vanish on $T$, the defect formula follows, proving (3).

The displayed trivial norm is ultrametric and is independent of a basis:
for any basis it is the associated maximum norm over the trivially valued
field.  If $g\in L^1(\Omega,\nu)$ dominates the distance from $\xi$ to a
basis norm, extend $g$ by zero on $T$.  The extended function remains
integrable and the new local distance is zero on $T$.  Measurability is
immediate because every global nonzero vector has norm $1$ on $T$.
Finally, the additional logarithmic norm is zero, proving the degree
identity and the final item.
\end{proof}

\begin{corollary}[The neutral Boolean cube over $\Q$]
\label{cor:q-neutral-support-cube}
For the standard adelic curve $S_\Q$ and any measure space $T$, the balanced
supports of $S_\Q\boxplus T$ are exactly
\[
\bigl\{\varnothing,\Omega_\Q\bigr\}\times\mathscr B_T.
\]
If $T$ is a finite set of $r$ positive-measure atoms, the added
height-neutral support directions form the Boolean cube $\B^r$.
\end{corollary}

\begin{proof}
Theorem~\ref{thm:trivial-place-neutrality} identifies the defect with its
$S_\Q$ component, and Theorem~\ref{thm:q-balanced-supports} classifies the
kernel on that component.
\end{proof}

\subsection{Global intersections and the exact contribution of a trivial place}

The preceding neutrality statements concern scalar heights and canonical
trivial norms.  A trivially valued place can carry nontrivial Berkovich
metrics and nonzero local intersection theory; this is an essential part of
Arakelov geometry over trivially valued fields
\cite{ChenMoriwakiDirichlet}.  The global effect is nevertheless exact.

\begin{theorem}[Variation of arithmetic intersection under a trivial-place
extension]
\label{thm:intersection-trivial-place-variation}
Let $X$ be a projective $K$-scheme of dimension $d$, and let
$\overline D'_0,\dots,\overline D'_d$ be integrable adelic Cartier divisors
on $X$ over $S\boxplus T$, with properly intersecting underlying Cartier
divisors.  Let $\overline D_i$ be their restrictions to $S$.  Then
\[
(\overline D'_0\cdots\overline D'_d)_{S\boxplus T}
=
(\overline D_0\cdots\overline D_d)_S
+
\int_T
(\overline D'_{0,t}\cdots\overline D'_{d,t})_t\,\mu(dt).
\]
Consequently, the global intersection number is unchanged exactly when the
added local intersection function has integral zero; in particular it is
unchanged when that function vanishes almost everywhere.
\end{theorem}

\begin{proof}
Chen--Moriwaki define the global intersection number as the integral of the
integrable local intersection function.  The measure space of
$S\boxplus T$ is a disjoint union, so its integral is the sum of the two
integrals displayed above.
\end{proof}

\subsection{Absolute values, zero seminorms, and external absence}

Let
\[
\sigma:\R_{\geq0}\longrightarrow\B,
\qquad
\sigma(x)=0\iff x=0.
\]

\begin{theorem}[Place-independent scalar Booleanization]
\label{thm:absolute-value-boolean-no-go}
For every field $K$ and every absolute value $|\cdot|$ on $K$,
\[
\sigma(|a|)=
\begin{cases}
0,&a=0,\\
1,&a\neq0.
\end{cases}
\]
Hence scalar Booleanization is the same at every place.  A trivial absolute value therefore lies in the nondegenerate Boolean state, whereas a zero seminorm lies in the degenerate state.
\end{theorem}

\begin{proof}
By definition of an absolute value, $|a|>0$ for every nonzero field element.
The formula follows immediately.  A zero seminorm vanishes on nonzero
vectors and therefore fails the positivity axiom for a norm or absolute
value.
\end{proof}

\begin{proposition}[Two Boolean shadows of the three local states]
\label{prop:three-metric-support-states}
Let $\ell$ be a one-dimensional real vector space.  Modulo positive
rescaling, the local line data with an adjoined absence state form the
three-element chain
\[
C_2=\{0<1<2\},
\]
where
\[
0=\text{external absence},
\qquad
1=\text{zero seminorm},
\qquad
2=\text{nondegenerate seminorm}.
\]
Equip $C_2$ with join and meet.  The two threshold characters
\[
\theta_{\mathrm{pres}}(j)=\mathbf1_{\{j\ge1\}},
\qquad
\theta_{\mathrm{ndeg}}(j)=\mathbf1_{\{j\ge2\}}
\]
are semiring homomorphisms $C_2\to\B$, and their product
\[
\Theta_{\mathrm{met}}:C_2\longrightarrow\B^2,
\qquad
j\longmapsto
(\theta_{\mathrm{pres}}(j),\theta_{\mathrm{ndeg}}(j)),
\]
is injective with image
\[
\{(s,d)\in\B^2:d\leq s\}
=\{(0,0),(1,0),(1,1)\}.
\]
Tensor product of local line data corresponds to meet in $C_2$.
The first Boolean coordinate is split presence; the second is the
Connes--Consani nondegeneracy coordinate.  Together they reconstruct the
full three-state local datum.
\end{proposition}

\begin{proof}
Every seminorm under consideration is either zero or $\lambda|\cdot|$ with
$\lambda>0$, and positive rescaling identifies all nonzero values of
$\lambda$.  Tensoring with an absent datum remains absent, tensoring a zero
seminorm with a present datum gives a zero seminorm, and tensoring two
nondegenerate seminorms gives a nondegenerate seminorm; this is the meet
table of $C_2$.  The two maps are the threshold characters from
Corollary~\ref{cor:chain-threshold-reconstruction}.  Their values are
\[
0\mapsto(0,0),\qquad
1\mapsto(1,0),\qquad
2\mapsto(1,1),
\]
which proves injectivity and the image formula.  The interpretation of the
second coordinate follows from the two Connes--Consani archimedean metric
orbits.
\end{proof}

\begin{theorem}[Measurable three-state metric lattice]
\label{thm:global-three-state-metric-lattice}
Let $C_2=\{0<1<2\}$ and let
\[
\mathscr C_S
:=L^0(\Omega,C_2)/{\sim_\nu}
\]
be the lattice of measurable $C_2$-valued functions modulo almost-everywhere
equality, with pointwise join and meet.  For $c\in\mathscr C_S$, set
\[
P_c=[\{\omega:c(\omega)\geq1\}],
\qquad
N_c=[\{\omega:c(\omega)=2\}].
\]
Then
\[
\boxed{
\mathscr C_S
\xrightarrow{\ \sim\ }
\operatorname{Inc}(\mathscr B_S)
:=\{(P,N)\in\mathscr B_S^2:N\leq P\},
\qquad
c\longmapsto(P_c,N_c),
}
\]
is an isomorphism of bounded distributive lattices.  The two coordinate
maps are the global presence and nondegeneracy characters.  Under tensor
product of local line states, the corresponding incidence pairs multiply
by coordinatewise meet.

The map
\[
\Delta_S^{\mathrm{met}}:
\operatorname{Inc}(\mathscr B_S)
\longrightarrow\Hom_{\mathrm{gp}}(K^\times,\R),
\qquad
\Delta_S^{\mathrm{met}}(P,N)=\Delta_S(N),
\]
is a Boolean-lattice valuation and is natural under measure-preserving
refinements.  It records the product-formula defect of the nondegenerate
metric locus while retaining the independent presence locus.
\end{theorem}

\begin{proof}
The inclusion $N_c\leq P_c$ is immediate.  Conversely, an incident pair
$N\leq P$ determines the measurable function which is $2$ on $N$, $1$
on $P\setminus N$, and $0$ outside $P$.  These constructions are inverse
modulo null sets and preserve pointwise joins and meets.  The tensor-product
table is the meet table of $C_2$, hence becomes coordinatewise intersection.
Finally, projection $(P,N)\mapsto N$ preserves joins and meets, so the
valuation and naturality statements follow from
Theorems~\ref{thm:adelic-support-defect} and
\ref{thm:adelic-defect-naturality}.
\end{proof}

\begin{remark}[Faithful metric-state coordinates]
The three local structures are connected by two independent Boolean
projections:
\[
C_2\xrightarrow{\theta_{\mathrm{pres}}}\B,
\qquad
C_2\xrightarrow{\theta_{\mathrm{ndeg}}}\B.
\]
Their product is faithful.  This is the local metric instance of the
multi-Boolean normal form of Theorem~\ref{thm:birkhoff-support-coordinates}.
\end{remark}


\section{Rigidified denominator points and universal Frobenius}
\label{sec:rigidified-denominators}

Finite denominator refinement is invisible on isomorphism classes of
arithmetic-site points.  It becomes visible after retaining the embedding
of the arithmetic lattice.

\begin{definition}[Rigidified arithmetic-site point]
\label{def:rigidified-point}
A rigidified point is a pair $(H,\iota)$ consisting of an ordered rank-one
subgroup $H\subseteq\Q$ and an order-preserving embedding
\[
\iota:\Z\hookrightarrow H.
\]
A morphism
\[
f:(H,\iota_H)\longrightarrow(K,\iota_K)
\]
is an injective order-preserving homomorphism satisfying
$f\circ\iota_H=\iota_K$.

For $Q\geq1$, put
\[
H_Q=\frac1Q\Z,
\qquad
\mathfrak x_Q=(H_Q,\iota_Q),
\qquad
\iota_Q(n)=n.
\]
\end{definition}

\begin{theorem}[Divisibility classification]
\label{thm:rigidified-divisibility}
For positive integers $Q,M$,
\[
\operatorname{Hom}(\mathfrak x_Q,\mathfrak x_M)
=
\begin{cases}
\{\text{one morphism}\},&Q\mid M,\\
\varnothing,&Q\nmid M.
\end{cases}
\]
Consequently,
\[
\mathfrak x_Q\cong\mathfrak x_M
\quad\Longleftrightarrow\quad
Q=M.
\]
\end{theorem}

\begin{proof}
Let $f:\mathfrak x_Q\to\mathfrak x_M$ be a morphism.  Rigidification gives
$f(1)=1$.  For every $a\in\Z$,
\[
Qf(a/Q)=f(a)=a.
\]
Since $H_M\subseteq\Q$ is torsion free, $f(a/Q)=a/Q$.  Hence $f$ is forced
to be the inclusion $H_Q\hookrightarrow H_M$.  Such an inclusion exists
if and only if $1/Q\in(1/M)\Z$, equivalently $Q\mid M$.  The isomorphism
criterion follows by applying this condition in both directions.
\end{proof}

\begin{corollary}[The framed fibre over the arithmetic point $\Z$]
\label{cor:framed-fibre}
After forgetting the rigidification, every $\mathfrak x_Q$ represents the
same arithmetic-site point $[\Z]$.  The full divisibility category
$(\N^\times,\mid)$ is therefore recovered as a category of distinct
framings of one point.
\end{corollary}

\begin{proof}
Multiplication by $Q$ gives an order isomorphism
$H_Q\cong\Z$.  Theorem \ref{thm:rigidified-divisibility} shows that the
framings remain distinct and that their morphisms are exactly divisibility.
\end{proof}

Connes--Consani also realize the full adele class space by framed arithmetic
divisors $(L,\xi,\tau_\infty)$ \cite{CCJacobian}.  The denominator group
has a canonical such frame:
\[
\mathfrak D_Q
:=\left(\frac1Q\Z,\ \xi_Q(x)=Qx\in\widehat\Z,\
\tau_Q(x)=x\in\R\right).
\]

\begin{theorem}[The two denominator laws]
\label{thm:two-denominator-laws}
For positive integers $Q,M$:
\begin{enumerate}[label=\textup{(\roman*)}]
\item tensor product of framed arithmetic divisors gives
\[
\mathfrak D_Q\otimes\mathfrak D_M\cong\mathfrak D_{QM};
\]
\item least common refinement inside $\Q$ gives
\[
H_Q+H_M=H_{\operatorname{lcm}(Q,M)}.
\]
Thus multiplication of denominator indices is Picard multiplication,
whereas least common multiple is join in the embedded-root refinement
poset.
\end{enumerate}
\end{theorem}

\begin{proof}
Multiplication in $\Q$ induces
\[
H_Q\otimes_\Z H_M\xrightarrow{\sim}H_{QM},
\qquad
\frac aQ\otimes\frac bM\longmapsto\frac{ab}{QM}.
\]
Under this map the finite frames multiply as
$(Qa/Q)(Mb/M)=ab=QM(ab/QM)$, and the infinite frames multiply as ordinary
real numbers.  This proves (i).

Put $L=\operatorname{lcm}(Q,M)$.  Both $H_Q$ and $H_M$ lie in $H_L$.
Writing $d_Q=L/Q$ and $d_M=L/M$, one has
$\gcd(d_Q,d_M)=1$; choose $a,b\in\Z$ with $ad_Q+bd_M=1$.  Then
\[
\frac1L=\frac aQ+\frac bM,
\]
so the subgroup generated by $H_Q$ and $H_M$ contains $H_L$.  This proves
(ii).
\end{proof}

\begin{remark}[The two denominator monoidal laws]
The operation $Q\mapsto\Lambda_Q$ adjoins all roots of order at most $Q$
and is therefore a join/perfection operation.  It is not iteration of the
adelic tensor product, which would multiply denominator indices.  Keeping
these laws separate is necessary for the arithmetic-site interpretation of
the Beurling tower.
\end{remark}

For each $Q$, let
\[
A_Q=\Fone[T^{H_Q^+}]
=\Fone[T^{\frac1Q\N}].
\]
There is an abstract isomorphism
\[
\phi_Q:\Fone[T]\xrightarrow{\sim}A_Q,
\qquad
\phi_Q(T^n)=T^{n/Q}.
\]
If $Q\mid M$, inclusion of exponent monoids induces
$j_{Q,M}:A_Q\hookrightarrow A_M$.

\begin{theorem}[Denominator refinement is universal Frobenius]
\label{thm:denominator-frobenius}
If $M=mQ$, then
\[
\boxed{
\phi_M^{-1}\circ j_{Q,M}\circ\phi_Q
=\operatorname{Frob}_m
}
\]
on $\Fone[T]$, where
\[
\operatorname{Frob}_m(T^n)=T^{mn}.
\]
\end{theorem}

\begin{proof}
For $n\in\N$,
\[
T^n\xmapsto{\phi_Q}T^{n/Q}
\xmapsto{j_{Q,M}}T^{n/Q}
=T^{mn/M}
\xmapsto{\phi_M^{-1}}T^{mn}.
\]
This is precisely the $m$-th universal Frobenius.
\end{proof}

\begin{remark}
The finite denominator stages are isomorphic as unframed stalks.  Their
arithmetic content lies in the transition maps, which become Frobenius
endomorphisms after choosing the canonical coordinates $\phi_Q$.  This is
why denominator refinement should be regarded as motion in a framed fibre,
not as motion among ordinary point isomorphism classes.
\end{remark}

\section{LCM perfection, Steinitz points, and emergent places}
\label{sec:lcm-perfection}

\begin{definition}[Frobenius-perfect rigidified point]
\label{def:frobenius-perfect-point}
Let $\Z\subseteq H\subseteq\Q$ be a rigidified arithmetic-site point and
put
\[
A_H=\Fone[T^{H_+}].
\]
For a positive integer $m$, the point is called \emph{$m$-perfect} if
$\Frob_m:A_H\to A_H$ is an automorphism.  It is called \emph{globally
Frobenius-perfect} if it is $m$-perfect for every $m\geq1$.
\end{definition}

\begin{theorem}[Characterization and uniqueness of Frobenius perfection]
\label{thm:unique-frobenius-perfect-point}
Let $\Z\subseteq H\subseteq\Q$.
\begin{enumerate}[label=\textup{(\roman*)}]
\item $\Frob_m$ is an automorphism of $A_H$ if and only if $H$ is
$m$-divisible.
\item The smallest $p$-perfect rigidified point containing $\Z$ is
$\Z[1/p]$.
\item The unique globally Frobenius-perfect rigidified point is $H=\Q$.
\end{enumerate}
\end{theorem}

\begin{proof}
The Frobenius acts on monomials by
\[
T^h\longmapsto T^{mh}.
\]
Multiplication by $m$ is injective on the torsion-free group $H$.
Surjectivity on monomials is equivalent to the condition that for every
$h\in H_+$ there is $h'\in H_+$ with $mh'=h$, which is exactly
$m$-divisibility of $H$.  This proves (i).

The smallest subgroup of $\Q$ containing $\Z$ and closed under division
by $p$ is
\[
\bigcup_{r\geq0}p^{-r}\Z=\Z[1/p],
\]
which proves (ii).  If $H$ is globally Frobenius-perfect, then
$1/m\in H$ for every $m\geq1$, hence $\Q\subseteq H$.  Since
$H\subseteq\Q$, equality follows.  Conversely, $\Q$ is divisible by every
positive integer, proving (iii).
\end{proof}

\begin{remark}[Frobenius perfection of the exponent monoid]
Theorem \ref{thm:unique-frobenius-perfect-point} concerns bijectivity of
universal Frobenius on an exponent monoid.  It does not provide a complete
nonarchimedean Tate ring, a tilt, or any of the analytic structures required
in perfectoid geometry.
\end{remark}

Let $Q_1,\dots,Q_r$ be positive integers and put
$L=\operatorname{lcm}(Q_1,\dots,Q_r)$.

\begin{theorem}[LCM as categorical join]
\label{thm:lcm-join}
Inside $\Q$ one has
\[
H_L=H_{Q_1}+\cdots+H_{Q_r}.
\]
Equivalently, $\mathfrak x_L$ is the join of the rigidified points
$\mathfrak x_{Q_i}$, and $A_L$ is the least cyclic-root stalk containing
all $A_{Q_i}$.
\end{theorem}

\begin{proof}
Because $Q_i\mid L$, one has $H_{Q_i}\subseteq H_L$.  Put $d_i=L/Q_i$.
The least-common-multiple property implies
$\gcd(d_1,\dots,d_r)=1$.  Hence integers $u_i$ exist with
$\sum_i u_id_i=1$.  Dividing by $L$ gives
\[
\frac1L=\sum_i\frac{u_i}{Q_i},
\]
so the subgroup generated by the $H_{Q_i}$ contains the generator $1/L$
of $H_L$.  This proves equality.  Passing to positive exponent monoids and
then to spherical monoid algebras proves the stalk statement.
\end{proof}

Define
\[
\Lambda_Q=\operatorname{lcm}(1,2,\dots,Q).
\]
Choose $Q_0\geq3$ and recursively set
\[
Q_{j+1}=\Lambda_{Q_j}.
\]

\begin{lemma}[Strict growth]
\label{lem:lcm-growth}
For $Q\geq3$,
\[
\Lambda_Q\geq Q(Q-1)>Q.
\]
\end{lemma}

\begin{proof}
The coprime integers $Q$ and $Q-1$ both occur in the defining least common
multiple, so
\[
\Lambda_Q\geq\operatorname{lcm}(Q,Q-1)=Q(Q-1).
\]
\end{proof}

\begin{theorem}[Global Frobenius perfection]
\label{thm:global-perfection}
The filtered colimits of the LCM tower are
\[
\varinjlim_jH_{Q_j}=\Q,
\qquad
\varinjlim_jA_{Q_j}=\Fone[T^{\Q_+}].
\]
Every universal Frobenius
\[
\operatorname{Frob}_m:\Fone[T^{\Q_+}]
\longrightarrow\Fone[T^{\Q_+}]
\]
is an automorphism.
\end{theorem}

\begin{proof}
The union of the groups $H_{Q_j}$ is contained in $\Q$.  Conversely, take
$a/b\in\Q$.  By Lemma \ref{lem:lcm-growth}, choose $j$ with $Q_j\geq b$.
Then $b\mid Q_{j+1}$, hence $a/b\in H_{Q_{j+1}}$.  This proves the first
identity.  The second follows by taking monoid algebras of the increasing
positive cones.  Multiplication by $m$ is a bijection of $\Q_+$ with
inverse $q\mapsto q/m$, so the induced Frobenius is invertible.
\end{proof}

A supernatural or Steinitz number is a formal product
\[
N=\prod_p p^{n_p},
\qquad
n_p\in\N\cup\{\infty\}.
\]
It determines the ordered subgroup
\[
H_N=
\{x\in\Q:v_p(x)\geq-n_p\text{ for every }p\}.
\]
This is the modern point-theoretic form of Steinitz's $G$-numbers
\cite{Steinitz}.

\begin{corollary}[The maximal supernatural endpoint]
\label{cor:maximal-supernatural}
For
\[
\Omega=\prod_p p^\infty,
\]
one has $H_\Omega=\Q$.  Thus the LCM tower converges to the maximal
supernatural point $x_\Omega$ of the arithmetic site, whose stalk is the
globally Frobenius-perfect algebra $\Fone[T^{\Q_+}]$.
\end{corollary}

\begin{theorem}[Places emerge only at perfection]
\label{thm:places-emerge}
Let $\Sigma_H$ be the Connes--Consani set of places at which $H$ embeds
densely.  Then
\[
\Sigma_{H_{Q_j}}=\varnothing
\quad\text{for every finite stage }j,
\]
whereas
\[
\Sigma_{\Q}=\{\infty\}\cup\{p:p\text{ prime}\}.
\]
In particular,
\[
\bigcup_j\Sigma_{H_{Q_j}}
\neq
\Sigma_{\varinjlim_jH_{Q_j}}.
\]
\end{theorem}

\begin{proof}
Each $H_{Q_j}$ is order-isomorphic to $\Z$, and Connes--Consani prove that
$\Sigma_H$ is empty for $H\cong\Z$.  The group $\Q$ is dense in $\R$ and
is $p$-divisible for every prime $p$, so the classification of places in
\cite{CCAbsolute} gives the second identity.
\end{proof}

\begin{remark}
Theorem \ref{thm:places-emerge} is not a formal paradox.  The functor
$H\mapsto\Sigma_H$ detects density and divisibility properties that need
not occur at any finite cyclic-root stage.  Coherent Frobenius perfection
creates local directions that are absent at every finite level.
\end{remark}

\section{Cyclotomic carry, finite-adelic duality, and the perfected denominator anomaly}
\label{sec:perfected-anomaly}

The globally Frobenius-perfect point carries two distinct kinds of data.
The first is the relative denominator quotient produced by the fixed
rigidification $\Z\subseteq\Q$.  The second is the inverse-limit character
data retained by the transition maps of the denominator tower.  They are
related by two canonical nonsplit extensions.  This section computes both
extensions, identifies the resulting profinite symmetry, and compares the
answer with the rooted arithmetic divisors of Connes--Consani.

The word \emph{perfect} in this section always means Frobenius perfection
of an arithmetic-site stalk.  It does not mean that
$\Fone[T^{\Q_+}]$ is a perfectoid Tate ring or that $x_\Omega$ is a
perfectoid space in the sense of Scholze \cite{ScholzePerfectoid}.

\subsection{Finite denominator phases and their extension classes}

\begin{definition}[Relative phase quotient]
\label{def:finite-phase-quotient}
For $Q\geq1$, define
\[
\Phi_Q:=H_Q/\Z=\frac{Q^{-1}\Z}{\Z}.
\]
The map
\[
\eta_Q:\Z/Q\Z\xrightarrow{\sim}\Phi_Q,
\qquad
[a]_Q\longmapsto \frac aQ+\Z,
\]
is a canonical isomorphism.  We write
\[
\mathcal E_Q:
0\longrightarrow\Z\longrightarrow H_Q
\longrightarrow\Phi_Q\longrightarrow0
\]
for the resulting short exact sequence.
\end{definition}

\begin{theorem}[Finite carry class]
\label{thm:finite-carry-class}
Under the identification
\[
\Ext^1_{\Z}(\Z/Q\Z,\Z)
\xrightarrow{\ \sim\ }
\operatorname{coker}(Q:\Z\to\Z)
=\Z/Q\Z
\]
induced by the projective resolution
$0\to\Z\xrightarrow{Q}\Z\to\Z/Q\Z\to0$,
the extension class of $\mathcal E_Q$ is
\[
[\mathcal E_Q]=1\pmod Q.
\]
In particular, $\mathcal E_Q$ is nonsplit for $Q>1$.

If $Q\mid M$, the canonical inclusion
$\iota_{Q,M}:\Phi_Q\hookrightarrow\Phi_M$ satisfies
\[
\iota_{Q,M}^{*}[\mathcal E_M]=[\mathcal E_Q].
\]
Under the displayed Ext-identifications, the induced map is reduction
$\Z/M\Z\to\Z/Q\Z$.
\end{theorem}

\begin{proof}
The projective resolution
\[
0\longrightarrow\Z\xrightarrow{\,Q\,}\Z
\longrightarrow\Z/Q\Z\longrightarrow0
\]
gives
\[
\Ext^1_{\Z}(\Z/Q\Z,\Z)
\cong\operatorname{coker}(Q:\Z\to\Z)
\cong\Z/Q\Z.
\]
Multiplication by $Q$ gives an isomorphism $H_Q\xrightarrow{\sim}\Z$.
Under this isomorphism, the inclusion $\Z\hookrightarrow H_Q$ becomes
multiplication by $Q$.  Thus $\mathcal E_Q$ is the pushout of the displayed
resolution along the identity of its kernel, and therefore represents the
class $1\pmod Q$.

Suppose $M=mQ$.  The inclusion $H_Q\hookrightarrow H_M$ induces
\[
\iota_{Q,M}:\Phi_Q\longrightarrow\Phi_M,
\qquad
\frac aQ+\Z\longmapsto\frac aQ+\Z.
\]
The pullback of $\mathcal E_M$ along this map consists exactly of those
$h\in H_M$ whose class modulo $\Z$ lies in the image of $\Phi_Q$; this
subgroup is $H_Q$.  Hence the pullback extension is $\mathcal E_Q$.
The corresponding contravariant map on Ext sends $1\pmod M$ to
$1\pmod Q$, and is therefore the ordinary reduction map.
\end{proof}

\begin{corollary}[Finite sign and phase shadows]
\label{cor:finite-phase-shadows}
The case $Q=2$ is the binary sign extension, the case $Q=3$ is the
ternary phase extension, and general $Q$ gives a canonical $Q$-phase
extension.  None of these extensions admits a homomorphic choice of a
representative phase.
\end{corollary}

\begin{proof}
Only the terminology is new.  Nonsplitting is Theorem
\ref{thm:finite-carry-class}.
\end{proof}

\begin{definition}[Finite carry cocycle]
\label{def:finite-carry-cocycle}
For $[a]_Q\in\Z/Q\Z$, let $0\leq a<Q$ be its standard representative and
put
\[
s_Q([a]_Q)=\frac aQ\in H_Q.
\]
Define
\[
c_Q([a]_Q,[b]_Q)
:=s_Q([a]_Q)+s_Q([b]_Q)-s_Q([a+b]_Q)
=\left\lfloor\frac{a+b}{Q}\right\rfloor.
\]
\end{definition}

\begin{proposition}[Twisted reconstruction of the denominator group]
\label{prop:finite-twisted-reconstruction}
The function $c_Q$ is a normalized symmetric $2$-cocycle with values in
$\Z$.  On the set $\Z\times\Z/Q\Z$, define
\[
(n,\bar a)+(m,\bar b)
=
\bigl(n+m+c_Q(\bar a,\bar b),\bar a+\bar b\bigr).
\]
Then
\[
\Z\times_{c_Q}\Z/Q\Z\xrightarrow{\sim}H_Q,
\qquad
(n,[a]_Q)\longmapsto n+\frac aQ,
\]
is an isomorphism of abelian groups.  The cocycle $c_Q$ represents
$[\mathcal E_Q]$.
\end{proposition}

\begin{proof}
Normalization and symmetry are immediate.  Associativity of the displayed
operation is equivalent to
\[
c_Q(x,y)+c_Q(x+y,z)=c_Q(y,z)+c_Q(x,y+z),
\]
which follows by expanding each side using the section $s_Q$ and
associativity in $H_Q$.  Euclidean division gives a unique expression
$h=n+a/Q$ with $0\leq a<Q$, so the displayed map is bijective.  Its
compatibility with addition is exactly the definition of $c_Q$.
\end{proof}

\subsection{The global carry class}

The filtered union of the finite phase groups is
\[
\Phi_\Omega:=\varinjlim_Q\Phi_Q=\Q/\Z.
\]
The rigidification $\Z\subseteq\Q$ gives the global extension
\[
\mathcal E_\Omega:
0\longrightarrow\Z\longrightarrow\Q
\longrightarrow\Q/\Z\longrightarrow0.
\]

\begin{theorem}[The global carry class is the profinite unit]
\label{thm:global-carry-class}
There is a canonical isomorphism
\[
\Ext^1_{\Z}(\Q/\Z,\Z)
\xrightarrow{\sim}
\varprojlim_Q\Ext^1_{\Z}(\Phi_Q,\Z)
\xrightarrow{\sim}
\varprojlim_Q\Z/Q\Z
=\widehat\Z.
\]
Under this isomorphism,
\[
\boxed{
[\mathcal E_\Omega]
=(1\bmod Q)_Q
=1_{\widehat\Z}.}
\]
Consequently, $\mathcal E_\Omega$ is nonsplit, and every finite carry
class is its pullback along $\Phi_Q\hookrightarrow\Q/\Z$.
\end{theorem}

\begin{proof}
The system $\{\Phi_Q\}$ is filtered by divisibility and has colimit
$\Q/\Z$.  The standard derived-limit exact sequence for Ext \cite{Weibel} gives
\[
0\longrightarrow
\varprojlim
olimits^{1}_Q\Hom(\Phi_Q,\Z)
\longrightarrow
\Ext^1_{\Z}(\varinjlim_Q\Phi_Q,\Z)
\longrightarrow
\varprojlim_Q\Ext^1_{\Z}(\Phi_Q,\Z)
\longrightarrow0.
\]
Every homomorphism from the finite group $\Phi_Q$ to the torsion-free group
$\Z$ is zero.  Hence the $\varprojlim^1$ term vanishes and the first arrow
displayed in the theorem is an isomorphism.  The second isomorphism is
Theorem \ref{thm:finite-carry-class}; its transition maps are reductions.
The pullback of $\mathcal E_\Omega$ to $\Phi_Q$ is $\mathcal E_Q$, so its
image is the compatible family $(1\bmod Q)_Q$, which is the multiplicative
unit of $\widehat\Z$.
\end{proof}

\begin{definition}[Global carry cocycle]
\label{def:global-carry-cocycle}
Let
\[
s:\Q/\Z\longrightarrow\Q\cap[0,1)
\]
be the standard set-theoretic section.  Define
\[
c_\Omega(x,y)=s(x)+s(y)-s(x+y)\in\Z.
\]
\end{definition}

\begin{proposition}[Explicit global history law]
\label{prop:global-carry-cocycle}
The function $c_\Omega$ is a normalized symmetric $2$-cocycle and
\[
\Q\cong\Z\times_{c_\Omega}(\Q/\Z)
\]
by the same twisted addition as in Proposition
\ref{prop:finite-twisted-reconstruction}.  Its restriction to $\Phi_Q$ is
$c_Q$.  Thus $c_\Omega$ simultaneously records every finite carry law.
\end{proposition}

\begin{proof}
The proof of Proposition \ref{prop:finite-twisted-reconstruction} applies
verbatim.  If $x=[a/Q]$ and $y=[b/Q]$ with $0\leq a,b<Q$, then
\[
c_\Omega(x,y)=\left\lfloor\frac{a+b}{Q}\right\rfloor=c_Q([a]_Q,[b]_Q).
\]
\end{proof}

\begin{corollary}[Classifying-space form]
\label{cor:carry-gerbe}
The extension $\mathcal E_\Omega$ induces a homotopy fibre sequence of
classifying spaces
\[
B\Z\longrightarrow B\Q\longrightarrow B(\Q/\Z).
\]
Equivalently, it determines a principal $B\Z$-bundle over $B(\Q/\Z)$,
classified by the central-extension class represented by $c_\Omega$ in
$H^2(B(\Q/\Z);\Z)$.  Its restriction to $B(\Z/Q\Z)$ is the generator
represented by $c_Q$.
\end{corollary}

\begin{proof}
A central extension of discrete groups gives a fibration of classifying
spaces with fibre the classifying space of the kernel.  The normalized
cocycle is the standard group-cohomological representative of that
extension.
\end{proof}

\subsection{Finite-adelic character duality}

For a discrete abelian group $H$, put
\[
H^\vee:=\Hom(H,\Q/\Z).
\]
Connes--Consani identify the character dual of $\Q$ with the additive group
of finite adeles \cite{CCJacobian}:
\[
\psi:\mathbb A_f\xrightarrow{\sim}\Hom(\Q,\Q/\Z),
\qquad
\psi(a)(q)=\pi(qa),
\]
where $\pi:\mathbb A_f\to\mathbb A_f/\widehat\Z\cong\Q/\Z$ is the
canonical quotient.

\begin{theorem}[The inverse limit of finite character systems is adelic]
\label{thm:adelic-history-limit}
For every $Q$, evaluation at $1/Q$ gives an isomorphism
\[
\operatorname{ev}_Q:H_Q^\vee\xrightarrow{\sim}\Q/\Z.
\]
If $Q\mid M=mQ$, the restriction map
\[
H_M^\vee\longrightarrow H_Q^\vee
\]
corresponds under these evaluations to multiplication by $m$ on
$\Q/\Z$.  Consequently,
\[
\boxed{
\varprojlim_Q H_Q^\vee
\cong
\Hom\!\left(\varinjlim_QH_Q,\Q/\Z\right)
\cong
\Hom(\Q,\Q/\Z)
\cong
\mathbb A_f.}
\]
A finite adele is therefore exactly a coherent family of denominator-phase
histories.
\end{theorem}

\begin{proof}
The group $H_Q$ is freely generated by $1/Q$, so evaluation at that
generator is an isomorphism.  If $M=mQ$ and $\varphi\in H_M^\vee$, then
\[
\varphi(1/Q)=\varphi(m/M)=m\varphi(1/M),
\]
which proves the transition formula.  Since $\Hom(-,\Q/\Z)$ sends
filtered colimits to inverse limits,
\[
\varprojlim_QH_Q^\vee
\cong\Hom(\varinjlim_QH_Q,\Q/\Z).
\]
Theorem \ref{thm:global-perfection} identifies the colimit with $\Q$, and
the final isomorphism is the Connes--Consani adelic duality theorem.
\end{proof}

\begin{theorem}[The finite-adelic dual extension]
\label{thm:adelic-history-extension}
Applying $\Hom(-,\Q/\Z)$ to $\mathcal E_\Omega$ gives the canonical exact
sequence
\[
\boxed{
0\longrightarrow\widehat\Z
\longrightarrow\mathbb A_f
\xrightarrow{\ \pi\ }\Q/\Z
\longrightarrow0.}
\]
Here
\[
\widehat\Z=\End(\Q/\Z)
\]
is the subgroup of histories whose restriction to the rigidifying copy of
$\Z$ is zero, and $\pi(a)=\psi(a)(1)$.  This extension is nonsplit, even as
an abstract extension of abelian groups.
\end{theorem}

\begin{proof}
The group $\Q/\Z$ is divisible and therefore injective in the category of
abelian groups.  Applying $\Hom(-,\Q/\Z)$ to
\[
0\to\Z\to\Q\to\Q/\Z\to0
\]
therefore gives
\[
0\to\Hom(\Q/\Z,\Q/\Z)
\to\Hom(\Q,\Q/\Z)
\to\Hom(\Z,\Q/\Z)
\to0.
\]
The three nonzero terms identify respectively with $\widehat\Z$,
$\mathbb A_f$, and $\Q/\Z$.  Under the Connes--Consani isomorphism, the
last arrow is evaluation at $1$ and is the canonical quotient $\pi$.

The additive group of $\mathbb A_f$ is torsion free.  A splitting would
embed the nonzero torsion group $\Q/\Z$ into $\mathbb A_f$, which is
impossible.
\end{proof}

\begin{figure}[t]
\centering
\begin{tikzcd}[column sep=large,row sep=large]
0 \arrow[r] & \Z \arrow[r] \arrow[d,equal]
& H_Q \arrow[r] \arrow[d,hook]
& \Phi_Q \arrow[r] \arrow[d,hook] & 0 \\
0 \arrow[r] & \Z \arrow[r]
& \Q \arrow[r]
& \Q/\Z \arrow[r] & 0
\end{tikzcd}
\qquad
\begin{tikzcd}[column sep=large,row sep=large]
0 \arrow[r] & \widehat\Z \arrow[r]
& \mathbb A_f \arrow[r,"\pi"]
& \Q/\Z \arrow[r] & 0
\end{tikzcd}
\caption{The finite carry extension embeds in the global carry extension;
character duality converts the global extension into the nonsplit adelic
history extension.}
\label{fig:carry-history-extensions}
\end{figure}

\begin{theorem}[Torsion transmutation at perfection]
\label{thm:torsion-transmutation}
For every finite denominator stage,
\[
\Tor(H_Q^\vee)=H_Q^\vee\cong\Q/\Z.
\]
At the globally perfect point,
\[
\Q^\vee\cong\mathbb A_f,
\qquad
\Tor(\Q^\vee)=0.
\]
Consequently,
\[
\boxed{
\Tor\!\left(\varprojlim_QH_Q^\vee\right)=0,
\qquad
\varprojlim_Q\Tor(H_Q^\vee)\cong\mathbb A_f.}
\]
Thus the torsion functor does not commute with the inverse limit of the
denominator tower: finite cyclotomic phase data is converted into
nontorsion finite-adelic character data.
\end{theorem}

\begin{proof}
Each $H_Q$ is infinite cyclic, hence $H_Q^\vee\cong\Q/\Z$ and is entirely
torsion.  Theorem \ref{thm:adelic-history-limit} identifies the inverse
limit with $\mathbb A_f$, whose additive group is torsion free.
\end{proof}

\subsection{Relative phases and rooted phases}

For a set $S$ of rational primes, put
\[
H_S:=\Z[S^{-1}]
=\Z[p^{-1}:p\in S]\subseteq\Q.
\]
Let $C_{p^\infty}=\Q_p/\Z_p$ denote the Pruefer $p$-group.

\begin{theorem}[Phase complementarity for saturated denominator points]
\label{thm:phase-complementarity}
The relative denominator quotient and the torsion of the Connes--Consani
character dual are
\[
H_S/\Z\cong\bigoplus_{p\in S}C_{p^\infty},
\qquad
\Tor(H_S^\vee)
\cong\bigoplus_{p\notin S}C_{p^\infty}.
\]
Hence the canonical primary decomposition of $\Q/\Z$ gives
\[
\boxed{
H_S/\Z\ \oplus\ \Tor(H_S^\vee)
\cong\Q/\Z.}
\]
At $S=\varnothing$, all cyclotomic phase is rooted torsion and no relative
phase is present.  At $S$ equal to the set of all primes, all cyclotomic
phase is relative denominator phase and the rooted torsion vanishes.
\end{theorem}

\begin{proof}
The $p$-primary subgroup of $H_S/\Z$ is zero when $p\notin S$ and is the
union of the cyclic groups $p^{-n}\Z/\Z$ when $p\in S$, hence is
$C_{p^\infty}$.  This proves the first formula.  The singular primes of
$H_S$ in the sense of Connes--Consani are exactly the primes in $S$.
Their character-dual theorem identifies the torsion subgroup with the
roots of unity whose order is prime to every singular prime, namely
$\bigoplus_{p\notin S}C_{p^\infty}$.  The final formula is the canonical
primary decomposition
$\Q/\Z=\bigoplus_pC_{p^\infty}$.
\end{proof}

\begin{corollary}[Root collapse at global perfection]
\label{cor:root-collapse-global-perfection}
For each finite $Q$, the singular-prime set of $H_Q$ is empty and
\[
\Tor(H_Q^\vee)\cong\Q/\Z.
\]
For $H_\Omega=\Q$, every prime is singular and
\[
\Tor(H_\Omega^\vee)=0.
\]
Thus a Connes--Consani root map
\[
\rho:\Q/\Z\longrightarrow H_\Omega^\vee
\]
whose image is required to lie in the torsion subgroup is necessarily
zero.  The finite root data does not survive as torsion at the endpoint;
it survives through the finite-adelic character group of Theorem
\ref{thm:adelic-history-limit}.
\end{corollary}

\begin{proof}
Finite exponents never make $H_Q$ divisible by a prime, whereas $\Q$ is
$p$-divisible for every $p$.  Apply the Connes--Consani torsion theorem and
Theorem \ref{thm:torsion-transmutation}.
\end{proof}

\begin{remark}[The exact location of the inverse-limit data]
The static perfect point does not contain a residual group of torsion roots.
The retained datum is the inverse-limit object $\mathbb A_f$ and, more precisely,
the nonsplit extension of Theorem \ref{thm:adelic-history-extension}.
This distinction prevents the relative quotient $\Q/\Z$ from being
mistaken for the torsion subgroup of $\Q^\vee$.
\end{remark}

\subsection{Cyclotomic symmetry and the Bost--Connes interface}

\begin{theorem}[Cyclotomic symmetry of the anomaly]
\label{thm:cyclotomic-symmetry-anomaly}
There are canonical isomorphisms
\[
\Q/\Z\cong\bigoplus_pC_{p^\infty},
\qquad
\End(\Q/\Z)\cong\prod_p\Z_p=\widehat\Z,
\]
and therefore
\[
\boxed{
\Aut(\Q/\Z)\cong\widehat\Z^{\,\times}.}
\]
The exponential map $x\mapsto e^{2\pi i x}$ identifies $\Q/\Z$ with the
group $\mu_\infty$ of all complex roots of unity.  Multiplication by
$u\in\widehat\Z^{\,\times}$ acts equivariantly on the adelic history
extension:
\[
\begin{tikzcd}[column sep=large]
0 \arrow[r] & \widehat\Z \arrow[r] \arrow[d,"\times u"]
& \mathbb A_f \arrow[r,"\pi"] \arrow[d,"\times u"]
& \Q/\Z \arrow[r] \arrow[d,"\times u"] & 0 \\
0 \arrow[r] & \widehat\Z \arrow[r]
& \mathbb A_f \arrow[r,"\pi"]
& \Q/\Z \arrow[r] & 0.
\end{tikzcd}
\]
\end{theorem}

\begin{proof}
The primary decomposition of a torsion divisible group gives the first
isomorphism.  The endomorphism ring of $C_{p^\infty}$ is $\Z_p$, and
homomorphisms between distinct primary components vanish.  Hence
\[
\End(\Q/\Z)\cong\prod_p\Z_p.
\]
Its invertible elements are exactly $\prod_p\Z_p^\times$.  The
exponential identification is immediate.  Finally, $u$ is a unit of the
maximal compact subring $\widehat\Z\subset\mathbb A_f$, and
$\pi(ua)=u\pi(a)$ under the identification of $\widehat\Z$ with
$\End(\Q/\Z)$.
\end{proof}

\begin{corollary}[Bost--Connes symmetry group]
\label{cor:bc-symmetry-interface}
The automorphism group of the perfected phase quotient is the arithmetic
symmetry group of the Bost--Connes system:
\[
\Aut(\Q/\Z)
\cong\widehat\Z^{\,\times}
\cong\Gal(\Q^{\mathrm{ab}}/\Q).
\]
Under the last isomorphism, the action on $\mu_\infty$ is the usual
cyclotomic Galois action.  This identifies the symmetry group of the
finite-adelic dual extension with the symmetry acting on the low-temperature
and ground-state sectors of the Bost--Connes system
\cite{BostConnes,CCBCCyclotomy,CMR}.  The split-support contribution is the coefficient and symmetry layer; the $C^*$-dynamics is obtained by adjoining the Bost--Connes endomorphisms and crossed-product structure.
\end{corollary}

\begin{proof}
The first isomorphism is Theorem
\ref{thm:cyclotomic-symmetry-anomaly}; the second is the Kronecker--Weber
class-field-theory identification used in the Bost--Connes system.
\end{proof}

\subsection{Steinitz history and Möbius recovery}

\begin{definition}[Steinitz divisor history]
\label{def:steinitz-history}
For a supernatural number $N$, let
\[
\operatorname{Div}_{\mathrm{fin}}(N)
=\{n\in\N^\times:n\mid N\}.
\]
For $\Omega=\prod_pp^\infty$,
\[
\operatorname{Div}_{\mathrm{fin}}(\Omega)=\N^\times.
\]
We call the divisibility category
$(\N^\times,\mid)$ the finite denominator history of $x_\Omega$.
\end{definition}

\begin{proposition}[Cofinality and reconstruction of the anomaly class]
\label{prop:history-cofinality}
Every LCM tower $Q_{j+1}=\Lambda_{Q_j}$ with $Q_0\geq3$ is cofinal in
$(\N^\times,\mid)$.  Consequently,
\[
\widehat\Z
\cong\varprojlim_j\Z/Q_j\Z,
\]
and the class $[\mathcal E_\Omega]$ is recovered from the finite history as
\[
[\mathcal E_\Omega]=(1\bmod Q_j)_j.
\]
\end{proposition}

\begin{proof}
Given $n\geq1$, choose $j$ with $Q_j\geq n$.  Then $n\mid\Lambda_{Q_j}
=Q_{j+1}$.  Thus the tower is cofinal.  Inverse limits are unchanged by
restriction to a cofinal subsystem, and Theorem
\ref{thm:global-carry-class} gives the class formula.
\end{proof}

\begin{theorem}[Incidence inversion on denominator history]
\label{thm:mobius-history}
Let $A$ be an abelian group and let $f:\N^\times\to A$.  Define
\[
g(n)=\sum_{d\mid n}f(d).
\]
Then
\[
\boxed{
f(n)=\sum_{d\mid n}\mu(d)\,g(n/d),}
\]
where $\mu$ is the classical Möbius function.  Equivalently, the zeta
function of the locally finite divisibility poset has convolution inverse
\[
\mu_{\mid}(d,n)=\mu(n/d).
\]
Thus cumulative data on all predecessor denominator stages can be
recovered exactly from the finite Steinitz history.
\end{theorem}

\begin{proof}
Substitute the definition of $g$:
\[
\sum_{d\mid n}\mu(d)g(n/d)
=
\sum_{d\mid n}\mu(d)
\sum_{e\mid n/d}f(e).
\]
The coefficient of $f(e)$ is
\[
\sum_{d\mid n/e}\mu(d),
\]
which equals $1$ if $e=n$ and $0$ otherwise.
\end{proof}

\begin{remark}[Primewise self-similarity]
The exponent-vector identification
\[
\N^\times\cong\bigoplus_p\N
\]
turns denominator history into a restricted product of one-prime chains.
Dually,
\[
\Q/\Z\cong\bigoplus_pC_{p^\infty},
\qquad
\widehat\Z\cong\prod_p\Z_p.
\]
This primewise factorization is the precise self-similar structure behind
the informal language of an infinite or fractal denominator memory.
\end{remark}

\subsection{Combination with split support and place emergence}

Let $\chi_\C:G(\C)\to\B$ be the Boolean support character.  Under the
exponential identification $\Q/\Z\cong\mu_\infty\subset\C^\times$,
\[
\chi_\C(\zeta)=1
\quad(\zeta\in\mu_\infty),
\qquad
\chi_\C(e)=1,
\qquad
\chi_\C(\tau)=0.
\]
Thus Boolean support remembers neither cyclotomic phase nor the distinction
between nonzero phase and supported zero.  The missing information is not
another scalar in $G(\C)$; it is the relative extension and adelic
inverse-limit data computed above.

\begin{definition}[Perfected denominator anomaly]
\label{def:perfected-denominator-anomaly}
For a cofinal LCM tower $H_{Q_j}\to\Q$, define its perfected denominator
anomaly to be the triple
\[
\mathfrak A_\Omega
:=
\left(
[\mathcal E_\Omega],\
\mathcal E_\Omega^\vee,\
\operatorname{Emerg}_\Sigma
\right),
\]
where
\[
[\mathcal E_\Omega]\in\Ext^1_{\Z}(\Q/\Z,\Z),
\]
\[
\mathcal E_\Omega^\vee:
0\to\widehat\Z\to\mathbb A_f\to\Q/\Z\to0,
\]
and
\[
\operatorname{Emerg}_\Sigma
:=
\Sigma_{\Q}\setminus\bigcup_j\Sigma_{H_{Q_j}}.
\]
\end{definition}

\begin{theorem}[Exact computation of the perfected anomaly]
\label{thm:perfected-anomaly-computation}
The perfected denominator anomaly is
\[
\boxed{
\mathfrak A_\Omega
=
\left(
1_{\widehat\Z},\
0\to\widehat\Z\to\mathbb A_f\to\Q/\Z\to0,\
\{\infty\}\cup\{p:p\text{ prime}\}
\right).}
\]
Both extensions in this formula are nonsplit.  The first component retains
all finite carry classes; the second retains all coherent phase histories;
the third is the complete set of local directions created at Frobenius
perfection.
\end{theorem}

\begin{proof}
The first component is Theorem \ref{thm:global-carry-class}; the second is
Theorem \ref{thm:adelic-history-extension}; and the third is Theorem
\ref{thm:places-emerge}.
\end{proof}

\begin{remark}[Arithmetic obstruction and field-theoretic realization]
Definition~\ref{def:perfected-denominator-anomaly} packages the carry,
finite-adelic, and place-emergence obstructions at the perfect denominator
point.  Section~\ref{sec:gerbe-quantization} maps its finite phase quotients
to quadratic gerbe actions, and Section~\ref{sec:uhf-history} maps their
surface transgressions to the strict Heisenberg matrix tower.  The term
\emph{anomaly} therefore refers to this explicit sequence of obstruction
classes and realization functors.
\end{remark}


\part{Quantization, spectral charts, and analytic realization}

\section{Finite quadratic phases, gerbe anomalies, and theta quantization}
\label{sec:gerbe-quantization}

The perfected carry extension determines finite phase groups
\[
\Phi_n=n^{-1}\Z/\Z\cong\Z/n\Z.
\]
A gerbe anomaly requires one further datum: a homogeneous quadratic
refinement.  The rigidified generator $1/n$ supplies such a refinement at
every level.

\begin{definition}[Rigidified cyclic quadratic form]
\label{def:cyclic-quadratic-form}
For $\bar a\in\Phi_n$, represented by $0\leq a<n$, define
\[
q_n(\bar a)=
\begin{cases}
\dfrac{a^2}{2n}\pmod\Z,&n\ \text{even},\\[6pt]
\dfrac{n+1}{2}\dfrac{a^2}{n}\pmod\Z,&n\ \text{odd}.
\end{cases}
\]
\end{definition}

\begin{theorem}[Nondegenerate quadratic refinement]
\label{thm:cyclic-quadratic-refinement}
The function $q_n:\Phi_n\to\Q/\Z$ is well defined and homogeneous:
$q_n(r\bar a)=r^2q_n(\bar a)$ for every integer $r$.  Its polarization is
the perfect pairing
\[
b_n(\bar a,\bar b)
=q_n(\bar a+\bar b)-q_n(\bar a)-q_n(\bar b)
=\frac{ab}{n}\pmod\Z.
\]
\end{theorem}

\begin{proof}
If $n$ is even, replacing $a$ by $a+n$ changes $a^2/(2n)$ by
$a+n/2\in\Z$.  If $n$ is odd, the change is
\[
\frac{n+1}{2n}\bigl((a+n)^2-a^2\bigr)
=(n+1)a+\frac{n(n+1)}2\in\Z.
\]
Thus $q_n$ is well defined, and homogeneity follows from the defining
quadratic formula.  For even $n$ the polarization is $ab/n$.  For odd
$n$ it is $(n+1)ab/n$, which is congruent to $ab/n$ modulo $\Z$.  If
$b_n(a,b)=0$ for every $b$, then $a/n\in\Z$, so $a=0$ in $\Phi_n$.
\end{proof}

\begin{theorem}[Even-lattice realization of the cyclic phase form]
\label{thm:cyclic-discriminant-lattice}
For every $n\geq1$, the finite quadratic module $(\Phi_n,q_n)$ is the
discriminant form of an explicit even integral lattice $\Lambda_n$:
\[
\Lambda_n=
\begin{cases}
\langle n\rangle,& n\ \text{even},\\
-A_{n-1},& n\ \text{odd}.
\end{cases}
\]
Here $\langle n\rangle$ is the rank-one lattice with Gram matrix $(n)$,
$A_{n-1}=\{(x_1,\dots,x_n)\in\Z^n:\sum_i x_i=0\}$ has the standard
Euclidean form, and $A_0=0$.  There is an isomorphism
\[
(\Lambda_n^\vee/\Lambda_n,q_{\Lambda_n})
\cong(\Phi_n,q_n),
\qquad
q_{\Lambda_n}(x+\Lambda_n)=\frac{(x,x)}2\pmod\Z.
\]
\end{theorem}

\begin{proof}
Suppose first that $n$ is even.  Let $e$ generate $\langle n\rangle$ with
$(e,e)=n$.  The lattice is even, its dual is generated by $e/n$, and
\[
\Lambda_n^\vee/\Lambda_n\cong\Z/n\Z,
\qquad
q_{\Lambda_n}(a e/n)=\frac{a^2}{2n}\pmod\Z=q_n(\bar a).
\]

Now let $n$ be odd.  The root lattice $A_{n-1}$ is even, has determinant
$n$, and its discriminant group is cyclic generated by the class of
\[
\omega_1=
\left(\frac{n-1}{n},-\frac1n,\dots,-\frac1n\right).
\]
Since $(\omega_1,\omega_1)=(n-1)/n$, the negative lattice gives
\[
q_{-A_{n-1}}(a\omega_1)
=-\frac{a^2(n-1)}{2n}
\equiv\frac{a^2(n+1)}{2n}
=q_n(\bar a)\pmod\Z.
\]
For $n=1$ both groups are trivial.  This proves the claim in every case.
\end{proof}

\begin{theorem}[Finite gerbe anomaly associated with a denominator level]
\label{thm:finite-gerbe-anomaly}
The quadratic form $q_n$ determines a homotopy class
\[
\widetilde q_n:K(\Phi_n,2)\longrightarrow K(\Q/\Z,4).
\]
For every closed oriented four-manifold $X$, composition with
$\widetilde q_n$ and evaluation on $[X]$ define a quadratic action
\[
q_{n,X}:H^2(X;\Phi_n)\longrightarrow\Q/\Z.
\]
The finite gerbe path integral
\[
\mathcal A_n(X)
=
\sum_{\alpha\in H^2(X;\Phi_n)}
\frac{|H^0(X;\Phi_n)|}{|H^1(X;\Phi_n)|}
\exp\bigl(2\pi i q_{n,X}(\alpha)\bigr)
\]
is the closed-four-manifold value of an invertible four-dimensional TQFT.
In the Freed--Hopkins--Lurie--Teleman framework, an associated
three-dimensional boundary theory is a truncated morphism
$\mathbf1\to\mathcal A_n$.
\end{theorem}

\begin{proof}
Eilenberg--Mac Lane identify homogeneous quadratic maps $A\to B$ with
homotopy classes of maps $K(A,2)\to K(B,4)$.  Applying this to $q_n$
gives $\widetilde q_n$.  Theorem~\ref{thm:cyclic-discriminant-lattice}
realizes $q_n$ as the discriminant form of an even lattice, placing it in
the finite quadratic setting of \cite[\S5.2]{FHLT}.  Their finite
path-integral theorem, applied to the $2$-groupoid of $\Phi_n$-gerbes,
gives the displayed groupoid-cardinality weight, the invertible
four-dimensional theory, and its boundary truncated morphism.
\end{proof}

\subsection{Surface transgression and the finite Heisenberg algebra}

For the oriented torus $T^2$, identify
\[
H^1(T^2;\Phi_n)\cong\Phi_n^2.
\]
The cup product and $b_n$ induce the alternating form
\[
\Omega_n\bigl((a,b),(c,d)\bigr)
=\frac{ad-bc}{n}\pmod\Z.
\]

\begin{definition}[Finite Weyl operators]
Let $\zeta_n=e^{2\pi i/n}$ and let
\[
\mathscr H_n=\C[\Phi_n]
\]
with basis $\delta_0,\dots,\delta_{n-1}$.  Define
\[
P_n\delta_j=\zeta_n^j\delta_j,
\qquad
Q_n\delta_j=\delta_{j+1},
\]
with indices modulo $n$.  Put
\[
\omega_n((a,b),(c,d)):=\zeta_n^{-bc}.
\]
\end{definition}

\begin{theorem}[Finite Stone--von Neumann algebra]
\label{thm:finite-stone-von-neumann}
The Weyl operators satisfy
\[
P_n^n=Q_n^n=1,
\qquad
P_nQ_n=\zeta_n Q_nP_n.
\]
The $n^2$ operators
\[
\{P_n^aQ_n^b:0\leq a,b<n\}
\]
form a basis of $\End_\C(\mathscr H_n)$.  Consequently,
\[
\C_{\omega_n}[\Phi_n^2]
\cong
\End_\C(\mathscr H_n)
\cong M_n(\C),
\]
and $\mathscr H_n$ is the irreducible Schr\"odinger module.  The
commutator of the Weyl operators is
\[
W_xW_yW_x^{-1}W_y^{-1}
=\exp\bigl(2\pi i\Omega_n(x,y)\bigr).
\]
\end{theorem}

\begin{proof}
The commutation relation follows on each basis vector.  It implies
\[
(P_n^aQ_n^b)(P_n^cQ_n^d)
=\zeta_n^{-bc}P_n^{a+c}Q_n^{b+d},
\]
which is the cocycle $\omega_n$ and gives the stated commutator.  Suppose
\[
\sum_{a,b}c_{a,b}P_n^aQ_n^b=0.
\]
Applying this operator to $\delta_j$ and projecting to
$\delta_{j+b}$ shows, for every fixed $b$, that
\[
\sum_a c_{a,b}\zeta_n^{aj}=0
\qquad\text{for all }j.
\]
The finite Fourier matrix is invertible, so every $c_{a,b}=0$.  Thus the
operators are linearly independent.  Their number equals
$\dim_\C M_n(\C)=n^2$, so they form a basis.  The image algebra is all of
$\End_\C(\mathscr H_n)$, which proves irreducibility.
\end{proof}

\subsection{The explicit bundle-gerbe realization}

Let $L_N\to T^2$ be the degree-$N$ Hermitian line bundle and let
$\mathcal G_N\to T^3$ be the explicit bundle gerbe constructed in
\cite{GerbeTheta}.  Its Dixmier--Douady class and curvature are
\[
DD(\mathcal G_N)
=N[dt]\smile[dx\wedge dy],
\qquad
H_N=2\pi iN\,dt\wedge dx\wedge dy.
\]

\begin{theorem}[Winding-sector realization of the arithmetic phase]
\label{thm:winding-theta-realization}
On the winding-$r$ loop sector,
\[
\Trg(\mathcal G_N)|_{\mathcal L_r}
\cong L_N^{\otimes r}.
\]
Put $d=rN>0$.  For each flat parameter $(\phi,\psi)$, the holomorphic
theta space has dimension $d$ and carries normalized operators
$P_{r,N},Q_{r,N}$ satisfying
\[
P_{r,N}Q_{r,N}=e^{2\pi i/d}Q_{r,N}P_{r,N}.
\]
For $N=1$ and $r=n$, there is an isomorphism of Heisenberg modules
\[
\boxed{
\mathscr H_n
\cong
H^0(E_\tau,L_1^{\otimes n}\otimes F_{\phi,\psi}).
}
\]
It is unique up to a nonzero scalar.
\end{theorem}

\begin{proof}
The transgression, dimension, theta-basis, and magnetic-translation
statements are proved in \cite{GerbeTheta}.  After the scalar
normalization which removes the flat holonomies, its theta basis
$\theta_0,\dots,\theta_{n-1}$ satisfies
\[
P\theta_j=\zeta_n^j\theta_j,
\qquad
Q\theta_j=\theta_{j+1}.
\]
Thus $\delta_j\mapsto\theta_j$ intertwines the generators in
Theorem~\ref{thm:finite-stone-von-neumann}.  Uniqueness up to scalar is
Schur's lemma for the irreducible finite Heisenberg representation.
\end{proof}

\begin{remark}[Bulk-to-boundary realization]
The finite gerbe theory $\mathcal A_n$ supplies the four-dimensional
anomaly object.  Transgression to $T^2$ produces the Heisenberg boundary
algebra, and Theorem~\ref{thm:winding-theta-realization} realizes its
irreducible module as the winding-$n$ theta space of $\mathcal G_1$.
\end{remark}

\begin{theorem}[Frobenius as winding dilation]
\label{thm:frobenius-winding-dilation}
Let $\rho_m:S^1\to S^1$ be the degree-$m$ map, and let
$c_m:\mathcal L_n\to\mathcal L_{mn}$ be precomposition of loops by
$\rho_m$.  Then
\[
c_m^*\bigl(\Trg(\mathcal G_N)|_{\mathcal L_{mn}}\bigr)
\cong
\bigl(\Trg(\mathcal G_N)|_{\mathcal L_n}\bigr)^{\otimes m}.
\]
Thus multiplication by $m$ on arithmetic Frobenius degree is realized by
multiplication of loop winding.
\end{theorem}

\begin{proof}
A winding-$n$ loop composed with the degree-$m$ circle map has winding
$mn$.  The two line bundles in the displayed formula identify respectively
with $L_N^{\otimes mn}$ and
$(L_N^{\otimes n})^{\otimes m}$, which are canonically isomorphic.
\end{proof}

\subsection{Divisibility as a system of Lagrangian correspondences}

Let
\[
\Phi_n^\vee=\Hom(\Phi_n,\Q/\Z),
\qquad
D_n=\Phi_n\oplus\Phi_n^\vee,
\]
and equip $D_n$ with the hyperbolic quadratic form
\[
q_n^{\mathrm{hyp}}(a,\chi)=\chi(a).
\]
For $n\mid m$, let $i_{n,m}:\Phi_n\hookrightarrow\Phi_m$ be the
canonical inclusion and let
$r_{m,n}:\Phi_m^\vee\to\Phi_n^\vee$ be restriction.

\begin{definition}[Denominator correspondence]
Define
\[
\mathcal C_{n,m}
=
\left\{
\bigl((a,\chi),(i_{n,m}a,\psi)\bigr):
\chi=r_{m,n}\psi
\right\}
\subset D_n^-\oplus D_m.
\]
\end{definition}

\begin{theorem}[Lagrangian denominator history]
\label{thm:lagrangian-denominator-history}
The subgroup $\mathcal C_{n,m}$ is Lagrangian.  If
$n\mid m\mid\ell$, then relational composition satisfies
\[
\mathcal C_{m,\ell}\circ\mathcal C_{n,m}
=\mathcal C_{n,\ell}.
\]
Hence the divisibility category maps functorially to the category of finite
quadratic groups and Lagrangian correspondences.

Moreover, $\widehat\Z^\times$ acts by
\[
(a,\chi)\longmapsto(ua,u^{-1}\chi),
\]
preserves $q_n^{\mathrm{hyp}}$, and preserves every
$\mathcal C_{n,m}$.  The Lagrangian history is therefore canonically
cyclotomic-equivariant.
\end{theorem}

\begin{proof}
On $\mathcal C_{n,m}$ one has
\[
-q_n^{\mathrm{hyp}}(a,\chi)
+q_m^{\mathrm{hyp}}(i_{n,m}a,\psi)
=-\chi(a)+\psi(i_{n,m}a)=0.
\]
Thus the subgroup is isotropic.  Its cardinality is
$|\Phi_n||\Phi_m^\vee|=nm$, which is the square root of
$|D_n^-\oplus D_m|=n^2m^2$; hence it is maximal isotropic.
Composition follows from functoriality of inclusion and restriction.
The unit action preserves the hyperbolic form because
$(u^{-1}\chi)(ua)=\chi(a)$, and inclusion and restriction commute with
the reductions of a profinite unit, so the correspondences are invariant.
\end{proof}

\section{Quantized denominator history and the universal UHF limit}
\label{sec:uhf-history}

The Lagrangian correspondence system is canonical at the Morita level.  A
strict inductive system of matrix algebras requires an additional choice.
For $n\mid m$, write $m=kn$ and define
\[
\iota_{n,m}:M_n(\C)\longrightarrow M_m(\C)
\]
on matrix units by
\[
\boxed{
\iota_{n,m}(E_{ab}^{(n)})
=\sum_{r=0}^{k-1}E_{a+nr,b+nr}^{(m)}.
}
\]
These embeddings are unital and satisfy
$\iota_{m,\ell}\iota_{n,m}=\iota_{n,\ell}$.

\begin{definition}[Perfected Heisenberg algebra]
\label{def:perfected-uhf}
Define
\[
\mathcal Q_\Omega
:=\varinjlim_{n\mid m}(M_n(\C),\iota_{n,m}).
\]
\end{definition}

\begin{theorem}[UHF algebra of the globally perfect denominator]
\label{thm:perfected-uhf}
The algebra $\mathcal Q_\Omega$ is the universal UHF algebra of
supernatural type $\Omega=\prod_pp^\infty$.  Its ordered $K$-theory is
\[
K_0(\mathcal Q_\Omega)
\cong(\Q,\Q_{\geq0},1),
\qquad
K_1(\mathcal Q_\Omega)=0.
\]
\end{theorem}

\begin{proof}
At stage $n$, identify $K_0(M_n(\C))$ with $\Z$ using a rank-one
projection as generator.  Under $\iota_{n,m}$ a rank-one projection maps
to rank $m/n$, so the connecting map is multiplication by $m/n$.  The
map
\[
\Z\ni r\text{ at stage }n
\longmapsto\frac rn\in\Q
\]
is compatible with all connecting maps and identifies the direct limit
with $\Q$.  Positive ranks give $\Q_{\geq0}$ and the identity projection
gives the order unit $1$.  Since $K_1(M_n(\C))=0$ at every stage and
$K$-theory commutes with inductive limits, $K_1$ vanishes.  The directed
set contains arbitrarily large powers of every prime, so the UHF
supernatural type is $\Omega$.
\end{proof}

\begin{theorem}[The profinite diagonal]
\label{thm:profinite-diagonal}
Let $\mathcal D_n\subset M_n(\C)$ be the diagonal algebra, identified with
$C(\Z/n\Z)$.  The embeddings restrict to pullback along reduction
$\Z/m\Z\to\Z/n\Z$.  Therefore
\[
\mathcal D_\Omega
:=\varinjlim_n\mathcal D_n
\cong C(\widehat\Z).
\]
By Pontryagin duality,
\[
\boxed{
\mathcal D_\Omega
\cong C^*(\Q/\Z)
\subset\mathcal Q_\Omega.
}
\]
The compact space $\widehat\Z$ is a Cantor space.
\end{theorem}

\begin{proof}
Continuous functions on
$\widehat\Z=\varprojlim_n\Z/n\Z$ are the norm completion of locally
constant functions, and those are exactly the direct limit of the finite
diagonal algebras.  Pontryagin duality identifies the group $C^*$-algebra
of the discrete group $\Q/\Z$ with continuous functions on its compact
dual $\widehat\Z$.  Finally, $\widehat\Z$ is compact, metrizable, totally
disconnected, and has no isolated points; hence it is a Cantor space.
\end{proof}

\subsection{The strict cyclotomic carry obstruction}

For an integer $x$ and $q\geq1$, let $\langle x\rangle_q$ denote the
representative of $x\bmod q$ in $\{0,\dots,q-1\}$.  Let $m=kn$ and let
$u\in(\Z/m\Z)^\times$.  Define
\[
u_n(a):=\langle ua\rangle_n,
\qquad
c_{u;n,m}(a)
:=\frac{\langle ua\rangle_m-\langle ua\rangle_n}{n}
\in\Z/k\Z.
\]
The quotient is integral because the two residues agree modulo $n$.

\begin{theorem}[Cyclotomic carry cocycle and failure of strict equivariance]
\label{thm:cyclotomic-carry-obstruction}
Let $U_q(u)$ be the permutation unitary on $\C[\Z/q\Z]$ defined by
$U_q(u)e_a=e_{\langle ua\rangle_q}$.  Then:
\begin{enumerate}[label=\textup{(\arabic*)}]
\item for $a\in\Z/n\Z$ and $r\in\Z/k\Z$,
\[
\langle u(a+nr)\rangle_m
=u_n(a)+n\langle ur+c_{u;n,m}(a)\rangle_k;
\]
\item the carry functions satisfy the crossed cocycle law
\[
c_{uv;n,m}(a)
=u\,c_{v;n,m}(a)+c_{u;n,m}(v_n(a))
\pmod k;
\]
\item on matrix units,
\[
\Ad(U_m(u))\iota_{n,m}(E_{ab})
=
\sum_{t\in\Z/k\Z}
E_{u_n(a)+nt,\,
  u_n(b)+n(t+c_{u;n,m}(b)-c_{u;n,m}(a))};
\]
\item the square
\[
\begin{tikzcd}[column sep=large,row sep=large]
M_n(\C) \arrow[r,"\iota_{n,m}"]
  \arrow[d,"\Ad(U_n(u))"']
& M_m(\C) \arrow[d,"\Ad(U_m(u))"]\\
M_n(\C) \arrow[r,"\iota_{n,m}"']
& M_m(\C)
\end{tikzcd}
\]
commutes on $E_{ab}$ exactly when
$c_{u;n,m}(a)=c_{u;n,m}(b)$.  It commutes on the entire diagonal for every
$u$, and on the entire matrix algebra exactly when
$c_{u;n,m}$ is identically zero.
\end{enumerate}
Consequently, $\widehat\Z^\times$ acts canonically on the profinite
diagonal $C(\widehat\Z)$, but the naive finite permutation operators do
not define a strict action on the chosen UHF inductive system.  A strict
lift requires coherent trivializations of the carry cocycle.
\end{theorem}

\begin{proof}
Write the residue of $ua$ modulo $m$ uniquely as
$u_n(a)+nc_{u;n,m}(a)$.  Adding $nur$ and reducing modulo $m=nk$ gives
(1).  Applying (1) successively to $v$ and then $u$, and comparing with
$uv$, gives (2).

Conjugating a summand of the matrix embedding gives
\[
U_m(u)E_{a+nr,b+nr}U_m(u)^{-1}
=E_{u_n(a)+n(ur+c_u(a)),\,
    u_n(b)+n(ur+c_u(b))},
\]
where fibre coordinates are taken modulo $k$.  Since $u$ is invertible
modulo $k$, reindex by $t=ur+c_u(a)$ to obtain (3).  The lower-left then
has zero relative fibre displacement, so equality on $E_{ab}$ is
precisely $c_u(a)=c_u(b)$.  For $a=b$ this always holds, proving diagonal
equivariance.  Equality for all matrix units forces $c_u$ to be constant;
since $c_u(0)=0$, it is identically zero.  Passing to the direct limit of
the diagonals gives the canonical action on $C(\widehat\Z)$.
\end{proof}


\begin{theorem}[The cohomology class of strict cyclotomic carry]
\label{thm:cyclotomic-carry-class}
Let $n\mid m$, put $k=m/n$, and set
\[
U_m=(\Z/m\Z)^\times,
\qquad
M_{n,m}=\operatorname{Map}(\Z/n\Z,\Z/k\Z).
\]
Give $M_{n,m}$ the left $U_m$-module structure
\[
(u\cdot f)(a)=u\,f(u_n^{-1}a),
\]
where $u$ acts on the codomain through reduction modulo $k$ and $u_n$ is
its reduction modulo $n$.  Define
\[
d_{u;n,m}(a)=c_{u;n,m}(u_n^{-1}a).
\]
Then:
\begin{enumerate}[label=\textup{(\arabic*)}]
\item $d_{uv}=d_u+u\cdot d_v$, so
\[
d_{n,m}\in Z^1(U_m,M_{n,m});
\]
\item changing the fibre origin over $a$ by a function
$t\in M_{n,m}$ changes $d_{n,m}$ by a $1$-coboundary;
\item the class
\[
[d_{n,m}]\in H^1(U_m,M_{n,m})
\]
vanishes exactly when the standard finite permutation actions admit a
coherent strict trivialization over the embedding
$M_n(\C)\hookrightarrow M_m(\C)$;
\item for $n=2$ and $m=4$, the class $[d_{2,4}]$ is nonzero.
\end{enumerate}
\end{theorem}

\begin{proof}
Using the crossed cocycle law of
Theorem~\ref{thm:cyclotomic-carry-obstruction}, evaluate at
$(uv)_n^{-1}a=v_n^{-1}u_n^{-1}a$:
\[
\begin{aligned}
d_{uv}(a)
&=c_{uv}(v_n^{-1}u_n^{-1}a)\\
&=u\,c_v(v_n^{-1}u_n^{-1}a)
  +c_u(u_n^{-1}a)\\
&=(u\cdot d_v)(a)+d_u(a).
\end{aligned}
\]
This proves (1).

Write a point of the fibre over $a$ as $(a,r)$ with
$r\in\Z/k\Z$.  The unit $u$ acts by
\[
(a,r)\longmapsto(u_na,ur+c_u(a)).
\]
If the fibre origin is changed by $r'=r+t(a)$, the new carry is
\[
c'_u(a)=c_u(a)+t(u_na)-u t(a).
\]
After the reindexing defining $d_u$, this is addition of a group
$1$-coboundary.  Hence a simultaneous zero-carry gauge exists exactly when
$[d_{n,m}]=0$, proving (2)--(3).

For $n=2,m=4$, one has $U_4=\{1,3\}$ and $k=2$.  The action of $U_4$ on
$M_{2,4}$ is trivial because $3\equiv1\pmod2$ on both the domain and the
codomain.  Directly,
\[
d_3(0)=0,
\qquad
d_3(1)=\frac{3-1}{2}=1\pmod2.
\]
Thus $d_3$ is a nonzero homomorphism $U_4\to M_{2,4}$.  For a trivial
action every $1$-coboundary is zero, so its cohomology class is nonzero.
\end{proof}

\begin{remark}[Why the class is useful]
The obstruction is represented by a functorial cohomology class rather
than by a single failed square.  It can therefore be pulled back, compared
between refinements, and tested for strictifiability by ordinary group
cohomology.
\end{remark}

\begin{remark}[Morita equivariance versus strict equivariance]
The Lagrangian correspondence system of
Theorem~\ref{thm:lagrangian-denominator-history} is strictly
$\widehat\Z^\times$-equivariant.  The obstruction appears only after one
chooses the matrix embeddings $\iota_{n,m}$.  This is the precise sense in
which denominator history is canonical at the correspondence/Morita level
but carries a nontrivial arithmetic correction at the strict operator
level.
\end{remark}

\begin{remark}[Bost--Connes coefficient layer]
The equality
\[
C^*(\Q/\Z)=C(\widehat\Z)=\mathcal D_\Omega
\]
identifies the Bost--Connes cyclotomic coefficient algebra with the
canonical diagonal of $\mathcal Q_\Omega$.  Adjoining the
$\N^\times$ endomorphisms and taking the corresponding crossed product
recovers the dynamical layer of the Bost--Connes construction.
\end{remark}

\subsection{Split support for anomaly lines}

For a complex vector space $V$, put
\[
\Spl(V)=V\sqcup\{\tau_V\},
\]
where $\tau_V$ is external absence and $0_V\in V$ is the supported zero.
A linear map $f:V\to W$ extends by
\[
\Spl(f)(\tau_V)=\tau_W,
\qquad
\Spl(f)(v)=f(v).
\]

\begin{proposition}[Defined vanishing versus absent amplitude]
\label{prop:split-anomaly-line}
The assignment $V\mapsto\Spl(V)$ is a functor from complex vector spaces
to pointed sets.  If $v$ is supplied and $f(v)=0_W$, then
\[
\Spl(f)(v)=0_W\neq\tau_W.
\]
Consequently, for an anomalous field theory whose value on a closed
$d$-manifold is a possibly zero vector in an anomaly line, split support
distinguishes
\[
\begin{array}{c|c}
\tau&\text{the tier or field assignment is absent},\\
0&\text{the path integral is defined and vanishes},\\
z\neq0&\text{a nonzero anomalous amplitude}.
\end{array}
\]
The ordinary projection $\Spl(V)\to V$ identifies the first two states,
whereas the support character separates them.
\end{proposition}

\begin{proof}
Identity maps and composition are preserved by the definition.  A supplied
vector always remains in the supported copy of the target, even when its
ordinary linear image is zero.  The final statement applies this observation
to the anomaly line assigned by a truncated morphism
$\mathbf1\to\mathcal A$.
\end{proof}

\subsection{Dual anomalies and Morita-level identity}

\begin{proposition}[Anomaly cancellation]
\label{prop:anomaly-cancellation}
The inverse of the finite quadratic theory associated with $q_n$ is the
theory associated with $-q_n$, and
\[
\mathcal A(q_n)\otimes\mathcal A(-q_n)\simeq\mathbf1.
\]
For the explicit bundle gerbe,
\[
DD(\mathcal G_N^\vee)=-DD(\mathcal G_N),
\]
so $\mathcal G_N\otimes\mathcal G_N^\vee$ is stably trivial.  On winding
sectors,
\[
L_N^{\otimes r}\otimes L_N^{\otimes(-r)}\cong\mathbf1.
\]
\end{proposition}

\begin{proof}
The finite theory is invertible and changing $q_n$ to $-q_n$ complex
conjugates every Gauss phase, hence gives the tensor inverse.  Dualizing a
bundle gerbe negates its Dixmier--Douady class.  The last isomorphism is the
evaluation pairing between a line bundle and its dual.
\end{proof}

\begin{theorem}[Morita identification of the finite models]
\label{thm:morita-univalent-finite}
After choosing the standard Lagrangian polarization of $\Phi_n^2$, there
are algebra isomorphisms
\[
\C_{\omega_n}[\Phi_n^2]
\cong
\End_\C(\mathscr H_n)
\cong
\End_\C H^0(E_\tau,L_1^{\otimes n}\otimes F_{\phi,\psi})
\cong M_n(\C).
\]
The matrix algebra is Morita equivalent to $\C$ through the bimodule
$\C^n$.  Therefore the finite phase, Heisenberg, theta, and matrix
presentations determine the same boundary object in the Morita bicategory.
In the univalent completion of that bicategory, this equivalence determines a path between the represented boundary objects, while the four concrete presentations retain their own elementwise coordinates.
\end{theorem}

\begin{proof}
The first and last algebra isomorphisms are
Theorem~\ref{thm:finite-stone-von-neumann}; the middle one is
Theorem~\ref{thm:winding-theta-realization}.  The standard
$(M_n(\C),\C)$-bimodule $\C^n$ and its dual implement the Morita
equivalence.  The final sentence is the defining interpretation of
equivalence in a univalent or Rezk completion.
\end{proof}

\begin{figure}[htbp]
\centering
\begin{tikzcd}[column sep=large,row sep=large]
(\Phi_n,q_n) \arrow[r,"\mathrm{deloop}"]
  \arrow[d,"\mathrm{transgress}_{T^2}"']
& \widetilde q_n:K(\Phi_n,2)\to K(\Q/\Z,4)
  \arrow[r,"\mathrm{finite\ path\ integral}"]
& \mathcal A_n \\
(\Phi_n^2,\Omega_n) \arrow[r,"\omega_n"']
& \C_{\omega_n}[\Phi_n^2]
  \arrow[r,"\rho_n"']
& \End_\C(\mathscr H_n)\cong M_n(\C)
\end{tikzcd}
\caption{Finite denominator phase, Eilenberg--Mac Lane delooping,
gerbe integration, surface transgression, and Heisenberg quantization.}
\label{fig:phase-gerbe-quantization}
\end{figure}


\section{Monomial spectral charts and geometric Langlands}
\label{sec:monomial-langlands}

Let $X$ be a smooth complete complex curve, let
$T_n\subset\GL_n$ be the diagonal torus, and let
$N_n=N_{\GL_n}(T_n)=T_n\rtimes S_n$.  The monomial identity
\[
\GL_n(G(K))\cong N_n(K)
\]
turns split linear geometry into a normalizer-of-torus chart on the spectral
side of geometric Langlands.

\subsection{Monomial local systems as line local systems on covers}

Write $\LocSys_H^{\mathrm{Betti}}(X)$ for the groupoid of representations
$\pi_1(X)\to H$ modulo conjugation.

\begin{theorem}[Monomial spectral-cover equivalence]
\label{thm:monomial-spectral-cover}
There is a natural equivalence of groupoids
\[
\LocSys_{N_n}^{\mathrm{Betti}}(X)
\simeq
\left\{
(\pi:Y\to X,\mathcal L):
\begin{array}{l}
\pi\text{ is an }n\text{-sheeted finite covering},\\
\mathcal L\text{ is a rank-one complex local system on }Y
\end{array}
\right\}.
\]
Under this equivalence, the morphism induced by
$N_n\hookrightarrow\GL_n$ is
\[
(\pi:Y\to X,\mathcal L)\longmapsto\pi_*\mathcal L.
\]
\end{theorem}

\begin{proof}
Let $\rho:\pi_1(X,x)\to T_n\rtimes S_n$.  Projection to $S_n$ gives an
action of $\pi_1(X,x)$ on $\{1,\dots,n\}$ and hence an $n$-sheeted
cover $Y\to X$.  The $T_n$ component is a multiplicative $1$-cocycle
with coefficients in $(\C^\times)^n$, where $\pi_1(X,x)$ permutes the
coordinates through the covering action.  Such a cocycle is exactly the
descent datum for a rank-one local system on $Y$.

Conversely, choose the basis of the fibre of $\pi_*\mathcal L$ indexed by
the sheets over $x$.  Monodromy permutes the sheets and multiplies each
basis line by its rank-one holonomy, hence lies in $T_n\rtimes S_n$.
Changing the sheet basis conjugates by $N_n$, and the two constructions are
inverse on objects and morphisms.
\end{proof}

\begin{corollary}[Split local-system morphism]
\label{cor:split-locsys-morphism}
The ring reflection and the inclusion of structure groups give natural
morphisms
\[
\LocSys_{\GL_n(G(\C))}^{\mathrm{Betti}}(X)
\cong
\LocSys_{N_n}^{\mathrm{Betti}}(X)
\longrightarrow
\LocSys_{\GL_n}^{\mathrm{Betti}}(X).
\]
The first stack retains the finite cover and rank-one sheet data; the second
retains only the induced rank-$n$ local system.
\end{corollary}

\subsection{The denominator phase tower inside rank-one geometric Langlands}

Let $E=\C/(\Z+\tau\Z)$ be an elliptic curve, equipped with the oriented
basis of $H_1(E,\Z)$ represented by $1$ and $\tau$.  Put
\[
\Phi_n=n^{-1}\Z/\Z,
\qquad
D_n=\Phi_n\oplus\Phi_n^\vee,
\]
with alternating pairing
\[
\Omega_n((a,\chi),(b,\psi))=\psi(a)-\chi(b).
\]

\begin{theorem}[Torsion spectral chart]
\label{thm:torsion-spectral-chart}
The chosen basis and the exponential map give compatible symplectic
isomorphisms
\[
D_n
\xrightarrow{\sim}
H^1(E,\mu_n)
\xrightarrow{\sim}
\LocSys_{\Gm}^{\mathrm{Betti}}(E)[n].
\]
Under these isomorphisms, $\Omega_n$ is the cup-product pairing evaluated
on $[E]$.  For $m=qn$, define the rigidified transition
\[
\jmath_{n,m}(r/n,\chi_s^{(n)})
=(r/n,\chi_{qs}^{(m)}).
\]
These transitions agree with $\mu_n\hookrightarrow\mu_m$ under the displayed
isomorphisms and satisfy
\[
\Omega_m(\jmath_{n,m}x,\jmath_{n,m}y)=q\,\Omega_n(x,y).
\]
Moreover,
\[
\colim_nD_n\cong H^1(E,\Q/\Z)\cong(\Q/\Z)^2,
\qquad
\liminv_nH_1(E,\Z/n\Z)\cong\widehat\Z^{\,2}.
\]
\end{theorem}

\begin{proof}
Write $a=r/n$ and let $\chi_s(r/n)=rs/n$ for $r,s\in\Z/n\Z$.
The map
\[
(a,\chi_s)\longmapsto(\zeta_n^r,\zeta_n^s)
\]
identifies $D_n$ with
$\Hom(H_1(E,\Z),\mu_n)=H^1(E,\mu_n)$, which is exactly the $n$-torsion
of the Betti rank-one local-system torus $(\C^\times)^2$.  The cup product
of $(r,s)$ and $(r',s')$, evaluated on the oriented fundamental class, is
$rs'-sr'$ modulo $n$, which becomes
$\chi_{s'}(a)-\chi_s(a')$.  For $m=qn$, the inclusion $\mu_n\hookrightarrow\mu_m$ sends
$(\zeta_n^r,\zeta_n^s)$ to
$(\zeta_m^{qr},\zeta_m^{qs})$, which is exactly the stated transition.
The pairing scales by $q$ after substitution in the determinant formula.
Filtered colimit and inverse limit give the last two formulas
coordinatewise.
\end{proof}

Let $L$ be a degree-one line bundle on $E$.  Its $n$th tensor power has
theta group
\[
1\longrightarrow\C^\times
\longrightarrow\mathcal G(L^{\otimes n})
\longrightarrow E[n]
\longrightarrow1.
\]

\begin{theorem}[Theta-group realization of the torsion spectral chart]
\label{thm:theta-group-spectral-chart}
The principal polarization identifies $E[n]$ with the torsion spectral
chart of Theorem~\ref{thm:torsion-spectral-chart}.  Under a symplectic level
structure, the commutator pairing of $\mathcal G(L^{\otimes n})$ is
\[
((a,b),(c,d))\longmapsto\exp\!\left(\frac{2\pi i(ad-bc)}n\right).
\]
The weight-one theta-group representation on
$H^0(E,L^{\otimes n})$ is the Schr\"odinger representation of the finite
Heisenberg group, and hence
\[
\End H^0(E,L^{\otimes n})
\cong
\C_{\omega_n}[D_n]
\cong M_n(\C).
\]
\end{theorem}

\begin{proof}
For a degree-one polarization, the homomorphism
$E\to\operatorname{Pic}^0(E)$ associated with $L$ is an isomorphism, and
the kernel associated with $L^{\otimes n}$ is $E[n]$.  Mumford's theta-group
theorem identifies the commutator with the Weil pairing.  In a symplectic
basis the Weil pairing is the displayed determinant phase.  The unique
irreducible weight-one representation has dimension $n$ and is the finite
Schr\"odinger representation.  The final algebra identities follow from
Theorem~\ref{thm:finite-stone-von-neumann}.
\end{proof}

Thus the finite phase spaces used in the gerbe and UHF constructions are
literal finite closed strata of the rank-one geometric Langlands spectral
stack, and their theta modules are the corresponding Heisenberg
quantizations.

\subsection{A support-refined Hitchin map}

Let $K_X$ denote the canonical line bundle and let
\[
\mathscr K_X^\bullet=\bigoplus_{d\ge0}K_X^{\otimes d}
\]
be its graded symmetric algebra.  Apply split globalization degreewise,
with the tensor product of supported sections and $\tau$ absorbing.

\begin{definition}[Split monomial Higgs bundle]
A split monomial Higgs bundle of rank $n$ is local data
$(g_{\alpha\beta},\Phi_\alpha)$ on an open cover $\{U_\alpha\}$ such that
\begin{enumerate}[label=\textup{(\roman*)}]
\item $g_{\alpha\beta}\in N_n(\mathcal O_X(U_{\alpha\beta}))$ are
transition functions;
\item $\Phi_\alpha\in M_n(\Spl(K_X)(U_\alpha))$;
\item on overlaps,
$\Phi_\beta=g_{\alpha\beta}\Phi_\alpha g_{\alpha\beta}^{-1}$.
\end{enumerate}
Morphisms are monomial gauge transformations preserving the local Higgs
matrices.  Denote the resulting groupoid by
$\Higgs_{N_n}^{\mathrm{sp}}(X)$.
\end{definition}

\begin{theorem}[Split Hitchin map]
\label{thm:split-hitchin-map}
The local split characteristic coefficients glue and define a natural map
\[
h_{\mathrm{sp}}:
\Higgs_{N_n}^{\mathrm{sp}}(X)
\longrightarrow
\prod_{k=1}^n\Gamma\bigl(X,\Spl(K_X^{\otimes k})\bigr),
\qquad
(E,\Phi)\longmapsto(c_1^{\mathrm{sp}}(\Phi),\dots,c_n^{\mathrm{sp}}(\Phi)).
\]
Ring reflection gives the commutative square
\[
\begin{tikzcd}[column sep=large,row sep=large]
\Higgs_{N_n}^{\mathrm{sp}}(X)
  \arrow[r,"h_{\mathrm{sp}}"] \arrow[d,"p"']
& \displaystyle\prod_{k=1}^n\Gamma(X,\Spl(K_X^{\otimes k}))
  \arrow[d,"p"] \\
\Higgs_{N_n}(X)
  \arrow[r,"h"]
& \displaystyle\bigoplus_{k=1}^nH^0(X,K_X^{\otimes k}),
\end{tikzcd}
\]
where $h$ is the restriction of the classical $\GL_n$ Hitchin map.
At every geometric point $x$, the $k$th split coefficient is
\begin{enumerate}[label=\textup{(\alph*)}]
\item $\tau$ exactly when the support digraph of $\Phi_x$ has no cycle
cover on any $k$-vertex subset;
\item $e$ exactly when such a cycle cover exists and the ordinary Hitchin
coefficient vanishes at $x$;
\item a nonzero supported scalar otherwise.
\end{enumerate}
\end{theorem}

\begin{proof}
On overlaps, the Higgs matrices are related by monomial conjugation.
Theorem~\ref{thm:split-characteristic} shows that every local coefficient is
unchanged, so the coefficients glue in the indicated tensor powers of
$K_X$.  Applying ring reflection locally gives the ordinary characteristic
polynomial and hence the classical Hitchin coefficients.  The three stalk
cases are Theorem~\ref{thm:split-characteristic} after passage to the
residue field at $x$.
\end{proof}

The classical Hitchin base is the affine quotient of the Higgs field by the
fundamental invariant polynomials of $\mathfrak{gl}_n$.  The split Hitchin
map retains, in addition, the cycle-cover incidence data erased by that
affine quotient.

\subsection{Restriction of the geometric Langlands spectral functor}

The 2024 construction of Gaitsgory--Raskin supplies the coarse Betti
geometric Langlands functor
\[
\mathbb L^{\mathrm{Betti}}_{G,\mathrm{coarse}}:
\operatorname{Shv}^{\mathrm{Betti}}_{1/2,\Nilp}(\Bun_G)
\longrightarrow
\QCoh\bigl(\LocSys_{\check G}^{\mathrm{Betti}}(X)\bigr).
\]
For $G=\GL_n$, the dual group is again $\GL_n$.  Let
\[
j_X:\LocSys_{N_n}^{\mathrm{Betti}}(X)
\longrightarrow
\LocSys_{\GL_n}^{\mathrm{Betti}}(X)
\]
be the morphism of Corollary~\ref{cor:split-locsys-morphism}.

\begin{theorem}[Monomial spectral restriction]
\label{thm:monomial-langlands-restriction}
The composite
\[
\mathbb L^{\mathrm{mon}}_n
:=
j_X^*\circ\mathbb L^{\mathrm{Betti}}_{\GL_n,\mathrm{coarse}}
:
\operatorname{Shv}^{\mathrm{Betti}}_{1/2,\Nilp}(\Bun_{\GL_n})
\longrightarrow
\QCoh\bigl(\LocSys_{N_n}^{\mathrm{Betti}}(X)\bigr)
\]
is a canonical monomial spectral restriction of the coarse geometric
Langlands functor.  On the affine character charts associated with a
presentation of $\pi_1(X)$, pullback of the trace-word coordinates is given
by the coloured cycle formulas of
Theorems~\ref{thm:split-trace-word} and
\ref{thm:cycle-invariant-ring}.
\end{theorem}

\begin{proof}
A morphism of spectral stacks induces pullback on quasi-coherent sheaves, so
the displayed composite is defined functorially.  On a representation
chart, $j_X$ is induced by the inclusion $N_n\hookrightarrow\GL_n$.
Corollary~\ref{cor:monomial-character-morphism} computes the pullback of
every trace-character coordinate as a cycle-orbit sum, and
Theorem~\ref{thm:split-trace-word} supplies its support refinement.
\end{proof}

The morphisms in this section isolate a tractable spectral chart: finite
covers provide the Weyl-group component, line local systems provide the
torus component, classical invariant theory supplies explicit coordinates,
and the split layer records which cycle contributions are absent, present
and cancelling, or nonzero.

\section{Beurling realization of the reciprocal-root skeleton}
\label{sec:beurling-realization}

We now connect the perfect stalk $\Fone[T^{\Q_+}]$ with the reciprocal
Nyman--Beurling system studied in \cite{Schultze}.  Let
\[
\cH=L^2((1,\infty),y^{-2}\,dy)
\]
and, for $0<\alpha<1$, define
\[
F_\alpha(y)=\alpha\lfloor y\rfloor-\lfloor\alpha y\rfloor.
\]
Since
\[
F_\alpha(y)=\{\alpha y\}-\alpha\{y\},
\]
these functions are bounded and therefore belong to $\cH$.

For $Q\geq2$, put
\[
V_Q=\operatorname{span}\{F_{a/Q}:1\leq a\leq Q-1\},
\]
let $P_Q$ be orthogonal projection onto $V_Q$, and define
\[
r_Q=(I-P_Q)1,
\qquad
x_Q=\|r_Q\|^2
=\operatorname{dist}_{\cH}(1,V_Q)^2.
\]
Define the reciprocal subspace and distance
\[
\cR_Q=\operatorname{span}\{F_{1/n}:2\leq n\leq Q\},
\qquad
 d_Q=\operatorname{dist}_{\cH}(1,\cR_Q).
\]

\begin{definition}[Beurling realization of the perfect stalk]
\label{def:beurling-realization}
On the complex span of monomials $T^q$, $q\in\Q\cap(0,1)$, define
\[
\mathfrak B_\Omega(T^q)=F_q.
\]
The reciprocal-root skeleton is
\[
\mathscr R_\Omega=\{T^{1/n}:n\geq2\}
\subset\Fone[T^{\Q_+}].
\]
This is a linear analytic realization of stalk monomials; it is not an
$\Fone$-algebra morphism.
\end{definition}

\begin{theorem}[Direct Frobenius-refinement domination]
\label{thm:direct-beurling-domination}
For every $Q\geq2$,
\[
\cR_Q\subseteq V_{\Lambda_Q}
\]
and consequently
\[
\boxed{x_{\Lambda_Q}\leq d_Q^2.}
\]
\end{theorem}

\begin{proof}
For $2\leq n\leq Q$,
\[
\frac1n=\frac{\Lambda_Q/n}{\Lambda_Q},
\]
where $1\leq\Lambda_Q/n\leq\Lambda_Q-1$.  Hence
$F_{1/n}\in V_{\Lambda_Q}$, proving the subspace inclusion.  Distance to
a larger closed subspace cannot increase, so
\[
\operatorname{dist}(1,V_{\Lambda_Q})
\leq\operatorname{dist}(1,\cR_Q)=d_Q.
\]
Squaring gives the result.
\end{proof}

\begin{corollary}[Strengthened refinement alternative]
\label{cor:strong-refinement-alternative}
For each $Q\geq2$, either
\[
x_Q\leq4d_Q^2,
\]
or
\[
x_{\Lambda_Q}<\frac14x_Q
\qquad\text{and}\qquad
x_Q-x_{\Lambda_Q}>\frac34x_Q.
\]
\end{corollary}

\begin{proof}
If the first inequality fails, then $d_Q^2<x_Q/4$.  Apply Theorem
\ref{thm:direct-beurling-domination}.
\end{proof}

The estimate is stronger than the quadratic decrement proved in
\cite{Schultze}; it follows from the exact root inclusion and does not
require the unit-step visibility lemma.

\begin{lemma}[Reciprocal atoms in Nyman form]
\label{lem:nyman-form}
The map
\[
U:\cH\longrightarrow L^2(0,1),
\qquad
(Uf)(x)=f(1/x),
\]
is an isometry and satisfies
\[
UF_{1/n}(x)
=
\rho\!\left(\frac1{nx}\right)
-\frac1n\rho\!\left(\frac1x\right),
\]
where $\rho(t)=t-\lfloor t\rfloor$.
\end{lemma}

\begin{proof}
The substitution $y=1/x$ gives $y^{-2}dy=-dx$, proving the isometry.
Furthermore,
\[
\begin{aligned}
F_{1/n}(1/x)
&=\frac1n\left(\frac1x-\rho(1/x)\right)
 -\left(\frac1{nx}-\rho(1/(nx))\right)\\
&=\rho(1/(nx))-\frac1n\rho(1/x).
\end{aligned}
\]
\end{proof}

\begin{theorem}[Arithmetic-site form of the integer
Nyman--B\'aez-Duarte criterion]
\label{thm:arithmetic-site-nyman}
Let $Q_{j+1}=\Lambda_{Q_j}$ with $Q_0\geq3$.  The following are equivalent:
\begin{enumerate}[label=\textup{(\roman*)}]
\item the Riemann hypothesis;
\item $d_Q\to0$ as $Q\to\infty$;
\item $x_{Q_j}\to0$ as $j\to\infty$;
\item
\[
1\in
\overline{
\mathfrak B_\Omega
\bigl(\operatorname{span}_{\C}\mathscr R_\Omega\bigr)
}^{\,\cH}.
\]
\end{enumerate}
\end{theorem}

\begin{proof}
By the integer-dilation form of the Nyman--Beurling criterion due to
B\'aez-Duarte \cite{Nyman,Beurling,BaezDuarte}, the Riemann hypothesis is
equivalent to density of the functions in Lemma \ref{lem:nyman-form} at
the target $1$.  This is exactly the equivalence of (i), (ii), and (iv).

Assume (ii).  By Theorem \ref{thm:direct-beurling-domination},
\[
x_{Q_{j+1}}\leq d_{Q_j}^2\longrightarrow0,
\]
so (iii) follows.

Conversely, every $V_{Q_j}$ is contained in the full Nyman--Beurling
space generated by the functions
$\rho(\alpha/x)-\alpha\rho(1/x)$, $0<\alpha<1$.  If (iii) holds, then $1$
lies in the closure of that space, and the Nyman--Beurling theorem implies
(i).
\end{proof}

\begin{remark}[Arithmetic-site form of the analytic endpoint]
Theorem~\ref{thm:arithmetic-site-nyman} identifies the Riemann hypothesis
with one precise approximation statement: density of the reciprocal-root
Beurling realization at the globally Frobenius-perfect stalk.  The estimate
$x_{\Lambda_Q}\leq d_Q^2$ transfers integer reciprocal approximation
directly into the LCM Frobenius tower.
\end{remark}

\begin{theorem}[Mellin transform and Frobenius characters]
\label{thm:mellin-frobenius-characters}
For $0<\alpha<1$ and $\Re(s)>1$,
\[
\int_1^\infty F_\alpha(y)y^{-s-1}\,dy
=
\frac{\zeta(s)}s\bigl(\alpha-\alpha^s\bigr).
\]
In particular,
\[
\mathcal MF_{1/n}(s)
=
\frac{\zeta(s)}s
\bigl(\chi_1(n)-\chi_s(n)\bigr),
\qquad
\chi_s(n)=n^{-s}.
\]
\end{theorem}

\begin{proof}
For $\Re(s)>1$,
\[
\lfloor y\rfloor=\sum_{m\geq1}\mathbf1_{[m,\infty)}(y)
\]
with absolute convergence after integration, whence
\[
\int_1^\infty\lfloor y\rfloor y^{-s-1}\,dy
=\frac1s\sum_{m\geq1}m^{-s}
=\frac{\zeta(s)}s.
\]
Similarly,
\[
\int_1^\infty\lfloor\alpha y\rfloor y^{-s-1}\,dy
=\frac1s\sum_{m\geq1}(m/\alpha)^{-s}
=\frac{\alpha^s\zeta(s)}s.
\]
Subtracting gives the first identity.  The second is the specialization
$\alpha=1/n$.
\end{proof}

\section{Frobenius local systems and the helix multiplier}
\label{sec:helix-local-system}

The helix construction of \cite{Lavery} is most naturally expressed on the
Connes--Consani periodic orbit.  Fix a prime $p$ and write
\[
C_p=\R_+^\times/p^\Z
\cong\R/(\log p)\Z.
\]
Let $u=\log\lambda$ be the coordinate on the universal cover.

\begin{definition}[Frobenius phase line]
\label{def:frobenius-phase-line}
For $t\in\R$, let $L_{p,t}\to C_p$ be the complex line bundle
\[
L_{p,t}=(\R\times\C)/\sim,
\]
where
\[
(u,z)\sim(u+k\log p,p^{-ikt}z),
\qquad k\in\Z.
\]
Let $\Delta_p^{1/2}$ denote the real one-dimensional factor of automorphy
$p^{1/2}$.  Define
\[
\cL^{\mathrm{hel}}_{p,t}=\Delta_p^{1/2}\otimes L_{p,t},
\qquad
\cL^{\mathrm{crit}}_{p,t}=\Delta_p^{-1/2}\otimes L_{p,t}.
\]
\end{definition}

\begin{theorem}[Helix phase as Frobenius holonomy]
\label{thm:helix-holonomy}
The function
\[
\psi_t(u)=e^{-itu}
\]
is an equivariant lift of a section of $L_{p,t}$ and satisfies
\[
\psi_t(u+\log p)=p^{-it}\psi_t(u).
\]
Consequently,
\[
\operatorname{Hol}(L_{p,t})=p^{-it},
\]
while the two twists have Frobenius factors
\[
p^{1/2-it}
\qquad\text{and}\qquad
p^{-1/2-it},
\]
respectively.
\end{theorem}

\begin{proof}
Direct calculation gives
\[
\psi_t(u+k\log p)
=e^{-it(u+k\log p)}
=p^{-ikt}\psi_t(u),
\]
which is exactly the factor-of-automorphy condition.  Tensoring with the
half-density factors multiplies the holonomy by $p^{\pm1/2}$.
\end{proof}

\begin{corollary}[Dual chirality and determinant one]
\label{cor:dual-chirality}
Complex conjugation gives
\[
L_{p,-t}\cong L_{p,t}^{\vee},
\]
and therefore
\[
\det(L_{p,t}\oplus L_{p,-t})\cong\mathbf1.
\]
In a local trivialization, Frobenius acts by
\[
\begin{pmatrix}
p^{-it}&0\\0&p^{it}
\end{pmatrix},
\]
which is unitary of determinant one.
\end{corollary}

\begin{proof}
The inverse of the character $p^{-it}$ is $p^{it}$, which is also its
complex conjugate.  Hence the second line is the dual of the first, and
the determinant line is their tensor product.
\end{proof}

\begin{proposition}[Descent and winding closure]
\label{prop:helix-descent}
The equivariant section $\psi_t$ descends to a scalar-valued function on
$C_p$ if and only if
\[
p^{-it}=1,
\]
equivalently
\[
t\in\frac{2\pi}{\log p}\Z.
\]
On the rectangular Tate torus
\[
E_p\cong C_p\times S^1
\]
from \cite{CCAbsolute}, the trajectory
\[
u\longmapsto([u],[-tu])
\]
is closed if $t\log p/(2\pi)\in\Q$ and dense if this quotient is
irrational.
\end{proposition}

\begin{proof}
Scalar descent requires invariance under the deck translation
$u\mapsto u+\log p$, which is equivalent to $p^{-it}=1$.  This gives the
stated integral condition.  The final assertion is the standard
classification of linear flows on a two-dimensional torus: the orbit is
closed precisely when the two periods are rationally dependent.
\end{proof}

\begin{corollary}[Split fixed-locus jump for Frobenius phase]
\label{cor:helix-fixed-locus-jump}
Let $u_{p,t}=p^{-it}\in\C^\times$.  Multiplication by $u_{p,t}$ on
$G(\C)$ has fixed locus
\[
\operatorname{Fix}_{G(\C)}(u_{p,t})=
\begin{cases}
G(\C),&t\log p\in2\pi\Z,\\
\{\tau,e\}\cong\B,&t\log p\notin2\pi\Z.
\end{cases}
\]
When the holonomy is nontrivial, the two chiral phase lines have fixed locus $\B^2$.
Thus scalar descent is exactly the parameter locus at which invariant
amplitude jumps from Boolean support to the full split complex fibre.
\end{corollary}

\begin{proof}
Combine Proposition~\ref{prop:helix-descent} with
Theorem~\ref{thm:fixed-locus-unit}.  The two-chirality statement is the
Cartesian product of the fixed loci for $u_{p,t}$ and $u_{p,t}^{-1}$.
\end{proof}

\begin{proposition}[Supernatural endpoint obstruction]
\label{prop:supernatural-holonomy-obstruction}
Consider the $p$-power ray
\[
1,p,p^2,\dots,p^\infty
\]
in the supernatural number space, and the unitary Frobenius character
\[
\chi_t(p^k)=p^{-ikt}.
\]
The sequence $\chi_t(p^k)$ has a limit in $S^1$ if and only if
$p^{-it}=1$.  Thus the same scalar is simultaneously:
\begin{enumerate}[label=\textup{(\roman*)}]
\item the holonomy of $L_{p,t}$;
\item the obstruction to scalar descent on $C_p$;
\item the obstruction to continuous extension of the phase along the
$p$-power ray to $p^\infty$.
\end{enumerate}
The full helix multiplier $p^{k(1/2-it)}$ never converges in $\C^\times$.
\end{proposition}

\begin{proof}
Put $z=p^{-it}\in S^1$.  If $z^k\to L$, then also
$z^{k+1}\to L$, while $z^{k+1}=zz^k\to zL$.  Since $L\neq0$, one has
$z=1$.  The converse is immediate.  Finally,
$|p^{k(1/2-it)}|=p^{k/2}\to\infty$.
\end{proof}

\begin{remark}
The three-dimensional growing helix is an extrinsic metric realization.
The invariant arithmetic-geometric object is the Frobenius-equivariant
line with factor $n^{1/2-it}$, together with its dual chirality.  This
reformulation preserves the Lean-verified determinant and multiplicativity
statements of \cite{Lavery} while separating them from any claim about the
location of all analytic zeros.
\end{remark}

\section{Split local moduli and Boolean place cubes}
\label{sec:split-local-moduli}

Connes--Consani show that, for $H_p=\Z[1/p]$ and each
$v\in\{p,\infty\}$, the nontrivial local points form a principal
homogeneous space $\cM_p^v$ under the local Weil group
\[
W_p=\Q_p^\times,
\qquad
W_\infty=\C^\times.
\]
The split-support refinement distinguishes absence from the neutral local
morphism before taking the Weil quotient.

\begin{definition}[Split completion of a torsor]
\label{def:split-torsor-completion}
Let a group $W$ act simply transitively on a nonempty set $M$.  Define
\[
\widetilde M=\{\tau_M\}\sqcup\{\mathbf1_M\}\sqcup M,
\]
where $W$ fixes $\tau_M$ and $\mathbf1_M$ and retains its torsor action on
$M$.  Define the support map
\[
\sigma_M:\widetilde M\to\B
\]
by
\[
\sigma_M(\tau_M)=0,
\qquad
\sigma_M(\mathbf1_M)=1,
\qquad
\sigma_M(m)=1\quad(m\in M).
\]
\end{definition}

\begin{theorem}[Univalent quotient of a split torsor]
\label{thm:split-torsor-groupoid}
For the action groupoid one has an equivalence
\[
\boxed{
[\widetilde M/W]\simeq BW\sqcup BW\sqcup *.
}
\]
The two copies of $BW$ correspond to absence and the supported neutral
point; the terminal component corresponds to the free torsor orbit.
\end{theorem}

\begin{proof}
The action preserves the disjoint decomposition
$\widetilde M=\{\tau_M\}\sqcup\{\mathbf1_M\}\sqcup M$.  A fixed singleton
has action groupoid $BW$.  Since the action on $M$ is simply transitive,
its action groupoid is equivalent to the terminal groupoid.  Taking the
disjoint union gives the result.
\end{proof}

\begin{corollary}[The split $p$--$\infty$ square]
\label{cor:p-infinity-square}
For $H_p=\Z[1/p]$, define
\[
\widetilde\cM_{p,\infty}
=\widetilde\cM_p^p\times\widetilde\cM_p^\infty.
\]
Its coarse orbit stratification has nine states
\[
\{\text{absent},\text{neutral},\text{nontrivial}\}^{\{p,\infty\}},
\]
and the support map
\[
\widetilde\cM_{p,\infty}\longrightarrow\B^2
\]
forgets only the distinction between neutral and nontrivial supported
states.  At groupoid level,
\[
[\widetilde\cM_{p,\infty}/(W_p\times W_\infty)]
\simeq
(BW_p\sqcup BW_p\sqcup *)
\times
(BW_\infty\sqcup BW_\infty\sqcup *).
\]
\end{corollary}

\begin{proof}
Apply Theorem \ref{thm:split-torsor-groupoid} independently at the two
places and take products.  The coarse states are the nine products of the
three one-place states, while the Boolean support quotient records only
whether each coordinate is absent or supported.
\end{proof}

\begin{theorem}[Finite Boolean place cubes]
\label{thm:boolean-place-cubes}
Let $H$ be an arithmetic-site point and let
$J\subseteq\Sigma_H$ be finite.  Suppose each $v\in J$ carries a
nontrivial local torsor $\cM_H^v$ under $W_v$.  Then the split local moduli
object
\[
\widetilde\cM_H^J
=\prod_{v\in J}\widetilde\cM_H^v
\]
has a canonical support map to $\B^J$.  The fibre over a mask
$A\subseteq J$ is the disjoint union, over all choices of neutral versus
nontrivial state on $A$, of products of the corresponding local torsors.
The quotient groupoid retains the stabilizer $BW_v$ at every absent or
neutral coordinate.
\end{theorem}

\begin{proof}
Take the product of the one-place support maps.  A coordinate with Boolean
value $0$ is forced to be absent.  A coordinate with Boolean value $1$ may
be neutral or belong to the nontrivial torsor.  The action-groupoid
statement follows from Theorem \ref{thm:split-torsor-groupoid} by finite
products.
\end{proof}

\begin{remark}[Why orbit sets are insufficient]
Passing directly to coarse orbit sets turns the three one-place components
into three points.  The action groupoid retains $W_v$ as the identity type
of the absent and neutral states.  This is the precise univalent content of
the split completion: the refinement is stacky, not merely set-theoretic.
\end{remark}

\section{Sheafwise split support, the split arithmetic site, and the split scaling site}
\label{sec:split-scaling-site}

The correct sheafwise construction is internal.  Let $\mathscr E$ be a
Grothendieck topos and let $A$ be a commutative semiring object of
$\mathscr E$.

\begin{definition}[Internal split globalization]
\label{def:internal-split}
Define the object
\[
\Spl_{\mathscr E}(A):=1\amalg A.
\]
On the $A$-summand retain the operations of $A$; let the point of the
$1$-summand be the new additive identity and multiplicative absorber.
\end{definition}

Because a topos is extensive, maps out of finite coproducts are determined
by their restrictions to the summands.  The semiring axioms can therefore
be checked on the finitely many support cases.

\begin{theorem}[Internal functoriality and universal maps]
\label{thm:internal-split-functor}
The assignment $A\mapsto\Spl_{\mathscr E}(A)$ is functorial on commutative
semiring objects.  It carries natural morphisms
\[
p_A:\Spl_{\mathscr E}(A)\to A,
\qquad
\chi_A:\Spl_{\mathscr E}(A)\to\B_{\mathscr E},
\]
where $p_A$ sends the new point to $0_A$ and $\chi_A$ records the coproduct
summand.  At every point $x$ of a topos with enough points,
\[
x^*\Spl_{\mathscr E}(A)\cong G(x^*A).
\]
\end{theorem}

\begin{proof}
For a semiring morphism $f:A\to A'$, define
$\Spl(f)$ to be the identity on the $1$-summand and $f$ on the supported
summand.  The case-by-case definitions make this a semiring morphism and
make composition immediate.  The maps $p_A$ and $\chi_A$ are likewise
specified on the two summands and preserve both operations.  An inverse
image functor of a geometric morphism preserves finite coproducts and
finite limits, hence transports the construction to
$1\amalg x^*A=G(x^*A)$.
\end{proof}

For sheaves on a topological space this internal coproduct has the explicit
form previously obtained by sheafification.

\begin{corollary}[Clopen support formula]
\label{cor:clopen-support-formula}
If $A=\cA$ is a sheaf of semirings on a topological space $X$, then
\[
\Spl(\cA)(U)
\cong
\coprod_{V\in\Clop(U)}\cA(V).
\]
A pair $(V,a)$ is the section equal to $a$ on the clopen support locus $V$
and absent on $U\setminus V$.  On connected $U$ this reduces to
$G(\cA(U))$.
\end{corollary}

\begin{proof}
A map $U\to1\amalg\cA$ has a locally constant support map
$U\to\{0,1\}$; its supported inverse image is clopen, and on that locus it
is exactly a section of $\cA$.  This correspondence is reversible and
compatible with the semiring operations.
\end{proof}

\subsection{The split arithmetic site}

The original arithmetic site is
\[
\mathfrak A=(\widehat{\N^\times},\Nbar),
\qquad
\Nbar=(\N\cup\{\infty\},\min,+),
\]
with the $n$-th Frobenius acting by multiplication by $n$ on finite
exponents and fixing $\infty$ \cite{CCArithmeticSite}.  Here $\infty$ is
the additive zero of the tropical semiring.

\begin{definition}[Split arithmetic site]
\label{def:split-arithmetic-site}
Define
\[
\mathfrak A^{\mathrm{sp}}
:=\bigl(\widehat{\N^\times},\Spl(\Nbar)\bigr).
\]
Write $\tau_{\mathrm{ar}}$ for the new zero and
$e_{\mathrm{ar}}:=\infty\in\Nbar$ for the supported tropical zero.
\end{definition}

\begin{theorem}[Frobenius-equivariant two-zero arithmetic site]
\label{thm:split-arithmetic-site}
Every universal Frobenius of $\Nbar$ extends uniquely to
$\Spl(\Nbar)$ by fixing $\tau_{\mathrm{ar}}$.  The support character
\[
\chi:\Spl(\Nbar)\to\B
\]
is $\N^\times$-equivariant.  At every point of the arithmetic topos the
stalk separates
\[
\tau_{\mathrm{ar}},
\qquad
e_{\mathrm{ar}}=\infty,
\qquad
n\in\N.
\]
Localizing at $e_{\mathrm{ar}}$ gives the Boolean support residue, whereas
localizing at $\tau_{\mathrm{ar}}$ gives the terminal semiring.
\end{theorem}

\begin{proof}
Functoriality in Theorem~\ref{thm:internal-split-functor} extends each
Frobenius and makes $\chi$ equivariant.  The stalk statement follows from
inverse-image compatibility.  The localization statement is the
idempotent division trichotomy applied stalkwise; the old tropical zero is
idempotent for multiplication because $\infty+\infty=\infty$.
\end{proof}

\begin{remark}[Synchronized and independent squares]
Theorem~\ref{thm:boolean-fibre-product} internalizes to semiring objects:
\[
\Spl(A\times B)\cong\Spl(A)\times_{\B}\Spl(B).
\]
This is the synchronized-support square.  The Cartesian product
$\Spl(A)\times\Spl(B)$ is larger and carries the full independent Boolean
square.  This distinction is separate from the Connes--Consani reduced and
unreduced tensor squares of the arithmetic site, which use
$\Nbar\otimes_{\B}\Nbar$ and its Newton-polygon reduction rather than a
Cartesian product.
\end{remark}

\subsection{The split scaling site}

Let
\[
\mathscr S=([0,\infty)\rtimes\N^\times,\cO)
\]
be the Connes--Consani scaling site, whose characteristic-one structure
sheaf consists of convex piecewise affine functions with integral slopes
\cite{CCScalingSite,CCAbsolute}.

\begin{definition}[Split scaling site]
\label{def:split-scaling-site}
Define
\[
\mathscr S^{\mathrm{sp}}
:=([0,\infty)\rtimes\N^\times,\Spl(\cO)).
\]
\end{definition}

\begin{theorem}[Three local scaling states]
\label{thm:three-scaling-states}
At every point $x$ of the scaling site, the split stalk separates
\[
\tau_x,
\qquad
e_x=0_{\cO_x},
\qquad
f\in\cO_x\setminus\{0\}.
\]
These are, respectively, absence of a local valuation profile, a present
profile equal to tropical zero, and a nonzero profile.  Boolean
localization at $e_x$ forgets all slope and amplitude data and retains only
support.
\end{theorem}

\begin{proof}
Apply Theorem~\ref{thm:internal-split-functor} at the point $x$ and then
the scalar localization theorem.
\end{proof}

\begin{remark}[Coefficient levels]
The semiringed split arithmetic site uses $\Spl(\Nbar)$, while the absolute
Frobenius-root stalks use the spherical algebras $\Fone[T^{H_+}]$.  The
Boolean semiring is the support coefficient, the tropical sheaf records
valuation profiles, and the spherical algebra records absolute
multiplicative generators.  The constructions in this paper preserve this
three-level coefficient hierarchy.
\end{remark}

\part{Theorem architecture and synthesis}

\section{Logical architecture and source interfaces}
\label{sec:construction-architecture}

The paper is organized as a sequence of exact functors and realization
maps.  The source of each layer is recorded here so that the derived
statements can be reused independently.

\begin{enumerate}[label=\textup{(\arabic*)},leftmargin=3.2em]
\item The split-zero algebra, Boolean fibre-product theorem, fixed-locus
formula, localizations, projective decompositions, projective semimodule
constraints, conservative ideal arithmetic, monad formulas, surcomplex
valuation results, and internal split-site construction are derived from
the definitions and the two split-zero source notes
\cite{Globalization,Secondary}.

\item The distributive-lattice globalization and the Boolean-character
classification are proved in Sections~\ref{sec:lattice-support}--
\ref{sec:nonassociative-support}.  Finite Birkhoff duality supplies the
multi-Boolean normal form.  The Cayley--Dickson threshold uses the
classical division-algebra and zero-divisor results in
\cite{BaezOctonions,MorenoCD,ImaedaCD,WilmotCD}; the tensor-category
associator is organized as in \cite{BHMTensor}.

\item The adelic-curve definitions, properness criterion, scalar height,
strongly adelic vector bundles, and global intersection integral come from
Chen--Moriwaki \cite{ChenMoriwakiBook,ChenMoriwakiIntersection}.  Section
\ref{sec:adelic-support-defect} derives the measure-algebra defect
valuation, its $L^\infty$ linearization and naturality, the
principal-logarithm pairing, weighted projective-height descent, weighted
Arakelov degree descent, global-field rigidity, the kernel splitting for
trivial-place extensions, canonical trivial norm extensions, and the
faithful local and measurable three-state metric lattices.  Arakelov geometry over trivially valued fields
is developed in \cite{ChenMoriwakiDirichlet}.

\item The absolute $\Fone$-site, its points and stalks, the map $\Theta$,
place sets, local torsors, periodic orbits, and scaling-site geometry come
from Connes--Consani
\cite{CCAbsolute,CCArithmeticSite,CCScalingSite}.  Their arithmetic Picard,
Jacobian, framed-divisor, rooted-divisor, and character-dual theorems come
from \cite{CCJacobian}.  The split product, reduction quotient,
denominator laws, and phase complementarity are derived here from those
structures.

\item The finite and global carry extensions, their Ext classes, carry
cocycles, cofinality, and the inverse-limit torsion jump are proved in
Sections~\ref{sec:lcm-perfection}--\ref{sec:perfected-anomaly}.  The
identification $\Hom(\Q,\Q/\Z)\cong\mathbb A_f$ and the rooted-torsion
formula are imported from \cite{CCJacobian}.  Cyclotomic symmetry and the
class-field-theory action are used in the form recorded in
\cite{BostConnes,CCBCCyclotomy,CMR}.

\item The ideal, prime, and spectral ordinal-sum theorems for $G_L(R)$
are proved directly in
Theorems~\ref{thm:lattice-ideal-classification}--
\ref{thm:spectral-ordinal-sum}.  The projective profile decomposition,
finite-field count, lattice determinant, and lattice-independent linear
group are Theorems~\ref{thm:lattice-projective-profile},
\ref{thm:lattice-finite-field-count},
\ref{thm:lattice-split-determinant}, and
\ref{thm:lattice-monomial-gl}.

\item Section~\ref{sec:split-matrix-invariants} proves the split determinant,
characteristic-coefficient, trace-word, cycle-generator, and UHF-compatibility
theorems.  The classical one-matrix and simultaneous-matrix invariant
theorems are used in the forms of Weyl--Procesi and Le Bruyn--Procesi
\cite{WeylClassical,ProcesiMatrices,LeBruynProcesi}.

\item The weighted-cycle normal form, conjugacy classification, trace and
Euler-factor formulas, and periodic-orbit factorization are proved in
Theorems~\ref{thm:weighted-cycle-classification} and
\ref{thm:monomial-periodic-euler-factor}.  Their scaling-site specialization
uses the Connes--Consani mapping-torus theorem \cite{CCKnotsPrimes}.

\item Section~\ref{sec:monomial-langlands} derives the equivalence between
normalizer local systems and line local systems on finite covers, identifies
the denominator phase tower with the torsion locus of
$\LocSys_{\Gm}^{\mathrm{Betti}}(E)$, and constructs the split Hitchin map.
The theta-group theorem is used from \cite{MumfordTheta}; the classical
Hitchin map and its quantization are used from
\cite{HitchinSystems,EtingofLiuHitchin}.  The coarse Betti geometric
Langlands functor is imported from \cite{GaitsgoryRaskinGLI}, and the
monomial spectral restriction is its pullback along the explicitly
constructed normalizer chart.

\item The finite quadratic-gerbe theory is the FHLT finite path-integral
theorem \cite{FHLT} applied to the explicitly constructed forms $q_n$ and
their even-lattice realizations in
Theorem~\ref{thm:cyclic-discriminant-lattice}.
The bundle-gerbe transgression and theta modules come from
\cite{GerbeTheta}; their identification with the finite Schr\"odinger
module is proved here.  The Lagrangian correspondence system, UHF limit,
ordered $K$-theory, profinite diagonal, crossed carry cocycle, and its
$H^1$ obstruction class are derived from the displayed constructions.

\item The helix development \cite{Lavery} supplies the formalized carrier,
multiplier, unitary block, determinant, and phasor identities.  Section
\ref{sec:helix-local-system} extracts their factor-of-automorphy form and
combines it with the fixed-locus theorem.

\item The reciprocal Beurling refinement originates in \cite{Schultze}.
The exact inclusion $\cR_Q\subseteq V_{\Lambda_Q}$ gives the strengthened
estimate $x_{\Lambda_Q}\leq d_Q^2$, and the Nyman--B\'aez-Duarte theorem
identifies density at the target section with the Riemann hypothesis.
\end{enumerate}

The principal synthesis chain is
\[
\begin{aligned}
\text{support lattice}
&\longrightarrow\text{split matrix invariants}
\longrightarrow\text{monomial spectral chart},\\
&\longrightarrow\text{adelic defect transform}
\longrightarrow\text{rigidified Frobenius tower},\\
&\longrightarrow\text{quadratic gerbe}
\longrightarrow\text{Heisenberg/theta boundary}
\longrightarrow\text{UHF limit}.
\end{aligned}
\]
Every arrow in this chain is defined in the corresponding section and is
compatible with the support projections developed in the opening algebraic
part of the paper.

\section{Synthesis: the support-amplitude theorem package}
\label{sec:final-synthesis}

\begin{theorem}[Support-amplitude synthesis]
\label{thm:final-synthesis-v11}
The constructions of this paper satisfy the following five linked
statements.

\begin{enumerate}[label=\textup{(\arabic*)},leftmargin=3.2em]
\item \textbf{Universal support geometry.}
For every nonzero commutative ring $R$ and nontrivial bounded distributive
lattice $L$, the semiring $G_L(R)$ has universal ring reflection $p_L$,
support projection $\chi_L$, and zero fibre $L$.  Its ideals and primes are
classified by Theorems~\ref{thm:lattice-ideal-classification} and
\ref{thm:lattice-prime-classification}, and
\[
\Spec G_L(R)\cong\SpecLat(L)\blacktriangleright\Spec R.
\]
Support localization is principal-interval restriction, while ordinary
nonzero localization preserves the amplitude field or ring.  The support
coefficient object is pre-Boolean: Boolean characters are prime branches,
branch base change selects one support generization, and the coefficient
hierarchy is
\[
\Fone\longrightarrow\Fnot(L)\xrightarrow{\epsilon}\B.
\]

\item \textbf{Projective and invariant geometry.}
Over a field $K$, projective space decomposes by lattice-valued support
profiles, with finite-field count
\[
\#\Pum^n(G_L(\F_q))
=\frac{(q+|L|-1)^{n+1}-|L|^{n+1}}{q-1}.
\]
The determinant and characteristic coefficients have simultaneous
amplitude and lattice-support values.  Every invertible matrix is monomial,
independently of $L$, and monomial conjugacy is classified by weighted
permutation cycles.  Those cycles compute characteristic polynomials,
traces, and Euler factors.

\item \textbf{Adelic and absolute arithmetic realization.}
For a proper adelic curve $S$, the measure algebra $\mathscr B_S$ is a
complete Boolean support lattice and the product formula is the vanishing
of the defect transform $\mathscr D_S$.  For every measurable support $A$,
\[
G_{\mathscr B_S}(K)[z_A^{-1}]
\cong\mathscr B_{S|_A},
\]
and this localization carries a proper restricted adelic curve exactly
when $A$ is balanced.  The Connes--Consani arithmetic Jacobian is
$\Picf\times\B$, the split spectrum resolves the zero Abel--Jacobi fibre,
and rigidified denominator points form the Frobenius divisibility tower.

\item \textbf{Frobenius perfection and quantization.}
The LCM colimit is the globally Frobenius-perfect point
$H_\Omega=\Q$ with stalk $\Fone[T^{\Q_+}]$, carry extension
$0\to\Z\to\Q\to\Q/\Z\to0$, and finite-adelic character dual.  Finite
quadratic refinements define gerbe theories whose torus transgression is the
finite Heisenberg algebra and whose theta module is the Schr\"odinger
representation.  The strict matrix system has UHF limit
$\mathcal Q_\Omega$ with ordered $K_0$ equal to $\Q$ and diagonal
$C(\widehat\Z)\cong C^*(\Q/\Z)$.

\item \textbf{Spectral descent and analytic endpoint.}
The identity
$\GL_n(G_L(K))=N_{\GL_n}(T)(K)$ identifies split local systems with line
local systems on finite covers.  Their loop determinants factor over
covering cycles, and the scaling-site specialization factors over
Frobenius cycles.  Split characteristic coefficients define a
support-refined Hitchin map, and pullback along the normalizer chart gives
the monomial geometric Langlands restriction.  At the perfect arithmetic
stalk, reciprocal roots have Beurling realization and satisfy
$x_{\Lambda_Q}\le d_Q^2$; density at the target section is equivalent to
the Riemann hypothesis.
\end{enumerate}
\end{theorem}

\begin{proof}
Statement (1) is the content of
Theorems~\ref{thm:lattice-normal-form},
\ref{thm:lattice-boolean-characters},
\ref{thm:lattice-localization}, and
\ref{thm:spectral-ordinal-sum}.  Statement (2) follows from
Theorems~\ref{thm:lattice-projective-profile},
\ref{thm:lattice-finite-field-count},
\ref{thm:lattice-split-determinant},
\ref{thm:lattice-monomial-gl}, and
\ref{thm:weighted-cycle-classification}.  Statement (3) is
Theorems~\ref{thm:balanced-support-localization},
\ref{thm:jacobian-boolean-product},
\ref{thm:abel-jacobi-resolution}, and
\ref{thm:denominator-frobenius}.  Statement (4) is
Theorems~\ref{thm:global-perfection},
\ref{thm:global-carry-class},
\ref{thm:finite-gerbe-anomaly},
\ref{thm:finite-stone-von-neumann}, and
\ref{thm:perfected-uhf}.  Statement (5) is
Theorems~\ref{thm:monomial-spectral-cover},
\ref{thm:monomial-periodic-euler-factor},
\ref{thm:split-hitchin-map},
\ref{thm:monomial-langlands-restriction},
\ref{thm:direct-beurling-domination}, and
\ref{thm:arithmetic-site-nyman}.
\end{proof}

\subsection*{Mathematical consequences}

The theorem package supplies three reusable mechanisms.

First, support should be computed before it is Booleanized.  The
pre-Boolean object $\Fnot(L)$ retains all branches at once; applying a
Boolean character chooses one branch and then recovers the usual
split-zero arithmetic geometry.  Second, a support refinement can be
computed before amplitude is forgotten.
This produces exact incidence invariants for determinants, cycle covers,
adelic place sets, and local sectors.  Third, localization has a geometric
interpretation: in finite lattice geometry it restricts to a principal
support interval, while on an adelic curve it restricts the measured place
family and tests properness through the product formula.  Fourth, monomial
spectral descent converts finite-cover data into explicit cycle factors,
which link classical invariant coordinates, Frobenius monodromy, theta
quantization, and local Euler determinants.

The resulting framework is therefore a support-amplitude geometry with
three mutually compatible realizations:
\[
\begin{tikzcd}[column sep=huge,row sep=large]
\text{lattice and semiring geometry}
  \arrow[r,"\text{adelic localization}"]
  \arrow[dr,"\text{matrix invariants}"']
& \text{arithmetic curves and Frobenius towers}
  \arrow[d,"\text{finite quantization}"]\\
& \text{monomial spectral charts and Euler factors}.
\end{tikzcd}
\]
The arrows are the explicit functors and formulas established in the body
of the paper.


\begin{thebibliography}{99}

\bibitem{Globalization}
ChatGPT 5.4 and Gemini 2.5 DeepThink,
\emph{A Note on the Globalization Semiring $G(\Z)$},
Zenodo, 2026.
DOI: \href{https://doi.org/10.5281/zenodo.18976982}{10.5281/zenodo.18976982}.

\bibitem{Secondary}
The Clankers,
\emph{A Secondary Note on Split-Zero Globalization: Semimodule
Classification, Multiplicative Reconstruction, and All-Orders Zeta
Deformation},
Zenodo, 2026.
DOI: \href{https://doi.org/10.5281/zenodo.19120905}{10.5281/zenodo.19120905}.

\bibitem{Homogenization}
The Clankers,
\emph{Support-Idempotent Homogenization for a Split-Zero Collatz F-Series
Frame},
Zenodo, 2026.
DOI: \href{https://doi.org/10.5281/zenodo.19900461}{10.5281/zenodo.19900461}.

\bibitem{CCAbsolute}
A.~Connes and C.~Consani,
\emph{On the Absolute Geometry of $\operatorname{Spec}\Z$},
arXiv:2606.06604 [math.AG], 2026.
DOI: \href{https://doi.org/10.48550/arXiv.2606.06604}{10.48550/arXiv.2606.06604}.

\bibitem{CCArithmeticSite}
A.~Connes and C.~Consani,
\emph{Geometry of the Arithmetic Site},
Adv. Math. \textbf{291} (2016), 274--329.

\bibitem{CCScalingSite}
A.~Connes and C.~Consani,
\emph{Geometry of the Scaling Site},
Selecta Math. (N.S.) \textbf{23} (2017), no.~3, 1803--1850.


\bibitem{CCJacobian}
A.~Connes and C.~Consani,
\emph{On the Jacobian of $\operatorname{Spec}\Z$},
arXiv:2602.15941 [math.NT], 2026.
DOI: \href{https://doi.org/10.48550/arXiv.2602.15941}{10.48550/arXiv.2602.15941}.

\bibitem{CCAbsoluteAlgebra}
A.~Connes and C.~Consani,
\emph{Absolute algebra and Segal's $\Gamma$-rings: au-dessous de
$\operatorname{Spec}(\Z)$},
J. Number Theory \textbf{162} (2016), 518--551.

\bibitem{CCAffineAbsolute}
A.~Connes and C.~Consani,
\emph{On absolute algebraic geometry: the affine case},
Adv. Math. \textbf{390} (2021), Article 107909.

\bibitem{CCGammaUniversal}
A.~Connes and C.~Consani,
\emph{Segal's gamma rings and universal arithmetic},
Q. J. Math. \textbf{72} (2021), 1--29.


\bibitem{BostConnes}
J.-B.~Bost and A.~Connes,
\emph{Hecke algebras, type III factors and phase transitions with
spontaneous symmetry breaking in number theory},
Selecta Math. (N.S.) \textbf{1} (1995), no.~3, 411--457.

\bibitem{CCBCCyclotomy}
A.~Connes and C.~Consani,
\emph{BC-system, absolute cyclotomy and the quantized calculus},
EMS Surv. Math. Sci. \textbf{9} (2022), 447--475.
DOI: \href{https://doi.org/10.4171/EMSS/64}{10.4171/EMSS/64}.

\bibitem{CMR}
A.~Connes, M.~Marcolli, and N.~Ramachandran,
\emph{KMS states and complex multiplication},
Selecta Math. (N.S.) \textbf{11} (2005), no.~3--4, 325--347.

\bibitem{ScholzePerfectoid}
P.~Scholze,
\emph{Perfectoid spaces},
Publ. Math. Inst. Hautes Etudes Sci. \textbf{116} (2012), 245--313.

\bibitem{Weibel}
C.~A.~Weibel,
\emph{An Introduction to Homological Algebra},
Cambridge Studies in Advanced Mathematics \textbf{38},
Cambridge University Press, Cambridge, 1994.

\bibitem{FarguesFontaine}
L.~Fargues and J.-M.~Fontaine,
\emph{Courbes et fibr\'es vectoriels en th\'eorie de Hodge $p$-adique},
Ast\'erisque \textbf{406} (2018), xiii+382 pp.

\bibitem{Lavery}
S.~Lavery,
\emph{The Double-Ended Helix: An Intrinsically On-Axis Geometric State
Space for Dirichlet $L$-Function Phasors},
manuscript with Lean~4 formalization, 2026,
\url{https://github.com/samlavery/helix_frobenius}.

\bibitem{Schultze}
K.~Schultze,
\emph{A Dyadic-to-Integer Beurling Reduction for a Reciprocal
Nyman--Beurling Refinement},
preliminary research note, v0.1.2, 2026,
\url{https://github.com/KarlSchultze/dyadic-beurling-reduction}.

\bibitem{Nyman}
B.~Nyman,
\emph{On the One-Dimensional Translation Group and Semi-Group in Certain
Function Spaces},
Ph.D. thesis, Uppsala University, 1950.

\bibitem{Beurling}
A.~Beurling,
\emph{A closure problem related to the Riemann zeta-function},
Proc. Nat. Acad. Sci. U.S.A. \textbf{41} (1955), 312--314.

\bibitem{BaezDuarte}
L.~B\'aez-Duarte,
\emph{A strengthening of the Nyman--Beurling criterion for the Riemann
hypothesis},
Rend. Mat. Acc. Lincei \textbf{14} (2003), 5--11.

\bibitem{Steinitz}
E.~Steinitz,
\emph{Algebraische Theorie der K\"orper},
J. Reine Angew. Math. \textbf{137} (1910), 167--309.

\bibitem{BerarducciMantova}
A.~Berarducci and V.~Mantova,
\emph{Surreal numbers, derivations and transseries},
J. Eur. Math. Soc. \textbf{20} (2018), no.~2, 339--390.
DOI: \href{https://doi.org/10.4171/JEMS/769}{10.4171/JEMS/769}.

\bibitem{Conway}
J.~H.~Conway,
\emph{On Numbers and Games},
2nd ed., A K Peters, Natick, MA, 2001.

\bibitem{Deitmar}
A.~Deitmar,
\emph{Schemes over $\Fone$},
in \emph{Number Fields and Function Fields---Two Parallel Worlds},
Progr. Math. \textbf{239}, Birkh\"auser, Boston, 2005, 87--100.

\bibitem{Golan}
J.~S.~Golan,
\emph{Semirings and Their Applications},
Kluwer Academic Publishers, Dordrecht, 1999.

\bibitem{Goldblatt}
R.~Goldblatt,
\emph{Lectures on the Hyperreals: An Introduction to Nonstandard Analysis},
Graduate Texts in Mathematics \textbf{188}, Springer, New York, 1998.

\bibitem{GrothendieckEGAII}
A.~Grothendieck,
\emph{\'El\'ements de g\'eom\'etrie alg\'ebrique II: \'Etude globale
\'el\'ementaire de quelques classes de morphismes},
Publ. Math. IH\'ES \textbf{8} (1961), 5--222.

\bibitem{HoTT}
The Univalent Foundations Program,
\emph{Homotopy Type Theory: Univalent Foundations of Mathematics},
Institute for Advanced Study, 2013,
\url{https://homotopytypetheory.org/book/}.

\bibitem{Lorscheid}
O.~Lorscheid,
\emph{The geometry of blueprints. Part I: Algebraic background and scheme
theory},
Adv. Math. \textbf{229} (2012), no.~3, 1804--1846.

\bibitem{Robinson}
A.~Robinson,
\emph{Non-standard Analysis},
North-Holland, Amsterdam, 1966.

\bibitem{Simons}
J.~Simons,
\emph{Meet the Surreal Numbers},
expository notes, 4 April 2017.


\bibitem{ChenMoriwakiBook}
H.~Chen and A.~Moriwaki,
\emph{Arakelov Geometry over Adelic Curves},
Lecture Notes in Mathematics \textbf{2258}, Springer, Singapore, 2020.

\bibitem{ChenMoriwakiIntersection}
H.~Chen and A.~Moriwaki,
\emph{Arithmetic intersection theory over adelic curves},
pre-publication monograph, Version 1, March 2021.

\bibitem{ChenMoriwakiDirichlet}
H.~Chen and A.~Moriwaki,
\emph{Sufficient conditions for the Dirichlet property},
preprint, Version 1.0, April 2017.

\bibitem{FHLT}
D.~S.~Freed, M.~J.~Hopkins, J.~Lurie, and C.~Teleman,
\emph{Topological quantum field theories from compact Lie groups},
arXiv:0905.0731v2, 2009.

\bibitem{GerbeTheta}
The Clankers,
\emph{Coupled Gerbes, Transgression, Explicit Theta Sectors, Hesse Cubics,
and Landau/BCM Structures on $T^3$},
pre-publication manuscript, March 2026.

\bibitem{BaezOctonions}
J.~C.~Baez,
\emph{The octonions},
Bull. Amer. Math. Soc. (N.S.) \textbf{39} (2002), 145--205.

\bibitem{MorenoCD}
G.~Moreno,
\emph{The zero divisors of the Cayley--Dickson algebras over the real
numbers},
Bol. Soc. Mat. Mexicana (3) \textbf{4} (1998), 13--28.

\bibitem{ImaedaCD}
K.~Imaeda and M.~Imaeda,
\emph{Sedenions: algebra and analysis},
Appl. Math. Comput. \textbf{115} (2000), 77--88.

\bibitem{WilmotCD}
G.~P.~Wilmot,
\emph{Structure of the Cayley--Dickson algebras},
arXiv:2505.11747v1, 2025.

\bibitem{BHMTensor}
P.~Bouwknegt, K.~C.~Hannabuss, and V.~Mathai,
\emph{$C^*$-algebras in tensor categories},
arXiv:math/0702802v2, 2009.

\bibitem{GlimmUHF}
J.~Glimm,
\emph{On a certain class of operator algebras},
Trans. Amer. Math. Soc. \textbf{95} (1960), 318--340.


\bibitem{BirkhoffLattice}
G.~Birkhoff,
\emph{Lattice Theory},
3rd ed., American Mathematical Society Colloquium Publications
\textbf{25}, American Mathematical Society, Providence, RI, 1967.

\bibitem{NeukirchANT}
J.~Neukirch,
\emph{Algebraic Number Theory},
Grundlehren der mathematischen Wissenschaften \textbf{322},
Springer, Berlin, 1999.

\bibitem{Stichtenoth}
H.~Stichtenoth,
\emph{Algebraic Function Fields and Codes},
2nd ed., Graduate Texts in Mathematics \textbf{254}, Springer, Berlin,
2009.


\bibitem{WeylClassical}
H.~Weyl,
\emph{The Classical Groups: Their Invariants and Representations},
2nd ed., Princeton University Press, Princeton, 1946.

\bibitem{ProcesiMatrices}
C.~Procesi,
\emph{The invariant theory of $n\times n$ matrices},
Adv. Math. \textbf{19} (1976), 306--381.

\bibitem{LeBruynProcesi}
L.~Le Bruyn and C.~Procesi,
\emph{Semisimple representations of quivers},
Trans. Amer. Math. Soc. \textbf{317} (1990), 585--598.

\bibitem{MumfordTheta}
D.~Mumford,
\emph{Tata Lectures on Theta I},
Progress in Mathematics \textbf{28}, Birkh\"auser, Boston, 1983.

\bibitem{HitchinSystems}
N.~J.~Hitchin,
\emph{Stable bundles and integrable systems},
Duke Math. J. \textbf{54} (1987), 91--114.

\bibitem{EtingofLiuHitchin}
P.~Etingof and H.~Liu,
\emph{Hitchin systems and their quantization},
arXiv:2409.09505, 2024.

\bibitem{GaitsgoryRaskinGLI}
D.~Gaitsgory and S.~Raskin,
\emph{Proof of the geometric Langlands conjecture I: construction of the
functor},
arXiv:2405.03599, 2024.


\bibitem{CCKnotsPrimes}
A.~Connes and C.~Consani,
\emph{Knots, Primes and the Adele Class Space},
preprint, 2026.

\end{thebibliography}

\end{document}
